% =============================================================================
%  最优传输的定性理论 — Villani《Optimal Transport: Old and New》第 I 部分 (第 4–8 章)
%  模式：被动自学（导师讲义 passive-study）  ·  主题：teal  ·  语言：中文
%  来源：C. Villani, Optimal Transport: Old and New, Springer 2009.
%        Part I "Qualitative description of optimal transport", Chapters 4–8.
%        (PDF pp. 55–210 / 印刷页 pp. 39–198)
%
%  所有示意图用 TikZ（teal 配色）重绘，标注 "据 Villani 图 x.y 重绘"。
%  不使用 beamer overlay 命令 (\pause/\onslide/\only/\uncover)。
%
%  CJK 字体：font-config 自动检测 .fonts/SourceHanSerifSC-Medium.otf（思源宋体）
%  数学字体：KpMath（模板默认）
%
%  本文件为「主控文件」：前置内容在此，各节 \input 自 slides_assets/ot/secN-*.tex
%  组装顺序由 lead-codex 指挥官统一维护；各节文件遵循 STYLE-CONTRACT.md。
% =============================================================================
\documentclass[aspectratio=169,10pt]{beamer}
\usepackage[teal]{template-lib/template-lib}
\uselayout{text}
\uselayout{eq}
\uselayout{twocol}
\uselayout{1img}
\uselayout{2img}
\uselayout{table}

\usetikzlibrary{arrows.meta,calc,positioning,patterns,decorations.pathmorphing,%
                fit,backgrounds,shapes.geometric,shapes.misc,intersections}

\graphicspath{{slides_assets/ot/}{slides_assets/ot/source_figures/}}

% ── 记号缩写（全讲义统一，secN 文件不得重复定义）──────────────────────────────
\newcommand{\R}{\mathbb{R}}
\newcommand{\N}{\mathbb{N}}
\newcommand{\E}{\mathbb{E}}
\newcommand{\Prob}{\mathbb{P}}
\newcommand{\X}{\mathcal{X}}
\newcommand{\Y}{\mathcal{Y}}
\newcommand{\Z}{\mathcal{Z}}
\newcommand{\Pp}{\mathcal{P}}                 % 概率测度集
\newcommand{\Ppp}[1]{\mathcal{P}_{#1}}        % P_p（p 阶矩有限）
\newcommand{\Pac}{\mathcal{P}^{\mathrm{ac}}}  % 绝对连续概率测度
\newcommand{\coup}[2]{\Pi(#1,#2)}             % 耦合（转移计划）集合
\newcommand{\law}{\operatorname{law}}
\newcommand{\Id}{\mathrm{Id}}
\newcommand{\supp}{\operatorname{supp}}
\newcommand{\dom}{\operatorname{dom}}
\newcommand{\pf}[2]{{#1}_{\#}#2}              % 推前 T_# mu
\newcommand{\Wp}[1]{W_{#1}}                   % Wasserstein 距离 W_p
\newcommand{\Wtwo}{W_{2}}
\newcommand{\ctr}[1]{{#1}^{c}}               % c-变换 psi^c
\newcommand{\cctr}[1]{{#1}^{cc}}            % 二次 c-变换
\newcommand{\inner}[2]{\langle #1,\,#2\rangle}
\newcommand{\abs}[1]{\lvert #1\rvert}
\newcommand{\norm}[1]{\lVert #1\rVert}
\newcommand{\dd}{\,\mathrm{d}}
\newcommand{\cost}{c}
\newcommand{\half}{\tfrac{1}{2}}
\newcommand{\eps}{\varepsilon}
\newcommand{\To}{\longrightarrow}
\newcommand{\wkto}{\rightharpoonup}           % 弱收敛
\newcommand{\subjto}{\;\text{s.t.}\;}

% ── 被动自学讲义的结构标记 ────────────────────────────────────────────────────
% 关键术语首次出现
% \TLterm{} 已由模板提供
% 难度星标（amssymb \bigstar），accent 色
\newcommand{\diffI}{\textcolor{TLaccent}{\ensuremath{\bigstar}}}
\newcommand{\diffII}{\textcolor{TLaccent}{\ensuremath{\bigstar\bigstar}}}
\newcommand{\diffIII}{\textcolor{TLaccent}{\ensuremath{\bigstar\bigstar\bigstar}}}
% 习题提示行
\newcommand{\hint}[1]{\emph{提示.}~#1}
% “定义 / 定理 / 命题 / 引理” 行内标签（plain-text-first：默认不进彩框）
\newcommand{\deflab}[1]{\textbf{\textcolor{TLprimaryDark}{#1}}}
% 文献注记标签（用于 References / Bibliographical notes 帧，含英文锚点供校验器识别）
\newcommand{\bibnote}[1]{\emph{文献注记 (bibliographical notes).}~#1}

% ── 本地化“要点”标签（要点。Key Takeaway）─────────────────────────────────────
\renewcommand{\TLtakeaway}[1]{%
  \vspace{0.4em}%
  \begin{tikzpicture}
    \node[fill=TLaccentLight, rounded corners=3pt, inner sep=8pt,
          text width=\linewidth-16pt, align=justify] {%
      \small\textcolor{TLaccent}{\textbf{要点。}} #1%
    };
  \end{tikzpicture}%
  \vspace{0.2em}%
}

\TLsetfoot{最优传输的定性理论 \textperiodcentered{} Villani 第 I 部分（第 4--8 章）}

% =============================================================================
\begin{document}

% ── 封面 ─────────────────────────────────────────────────────────────────────
\TLtitlecenter
  {最优传输的定性理论}
  {耦合 · Kantorovich 对偶 · Wasserstein 距离 · 位移插值 · 缩短原理}
  {自学导师讲义 \textperiodcentered{} 据 Villani《Optimal Transport: Old and New》第 4--8 章}
  {Qualitative Description of Optimal Transport}

% ── 如何使用本讲义 ───────────────────────────────────────────────────────────
\begin{frame}{如何使用本讲义 · How to Use}
  这是一份\TLterm{被动自学（passive-study）}的\emph{导师讲义}：它会把背景、动机、
  定义、证明步骤、可算的例子和习题\emph{讲全}，使你无需翻开原书即可读懂。
  \begin{itemize}
    \item \textbf{先动机，后形式。} 每个概念先回答“\emph{为什么需要它}”，再给形式定义，
          并在两页内给出一个\textbf{可算的例子 (worked example)}。
    \item \textbf{自足。} 第 4--8 章所需的测度论、耦合与两个传输问题的背景，
          统一在第 0 节补齐；专业术语首次出现时给出中文直觉。
    \item \textbf{常见陷阱}用警示框标出；\textbf{习题 (exercises)} 附\emph{提示 (hints)}，
          完整\emph{解答 (solutions)} 见附录 (appendix)。
    \item \textbf{记号}集中在末尾的\emph{术语表与记号表 (glossary \& notation)}。
  \end{itemize}
  \TLtakeaway{读法：遇到定义先想例子，遇到定理先想“它在说什么几何事实”，再看证明。}
\end{frame}

% ── 全局路线图 ───────────────────────────────────────────────────────────────
\begin{frame}{全局路线图：定性理论的五块拼图}
  \begin{center}
  \begin{tikzpicture}[>={Stealth[length=5pt]},node distance=4mm,
      box/.style={draw=TLprimary,fill=TLprimaryTint,rounded corners=3pt,
        align=center,inner sep=5pt,text width=2.35cm,font=\scriptsize},
      lab/.style={font=\scriptsize\itshape,text=TLinkSoft}]
    \node[box] (c4) {\textbf{§1 · Ch.4}\\最优耦合\\\emph{存在性}};
    \node[box,right=of c4] (c5) {\textbf{§2 · Ch.5}\\Kantorovich\\\emph{对偶}};
    \node[box,right=of c5] (c6) {\textbf{§3 · Ch.6}\\Wasserstein\\\emph{距离}};
    \node[box,right=of c6] (c7) {\textbf{§4 · Ch.7}\\位移插值\\\emph{测地线}};
    \node[box,right=of c7] (c8) {\textbf{§5 · Ch.8}\\缩短原理\\\emph{正则性}};
    \draw[->,TLaccent,thick] (c4)--(c5);
    \draw[->,TLaccent,thick] (c5)--(c6);
    \draw[->,TLaccent,thick] (c6)--(c7);
    \draw[->,TLaccent,thick] (c7)--(c8);
    \node[lab,below=5mm of c4.south] {解存在吗？};
    \node[lab,below=5mm of c5.south] {如何刻画最优？};
    \node[lab,below=5mm of c6.south] {测度间的几何};
    \node[lab,below=5mm of c7.south] {最优如何随时间演化};
    \node[lab,below=5mm of c8.south] {轨迹有多光滑};
  \end{tikzpicture}
  \end{center}
  \vspace{2pt}
  \begin{itemize}\small
    \item 主线：\emph{存在} $\to$ \emph{对偶刻画} $\to$ \emph{几何（距离）} $\to$
          \emph{动力学（测地线）} $\to$ \emph{正则性}。
    \item 第 0 节先补齐前置背景（测度、推前、耦合、Monge 与 Kantorovich 两个问题）。
  \end{itemize}
  \TLtakeaway{本讲义只覆盖 Villani 第 I 部分“定性描述”的核心五章；第 9--13 章（Monge 问题可解性、Jacobian、光滑性）不在本范围。}
\end{frame}

% =============================================================================
%  各节（指挥官维护组装顺序；secN 文件遵循 STYLE-CONTRACT.md）
% =============================================================================
% =============================================================================
%  §0  预备：测度、耦合与两个传输问题
%  指挥官撰写 —— 本节为「风格范本」。Codex 各节须在记号、版式、密度、图风、
%  worked-example 与 takeaway 习惯上与本节保持一致。
%  来源：Villani 引言 + 第 1–3 章（不在 4–8 范围内，但为其前置背景）。
% =============================================================================
\TLsection{0 \ 预备：测度、耦合与两个传输问题}

% ── 本节定位 ─────────────────────────────────────────────────────────────────
\begin{frame}{本节要补齐什么？}
  第 4--8 章默认你熟悉\emph{概率测度}、\emph{弱收敛}、\emph{推前}以及\emph{两个传输问题}。
  本节把这些背景一次补齐，使全讲义\textbf{自足}。
  \begin{itemize}
    \item \textbf{对象。} 在 Polish 空间上的概率测度 $\Pp(\X)$，以及如何度量它们的“靠近”。
    \item \textbf{动作。} 用映射搬运质量（\emph{推前} $\pf{T}{\mu}$），以及质量守恒。
    \item \textbf{两个问题。} Monge 的\emph{确定性}搬运 vs.\ Kantorovich 的\emph{可分裂}搬运。
    \item \textbf{一个伏笔。} 对偶（“面包店”经济直觉）——第 2 节的主角。
  \end{itemize}
  \TLtakeaway{把“最优传输”读成一句话：\emph{在所有把 $\mu$ 搬成 $\nu$ 的方案里，找总成本最小的那个}。本节把这句话里的每个词讲清楚。}
\end{frame}

% ── 概率测度空间 ─────────────────────────────────────────────────────────────
\begin{frame}{对象：Polish 空间上的概率测度 $\Pp(\X)$}
  \textbf{为什么。} 传输的两端是\emph{质量分布}，不是单个点；最自然的语言是\emph{概率测度}。

  \medskip
  \deflab{约定。} $\X,\Y$ 为\TLterm{Polish 空间}（完备可分度量空间，如 $\R^n$）。
  $\Pp(\X)$ 记其上全体\emph{概率测度}。一个 $\mu\in\Pp(\X)$ 给每个（可测）集合 $A$ 一个质量
  $\mu[A]\in[0,1]$，且 $\mu[\X]=1$。积分 $\int_{\X} \varphi \dd\mu$ 表示“按密度 $\mu$
  对检验函数 $\varphi\in C_b(\X)$ 取平均”：
  \[
    \mu \;\longmapsto\; \int_{\X} \varphi \dd\mu, \qquad \varphi\in C_b(\X).
  \]
  \begin{itemize}\small
    \item \emph{例 1（点质量）。} $\delta_{x_0}$：全部质量集中在 $x_0$，$\int\varphi\dd\delta_{x_0}=\varphi(x_0)$。
    \item \emph{例 2（经验测度）。} $\frac1n\sum_{i=1}^n\delta_{x_i}$：$n$ 个等重样本点。
    \item \emph{例 3（密度）。} $\mu=f\,\mathrm{Leb}$，$f\ge0,\ \int f=1$（如高斯）。
  \end{itemize}
  \TLtakeaway{$\Pp(\X)$ 同时容纳“离散沙堆”与“连续密度”——这正是传输需要的统一框架。}
\end{frame}

% ── 弱收敛 + 紧性 ────────────────────────────────────────────────────────────
\begin{frame}{度量“靠近”：弱收敛与 Prokhorov 紧性}
  \textbf{为什么。} 要证“最优方案\emph{存在}”，标准武器是\emph{紧性 + 下半连续}（第 1 节会用到）。
  先把“测度列收敛”说清楚。

  \medskip
  \deflab{弱（窄）收敛。} $\mu_k\wkto\mu$ 指对每个有界连续 $\varphi$，
  \[
    \int_{\X} \varphi\dd\mu_k \;\To\; \int_{\X} \varphi\dd\mu .
  \]
  直觉：沙堆\emph{整体形状}收敛，允许质量\emph{连续地移动}，但不允许质量\emph{逃逸到无穷}。

  \medskip
  \deflab{Prokhorov 定理（紧性判据）。} 一族测度\emph{胶着 (tight)} ——
  即对任意 $\eps>0$ 存在紧集 $K_\eps$ 使 $\mu[\X\setminus K_\eps]\le\eps$ 对全族成立——
  当且仅当它\emph{弱列紧}。
  \TLtakeaway{“胶着”= 质量不会偷偷跑到无穷远。它是后面所有\emph{存在性}证明的引擎。}
\end{frame}

% ── 推前测度 ─────────────────────────────────────────────────────────────────
\begin{frame}{动作：用映射搬运质量 —— 推前 $\pf{T}{\mu}$}
  \textbf{为什么。} “把 $\mu$ 搬成 $\nu$”最朴素的方式是给每个 $x$ 指定一个去处 $T(x)$。

  \medskip
  \deflab{推前（pushforward）。} 给可测 $T:\X\to\Y$，定义 $\nu=\pf{T}{\mu}$ 为
  \[
    \nu[B] \;=\; \mu\!\left[T^{-1}(B)\right]\quad(\forall B),
    \qquad\Longleftrightarrow\qquad
    \int_{\Y} \varphi\dd(\pf{T}{\mu}) = \int_{\X} \varphi\!\circ\! T \dd\mu .
  \]
  这恰是\emph{质量守恒}：$T$ 把落在 $T^{-1}(B)$ 的质量全部送进 $B$。

  \medskip
  \deflab{可算的例子（1D）。} 设 $\mu=\mathrm{Unif}[0,1]$，$T(x)=x^2$。则对 $y\in(0,1)$，
  \[
    \nu[(0,y)] = \mu[\,x^2<y\,] = \mu[\,x<\sqrt y\,] = \sqrt y
    \;\Rightarrow\; \nu \text{ 的密度 } g(y)=\frac{1}{2\sqrt y}.
  \]
  \TLtakeaway{推前 = “换元公式”的测度版。记号 $\pf{T}{\mu}=\nu$ 读作“$T$ 把 $\mu$ 搬成了 $\nu$”。}
\end{frame}

% ── Monge 问题 ───────────────────────────────────────────────────────────────
\begin{frame}{问题一：Monge —— 确定性搬运（可能无解）}
  \deflab{Monge 问题 (1781)。} 给定 $\mu\in\Pp(\X),\ \nu\in\Pp(\Y)$ 与成本 $c(x,y)\ge0$，求\emph{传输映射}
  \[
    \inf_{T}\ \int_{\X} c\big(x,\,T(x)\big)\dd\mu(x)
    \subjto \pf{T}{\mu}=\nu .
  \]
  即：每粒质量\emph{整粒}送到唯一去处 $T(x)$，且送达后的分布恰为 $\nu$。

  \medskip
  \deflab{它为何会“无解”。} Monge 不允许\emph{分裂}质量。若 $\mu=\delta_0$（一个点）而
  $\nu=\tfrac12\delta_{-1}+\tfrac12\delta_{1}$（两个点），则\emph{没有任何函数} $T$ 能把一个点
  送成两个点——可行集\emph{为空}。
  \begin{itemize}\small
    \item 即便可行，目标关于 $T$ \emph{高度非凸}，约束 $\pf{T}{\mu}=\nu$ 也非线性。
  \end{itemize}
  \TLwarnblock{陷阱}{“最优传输映射存在吗”是一个\emph{很难}的问题（第 9--12 章主题）。
  本讲义先绕开它——靠下面的松弛保证\emph{总有最优方案}。}
\end{frame}

% ── Kantorovich 松弛 ─────────────────────────────────────────────────────────
\begin{frame}{问题二：Kantorovich —— 允许分裂的搬运（总有解）}
  \deflab{耦合 / 转移计划。} $\mu,\nu$ 的\TLterm{耦合}是 $\X\times\Y$ 上的概率测度
  $\pi$，其两个\emph{边际}分别为 $\mu,\nu$：
  \[
    \pf{(\mathrm{proj}_{\X})}{\pi}=\mu,\quad \pf{(\mathrm{proj}_{\Y})}{\pi}=\nu .
    \qquad \text{记全体耦合为 } \coup{\mu}{\nu}.
  \]
  $\pi(dx\,dy)$ = “从 $x$ 附近搬一点质量到 $y$ 附近”的量；允许一个 $x$ 的质量\emph{分裂}去多个 $y$。

  \medskip
  \deflab{Kantorovich 问题 (1942)。}
  \[
    \boxed{\ C(\mu,\nu)\;=\;\inf_{\pi\in\coup{\mu}{\nu}}\ \int_{\X\times\Y} c(x,y)\dd\pi(x,y)\ }
  \]
  \begin{itemize}\small
    \item 目标对 $\pi$ \emph{线性}，约束\emph{线性凸}，$\coup{\mu}{\nu}$ 非空（含乘积测度 $\mu\otimes\nu$）且\emph{弱紧}。
    \item $\Rightarrow$ 最优 $\pi$ \emph{总存在}（第 1 节严格证明）。Monge 是 Kantorovich 中“不分裂”的特例。
  \end{itemize}
  \TLtakeaway{把“映射”松弛成“耦合”，问题从\emph{非凸、可能无解}变为\emph{线性、总有解}——这是全部理论的起点。}
\end{frame}

% ── 图：Monge vs Kantorovich ─────────────────────────────────────────────────
\begin{frame}{一图看懂：不可分裂 vs 可分裂}
  \begin{columns}[T]
    \begin{column}{0.5\textwidth}
      \centering {\small\textbf{Monge：整粒搬运}}\par\vspace{4pt}
      \begin{tikzpicture}[scale=0.82,>={Stealth[length=4pt]}]
        \foreach \y in {2,1.2,0.4}{\fill[TLprimary] (0,\y) circle (2pt);}
        \foreach \y in {2,1.2,0.4}{\fill[TLaccent] (3,\y) circle (2pt);}
        \draw[->,TLink] (0.15,2)--(2.85,2);
        \draw[->,TLink] (0.15,1.2)--(2.85,0.4);
        \draw[->,TLink] (0.15,0.4)--(2.85,1.2);
        \node[font=\scriptsize,TLinkSoft] at (0,2.6) {源 $\mu$};
        \node[font=\scriptsize,TLinkSoft] at (3,2.6) {汇 $\nu$};
      \end{tikzpicture}\par\vspace{3pt}
      {\scriptsize 每个源 $\to$ 唯一汇（一个排列）。}
    \end{column}
    \begin{column}{0.5\textwidth}
      \centering {\small\textbf{Kantorovich：质量可分}}\par\vspace{4pt}
      \begin{tikzpicture}[scale=0.82,>={Stealth[length=4pt]}]
        \foreach \y in {2,1.2,0.4}{\fill[TLprimary] (0,\y) circle (2pt);}
        \foreach \y in {2,1.2,0.4}{\fill[TLaccent] (3,\y) circle (2pt);}
        \draw[->,TLink!70,thick] (0.15,2)--(2.85,2);
        \draw[->,TLink!40] (0.15,2)--(2.85,1.2);
        \draw[->,TLink!70,thick] (0.15,1.2)--(2.85,0.4);
        \draw[->,TLink!40] (0.15,1.2)--(2.85,2);
        \draw[->,TLink!70,thick] (0.15,0.4)--(2.85,1.2);
        \node[font=\scriptsize,TLinkSoft] at (0,2.6) {源 $\mu$};
        \node[font=\scriptsize,TLinkSoft] at (3,2.6) {汇 $\nu$};
      \end{tikzpicture}\par\vspace{3pt}
      {\scriptsize 一个源的质量可\emph{分流}到多个汇（粗细 = 质量）。}
    \end{column}
  \end{columns}
  \vspace{4pt}
  \TLtakeaway{Kantorovich 把“箭头”换成“流量矩阵”：可行集变成凸的 $\coup{\mu}{\nu}$，于是最优解被紧性\emph{保送}出来。}
\end{frame}

% ── worked example：离散运输 ─────────────────────────────────────────────────
\begin{frame}{可算的例子：$3\times3$ 离散运输}
  设 $\mu=\tfrac13(\delta_{a_1}{+}\delta_{a_2}{+}\delta_{a_3})$，$\nu=\tfrac13(\delta_{b_1}{+}\delta_{b_2}{+}\delta_{b_3})$。
  耦合 $\pi=(\pi_{ij})$ 是\emph{双随机矩阵}的 $\tfrac13$ 倍：行和列和均为 $\tfrac13$。
  \begin{columns}[T]
    \begin{column}{0.52\textwidth}
      \[
        \min_{\pi\ge0}\ \sum_{i,j} c_{ij}\,\pi_{ij}
        \subjto \sum_j\pi_{ij}=\tfrac13,\ \sum_i\pi_{ij}=\tfrac13 .
      \]
      \begin{itemize}\small
        \item 这是一个\textbf{线性规划}；可行域是（缩放的）\emph{双随机矩阵}多胞形。
        \item \emph{Birkhoff 定理}：其顶点恰为\emph{置换矩阵} $\Rightarrow$ 存在\emph{最优的 Monge 映射}（一个排列）。
      \end{itemize}
    \end{column}
    \begin{column}{0.48\textwidth}
      \centering
      \[
        c=\begin{pmatrix} 0 & 2 & 2\\ 2 & 0 & 2\\ 2 & 2 & 0\end{pmatrix}
        \;\Rightarrow\;
        \pi^\star=\tfrac13\!\begin{pmatrix} 1&0&0\\0&1&0\\0&0&1\end{pmatrix}
      \]
      {\scriptsize 对角成本为 $0$，故恒等排列最优，$C(\mu,\nu)=0$。}
    \end{column}
  \end{columns}
  \vspace{2pt}
  \TLtakeaway{离散情形 = 线性规划 + Birkhoff：松弛（双随机）与原始（置换）在\emph{顶点}处重合。这预演了第 2 节“对偶”与第 9 章“Monge=Kantorovich”。}
\end{frame}

% ── 成本与对偶伏笔 ───────────────────────────────────────────────────────────
\begin{frame}{成本函数与对偶的“面包店”伏笔}
  \deflab{成本 $c(x,y)$。} “把一单位质量从 $x$ 送到 $y$ 的代价”。常见选择：
  \[
    c(x,y)=\abs{x-y}\ (\text{距离}),\qquad c(x,y)=\abs{x-y}^2\ (\text{二次，最重要}).
  \]
  最优总成本 $C(\mu,\nu)$ 由上页 Kantorovich 问题给出。

  \medskip
  \deflab{对偶直觉（第 2 节正式展开）。} 设想一家中介：在产地 $x$ 以价 $\psi(x)$ \emph{收货}，
  在销地 $y$ 以价 $\phi(y)$ \emph{售出}，只要价格“诚实”——
  \[
    \phi(y)-\psi(x)\le c(x,y)\quad(\forall x,y),
  \]
  其利润 $\int\phi\dd\nu-\int\psi\dd\mu$ 永不超过你自运的成本。\emph{最大利润} = \emph{最小成本}。
  \TLtakeaway{“成本视角”（搬多少）与“价格视角”（值多少）在最优处相等——这就是 Kantorovich 对偶，第 2 节的核心。}
\end{frame}

% ── 习题 ─────────────────────────────────────────────────────────────────────
\begin{frame}{习题 0 · Exercises（提示见本页，解答见附录）}
  \begin{itemize}
    \item \textbf{习题 0.1}~(\diffI)\ 设 $\mu=\mathrm{Unif}[0,1]$，$T(x)=2x$。求 $\nu=\pf{T}{\mu}$ 及其密度。
      \begin{itemize}\footnotesize
        \item \hint{用 $\nu[(0,y)]=\mu[\,2x<y\,]$，再对 $y$ 求导。}
      \end{itemize}
    \item \textbf{习题 0.2}~(\diffII)\ 证明：乘积测度 $\mu\otimes\nu$ 恒为一个耦合。由此论证 $\coup{\mu}{\nu}\neq\varnothing$。
      \begin{itemize}\footnotesize
        \item \hint{验证两个边际；这说明 Kantorovich 可行集\emph{总}非空。}
      \end{itemize}
    \item \textbf{习题 0.3}~(\diffIII)\ 取 $\mu=\delta_0$，$\nu=\tfrac12\delta_{-1}{+}\tfrac12\delta_{1}$，$c=\abs{x-y}^2$。
      说明 Monge 不可行，但写出唯一的 Kantorovich 耦合并算出 $C(\mu,\nu)$。
      \begin{itemize}\footnotesize
        \item \hint{$\delta_0$ 的质量必须分裂；唯一耦合是 $\tfrac12\delta_{(0,-1)}{+}\tfrac12\delta_{(0,1)}$。}
      \end{itemize}
  \end{itemize}
\end{frame}

% ── 小结 ─────────────────────────────────────────────────────────────────────
\begin{frame}{§0 小结：上路前的工具箱}
  \begin{itemize}
    \item \textbf{对象}：$\Pp(\X)$ 上的概率测度；\textbf{靠近}：弱收敛；\textbf{引擎}：Prokhorov 紧性。
    \item \textbf{动作}：推前 $\pf{T}{\mu}$（质量守恒 / 换元）。
    \item \textbf{两个问题}：Monge（确定性、可能无解）$\subset$ Kantorovich（耦合、线性、\emph{总有解}）。
    \item \textbf{伏笔}：成本视角 $\leftrightarrow$ 价格视角的对偶。
  \end{itemize}
  \TLtakeaway{记住主问题 $C(\mu,\nu)=\inf_{\pi\in\coup{\mu}{\nu}}\int c\dd\pi$。第 1 节起，我们追问：它\emph{何时取到}、\emph{长什么样}、\emph{如何随时间演化}。}
\end{frame}
      % §0  预备：测度、耦合与两个传输问题（指挥官撰写，风格范本）
% =============================================================================
%  §1  Ch.4 最优耦合的基本性质
%  来源：Villani, Optimal Transport: Old and New, Chapter 4 (pp. 43–48)
% =============================================================================
\TLsection{1 \ 最优耦合的基本性质}

% ── 动机：为什么存在性不是显然的 ───────────────────────────────────────────────
\begin{frame}{为什么最优耦合"存在"不是显然的？}
  我们已经知道 Kantorovich 问题的可行集 $\coup{\mu}{\nu}$ 非空（含乘积测度 $\mu\otimes\nu$），
  目标 $I[\pi]=\int_{\X\times\Y} c\dd\pi$ 对 $\pi$ 线性。

  \textbf{朴素想法。} 取一列"趋向下确界"的计划：存在 $(\pi_k)\subset\coup{\mu}{\nu}$ 使
  \[
    I[\pi_k] \;\To\; \inf_{\pi\in\coup{\mu}{\nu}} I[\pi].
  \]
  \textbf{问题。} $(\pi_k)$ 是否有收敛子列？极限是否还在 $\coup{\mu}{\nu}$ 里？

  \begin{columns}[T]
    \begin{column}{0.5\textwidth}
      \textbf{两个关键工具：}
      \begin{itemize}
        \item \textbf{紧性。} $\coup{\mu}{\nu}$ 在弱拓扑下紧——使得收敛子列存在。
        \item \textbf{下半连续。} $I[\pi]$ 在弱极限处不增——使得极限是最优。
      \end{itemize}
    \end{column}
    \begin{column}{0.5\textwidth}
      \textbf{对应结果：}
      \begin{itemize}
        \item Lemma 4.4 $+$ Prokhorov 定理 $\Rightarrow$ 紧性。
        \item Lemma 4.3 $\Rightarrow$ 下半连续。
        \item Theorem 4.1 组装：最优耦合存在。
      \end{itemize}
    \end{column}
  \end{columns}
  \TLtakeaway{存在性证明的骨架：下确界列 $\to$ 紧性（取极限） $\to$ 下半连续（极限仍最优）。}
\end{frame}

% ── 定理 4.1 陈述 ─────────────────────────────────────────────────────────────
\begin{frame}{定理 4.1：最优耦合的存在性}
  \deflab{定理 4.1（最优耦合存在，Villani Thm 4.1）。}

  设 $\X,\Y$ 为 Polish 空间，$\mu\in\Pp(\X)$，$\nu\in\Pp(\Y)$。
  设 $a:\X\to\R\cup\{-\infty\}$ 上半连续且 $a\in L^1(\mu)$，
  $b:\Y\to\R\cup\{-\infty\}$ 上半连续且 $b\in L^1(\nu)$。
  设成本函数 $c:\X\times\Y\to\R\cup\{+\infty\}$ 下半连续，且
  \[
    c(x,y) \;\ge\; a(x)+b(y) \quad \forall\, x\in\X,\; y\in\Y.
  \]
  则\emph{存在}一个耦合 $\pi^\star\in\coup{\mu}{\nu}$ 使得
  \[
    I[\pi^\star] = \inf_{\pi\in\coup{\mu}{\nu}} I[\pi], \quad I[\pi]:=\int_{\X\times\Y} c(x,y)\dd\pi(x,y).
  \]

  \deflab{注记 4.2。} 下界条件 $c\ge a+b$ 保证 $I[\pi]$ 在 $\R\cup\{+\infty\}$ 中\emph{有定义}（不出现 $+\infty-\infty$）。
  多数应用中可取 $a=b=0$，即 $c\ge0$。

  \TLtakeaway{定理 4.1 是"第一件好事"：在 Polish 空间上、下半连续成本满足下界条件时，最优耦合必然存在。}
\end{frame}

% ── 注记 4.5：存在性 ≠ 有限性 ────────────────────────────────────────────────
\begin{frame}{注记 4.5：存在性不等于有限最优成本}
  \deflab{注记 4.5。} 定理 4.1 断言最优耦合\emph{存在}，但不排除
  \[
    \inf_{\pi\in\coup{\mu}{\nu}} I[\pi] = +\infty.
  \]
  即：存在最优计划，但其成本可能无穷大。

  \textbf{何时成本有限？} 一个简单的充分条件：
  \[
    \iint_{\X\times\Y} c(x,y)\dd\mu(x)\dd\nu(y) < +\infty.
  \]
  此条件等价于\emph{独立耦合} $\mu\otimes\nu$ 的成本有限，从而保证至少有一个计划成本有限，
  进而最优成本有限。

  \textbf{更强的条件（后续章节用到）：} 若
  $c(x,y)\le c_{\X}(x)+c_{\Y}(y)$，其中 $c_{\X}\in L^1(\mu)$，$c_{\Y}\in L^1(\nu)$，
  则\emph{任意}耦合的成本均有限。

  \TLtakeaway{成本有限性是独立的条件：验证 $\mu\otimes\nu$ 的成本是否有限，是最简单的途径。}
\end{frame}

% ── Prokhorov 定理回顾 + 引理 4.4 ────────────────────────────────────────────
\begin{frame}{Prokhorov 定理与耦合集的胶着性（引理 4.4）}
  \deflab{Prokhorov 定理（紧性判据）。}
  设 $\X$ 为 Polish 空间，$\mathcal{F}\subset\Pp(\X)$。
  $\mathcal{F}$ 在弱拓扑下\emph{预紧}当且仅当它\TLterm{胶着 (tight)}：
  \[
    \forall\eps>0,\;\exists\text{ 紧集 } K_\eps\subset\X, \;
    \mu[\X\setminus K_\eps]\le\eps \;\forall\,\mu\in\mathcal{F}.
  \]
  直觉：无质量"逃逸到无穷"。

  \deflab{引理 4.4（转移计划的胶着性）。}
  设 $\mathcal{F}\subset\Pp(\X)$ 与 $\mathcal{G}\subset\Pp(\Y)$ 均胶着，
  则 $\Pi(\mathcal{F},\mathcal{G})$ 在 $\Pp(\X\times\Y)$ 中\emph{也胶着}。

  \textbf{证明思路。} 对任意 $\eps>0$，取紧集 $K_\eps\subset\X$，$L_\eps\subset\Y$ 使得
  $\mu[\X\setminus K_\eps]\le\eps$，$\nu[\Y\setminus L_\eps]\le\eps$。
  对任意耦合 $(X,Y)\sim\pi\in\Pi(\mathcal{F},\mathcal{G})$：
  \[
    \Prob\bigl[(X,Y)\notin K_\eps\times L_\eps\bigr]
    \;\le\; \Prob[X\notin K_\eps]+\Prob[Y\notin L_\eps]
    \;\le\; 2\eps.
  \]
  $K_\eps\times L_\eps$ 在 $\X\times\Y$ 中紧，上界 $2\eps$ 与耦合选取无关。\hfill$\square$

  \TLtakeaway{单个 $\mu,\nu$ 的胶着性（在 Polish 空间上自动成立）可"提升"到整个耦合集 $\coup{\mu}{\nu}$。}
\end{frame}

% ── 引理 4.3：成本泛函的下半连续性 ──────────────────────────────────────────
\begin{frame}{引理 4.3：成本泛函的下半连续性（陈述）}
  \deflab{引理 4.3（成本泛函的下半连续性）。}
  设 $\X,\Y$ 为 Polish 空间，$c:\X\times\Y\to\R\cup\{+\infty\}$ 下半连续，
  $h:\X\times\Y\to\R\cup\{-\infty\}$ 上半连续且 $c\ge h$。
  设 $\pi_k\wkto\pi$，且 $h\in L^1(\pi_k)$，$h\in L^1(\pi)$，
  $\int h\dd\pi_k\To\int h\dd\pi$。则
  \[
    \int_{\X\times\Y} c\dd\pi \;\le\; \liminf_{k\to\infty}\int_{\X\times\Y} c\dd\pi_k.
  \]

  \textbf{读法。} 泛函 $\pi\mapsto\int c\dd\pi$ 对弱收敛\emph{下半连续}——这正是 §0 中"下半连续"概念在测度空间 $\Pp(\X\times\Y)$ 上的体现（辅助函数 $h$ 仅用于把成本下界 $c\ge a+b$ 纳入框架）。

  \TLtakeaway{下半连续成本泛函对弱收敛也下半连续——这让"取极限不增加成本"成为可能（定理 4.1 的另一半引擎）。}
\end{frame}

% ── 引理 4.3 的证明 ──────────────────────────────────────────────────────────
\begin{frame}{引理 4.3 的证明（一）：为什么可令 $c\ge0$}
  \textbf{先解释归约。}\proofskipnote 因 $c\ge h$，函数 $c-h$ 非负；若已知引理对非负成本 $c-h$ 成立，则
  \[
    \begin{aligned}
      \int c\dd\pi
      &= \int(c-h)\dd\pi+\int h\dd\pi\\
      &\le \liminf_{k}\int(c-h)\dd\pi_k+\lim_{k}\int h\dd\pi_k\\
      &= \liminf_{k}\int c\dd\pi_k .
    \end{aligned}
  \]
  最后一行用的是假设 $\int h\dd\pi_k\To\int h\dd\pi$。所以只需证明非负下半连续情形。

  \textbf{合法性。} $h$ 的积分沿 $\pi_k\wkto\pi$ 收敛，表示辅助下界项不会在极限中丢失；
  因此给 $c-h$ 证明的下半连续不等式，可以通过加回同一个收敛的 $h$ 项传回原成本 $c$。

  \TLtakeaway{归约不是省略：$c\ge h$ 给出非负成本 $c-h$，而 $\int h\dd\pi_k\to\int h\dd\pi$ 保证加回 $h$ 后不等式仍成立。}
\end{frame}

\begin{frame}{引理 4.3 的证明：逼近 + 单调收敛}
  \deflab{补一个事实：单调收敛定理。}
  若 $0\le c_{\ell}\nearrow c$，则 $\sup_{\ell}\int c_{\ell}\dd\pi=\lim_{\ell}\int c_{\ell}\dd\pi=\int c\dd\pi$。\proofskipnote

  \textbf{非负情形。} 用 §0 下半连续性质 (a)（inf-卷积）取有界连续逼近
  $c_{\ell}\in C_b(\X\times\Y)$，$0\le c_{\ell}\nearrow c$。固定 $\ell$，弱收敛给
  $\int c_{\ell}\dd\pi_k\to\int c_{\ell}\dd\pi$，且 $c_{\ell}\le c$，故
  \[
    \int c_{\ell}\dd\pi
    = \lim_{k}\int c_{\ell}\dd\pi_k
    \le \liminf_{k}\int c\dd\pi_k
    \quad\Longrightarrow\quad
    \int c\dd\pi\le \liminf_{k}\int c\dd\pi_k .
  \]
  右端与 $\ell$ 无关，对 $\ell$ 取上确界，并用上面的单调收敛事实，即得结论。\hfill$\square$
  \TLtakeaway{非负情形的核心：lsc 给出连续递增逼近；固定 $\ell$ 用弱收敛和 $c_{\ell}\le c$，最后用单调收敛回到 $c$。}
\end{frame}

% ── 定理 4.1 的证明（一a）：预紧 ─────────────────────────────────────────────
\begin{frame}{定理 4.1 的证明（一a）：$\coup{\mu}{\nu}$ 为何\emph{预紧}}
  \textbf{总思路。}\proofskipnote 把 §0 的"紧集上 lsc 必达最小"搬到测度空间：取"紧集"$=\coup{\mu}{\nu}$、"lsc 目标"$=I[\pi]=\int_{\X\times\Y} c\dd\pi$。本帧先证可行集\emph{预紧}（弱闭包紧）。

  \deflab{补一个事实：Ulam 定理。} Polish 空间上\emph{每个}概率测度 $\mu$ 都自动胶着：完备可分度量空间里 $\mu$ 对紧集\TLterm{内正则 (inner regular)}——任给 $\eps>0$ 总有紧集 $K_\eps$ 使 $\mu[\X\setminus K_\eps]\le\eps$。故单点族 $\{\mu\},\{\nu\}$ 胶着。

  \textbf{推出预紧（三步，每步标注依据）：}
  \begin{itemize}\small
    \item $\{\mu\},\{\nu\}$ 胶着\hfill（Ulam 定理）
    \item $\Rightarrow$ $\coup{\mu}{\nu}$ 在 $\Pp(\X\times\Y)$ 中胶着\hfill（引理 4.4：两端胶着"提升"到耦合集）
    \item $\Rightarrow$ $\coup{\mu}{\nu}$ \emph{预紧}\hfill（Prokhorov：胶着 $\Leftrightarrow$ 弱预紧）
  \end{itemize}

  \TLtakeaway{预紧来自一条链：Ulam（单个测度胶着）$\to$ 引理 4.4（耦合集胶着）$\to$ Prokhorov（胶着即预紧）。}
\end{frame}

% ── 定理 4.1 的证明（一b）：闭性 → 紧，取极小化列 ────────────────────────────
\begin{frame}{定理 4.1 的证明（一b）：从预紧到\emph{紧}，并取极小化列}
  \textbf{预紧还不够：}极限未必落在 $\coup{\mu}{\nu}$ 里，故须再证它\emph{闭}（闭 $+$ 预紧 $=$ 紧）。\proofskipnote

  \deflab{闭性：把"边际是 $\mu$"传给极限。}
  设 $\pi_k\in\coup{\mu}{\nu}$，$\pi_k\wkto\pi$。任取 $\varphi\in C_b(\X)$，视 $(x,y)\mapsto\varphi(x)$ 为有界连续函数：
  \begin{align*}
    \int_{\X}\varphi\dd\big(\pf{(\mathrm{proj}_{\X})}{\pi}\big)
    &=\int_{\X\times\Y}\!\varphi(x)\dd\pi
    \;\overset{\pi_k\wkto\pi}{=}\;\lim_{k}\int_{\X\times\Y}\!\varphi(x)\dd\pi_k\\
    &=\lim_{k}\int_{\X}\varphi\dd\mu=\int_{\X}\varphi\dd\mu .
  \end{align*}
  对一切 $\varphi$ 成立 $\Rightarrow$ $\pf{(\mathrm{proj}_{\X})}{\pi}=\mu$，同理第二边际 $=\nu$，故 $\pi\in\coup{\mu}{\nu}$：集合\emph{闭}。\textbf{闭 $+$ 预紧 $\Rightarrow$ $\coup{\mu}{\nu}$ 弱紧。}

  \TLtakeaway{闭性核心只有一行：用检验函数 $\varphi(x)$ 测两边，弱收敛把"边际 $=\mu$"原样传到极限；再配预紧即得\emph{紧}。}
\end{frame}

% ── 定理 4.1 的证明（二）：下半连续锁定最优性 ────────────────────────────────
\begin{frame}{定理 4.1 的证明（二）：用下半连续\emph{锁定}最优性}
  \textbf{取极小化列。}\proofskipnote 取 $(\pi_k)$ 使 $I[\pi_k]\to\inf I$，（一b）弱紧性抽出子列 $\pi_k\wkto\pi^\star\in\coup{\mu}{\nu}$。难点：弱收敛\emph{不}保证 $I[\pi_k]\to I[\pi^\star]$；但 lsc（引理 4.3）给的\emph{一个方向}够用。

  \deflab{先核对引理 4.3 的前提。} 取 $h(x,y)=a(x)+b(y)$（下界条件 $c\ge a+b$ 中的 $h$，上半连续）。\textbf{关键：}所有 $\pi_k,\pi^\star$ 共享边际 $\mu,\nu$，故
  \[
    \int h\dd\pi_k=\int_{\X} a\dd\mu+\int_{\Y} b\dd\nu=\int h\dd\pi^\star .
  \]
  右端与 $k$ 无关（\emph{常数列}），故前提"$\int h\dd\pi_k\to\int h\dd\pi^\star$"自动成立，引理 4.3 适用。

  \textbf{收尾（夹逼）：}
  \begin{itemize}\small
    \item 引理 4.3 $\Rightarrow$ $I[\pi^\star]\le\liminf_k I[\pi_k]=\inf I$；
    \item 可行性 $\Rightarrow$ $I[\pi^\star]\ge\inf I$；夹逼得 $I[\pi^\star]=\inf I$，即 $\pi^\star$ \emph{最优}。\hfill$\square$
  \end{itemize}

  \TLtakeaway{lsc 只需"成本不会向上跳"这一个方向 $I[\pi^\star]\le\liminf I[\pi_k]=\inf I$，配可行性即夹出等号。}
\end{frame}

% ── 定理 4.6：最优性被子计划继承 ─────────────────────────────────────────────
\begin{frame}{第二件好事：最优性被子计划继承（定理 4.6）}
  \deflab{定理 4.6（限制性质，Villani Thm 4.6）。}
  设定同定理 4.1，并设最优传输成本 $C(\mu,\nu)<+\infty$，$\pi\in\coup{\mu}{\nu}$ 为最优计划。
  设 $\pi'$ 是 $\X\times\Y$ 上的非负测度，满足 $\pi'\le\pi$（作为测度）且 $\pi'[\X\times\Y]>0$。
  定义归一化限制：
  \[
    \tilde{\pi} \;:=\; \frac{\pi'}{\pi'[\X\times\Y]}.
  \]
  则 $\tilde{\pi}$ 是其边际 $\tilde{\mu}=\pf{(\mathrm{proj}_{\X})}{\tilde{\pi}}$，
  $\tilde{\nu}=\pf{(\mathrm{proj}_{\Y})}{\tilde{\pi}}$ 之间的\emph{最优}耦合。

  \textbf{通俗理解。} "最优传输计划的任意一个子计划，对其所搬运的那部分质量来说，也是最优的。"
  否则可以局部改进子计划，从而降低整体成本——矛盾。

  \TLtakeaway{限制性质说明：最优计划\emph{局部也优}。这是全书反复使用的简单而有力的原则。}
\end{frame}

% ── 定理 4.6 的矛盾证明（一）：构造与验证 ────────────────────────────────────
\begin{frame}{定理 4.6 的证明（一）：构造竞争计划}
  \textbf{反证。}\proofskipnote 令 $Z=\pi'[\X\times\Y]>0$，$\tilde{\pi}=\pi'/Z$。若 $\tilde{\pi}$ 不是最优的，
  则存在同边际计划 $\hat{\pi}\in\coup{\tilde{\mu}}{\tilde{\nu}}$，满足 $I[\hat{\pi}]<I[\tilde{\pi}]$。

  \deflab{补一个事实：同边际计划的差有零边际。}
  因 $\hat{\pi}$ 与 $\tilde{\pi}$ 共享边际 $(\tilde{\mu},\tilde{\nu})$，所以
  \[
    \pf{(\mathrm{proj}_{\X})}{(\hat{\pi}-\tilde{\pi})}=0,\quad
    \pf{(\mathrm{proj}_{\Y})}{(\hat{\pi}-\tilde{\pi})}=0,\qquad
    \hat{\pi}_{\mathrm{new}}:=(\pi-\pi')+Z\hat{\pi}=\pi+Z(\hat{\pi}-\tilde{\pi}).
  \]
  这里 $\pi-\pi'$ 非负（因为 $\pi'\le\pi$），$Z\hat{\pi}$ 非负，故 $\hat{\pi}_{\mathrm{new}}$ 是非负测度。
  又因上式右端的差项有零边际，$\hat{\pi}_{\mathrm{new}}$ 的两个边际仍是 $\mu,\nu$，即它是 $\pi$ 的竞争计划。

  \textbf{成本严格下降。} $C(\mu,\nu)<+\infty$ 保证下面的加减运算合法；于是
  \[
    I[\hat{\pi}_{\mathrm{new}}]
    = I[\pi]-Z I[\tilde{\pi}]+Z I[\hat{\pi}]
    < I[\pi].
  \]
  这与 $\pi$ 最优性矛盾，故 $\tilde{\pi}$ 是最优的。\hfill$\square$
  \TLtakeaway{隐藏步骤是边际核对：把新计划写成 $\pi+Z(\hat{\pi}-\tilde{\pi})$，同边际差项的边际为零，所以替换只改变局部成本，不改变全局边际。}
\end{frame}

% ── 定理 4.6 的证明（二）：唯一性的传递 ─────────────────────────────────────
\begin{frame}{定理 4.6 的证明（二）：唯一性的传递}
  \textbf{目标。}\proofskipnote 已知 $\pi$ 是 $(\mu,\nu)$ 的\emph{唯一}最优计划；要证 $\tilde{\pi}$ 是
  $(\tilde{\mu},\tilde{\nu})$ 的唯一最优计划。
  \textbf{逐步传递唯一性。}
  \begin{itemize}\small
    \item 任取一个 $(\tilde{\mu},\tilde{\nu})$ 的最优计划 $\hat{\pi}$。
    \item 仍令 $\hat{\pi}_{\mathrm{new}}=(\pi-\pi')+Z\hat{\pi}$。上一帧已验证：它是 $(\mu,\nu)$ 的竞争计划。
    \item 因 $\hat{\pi}$ 与 $\tilde{\pi}$ 都最优，它们成本相同，故
  \end{itemize}
  \[
    I[\hat{\pi}_{\mathrm{new}}]
    = I[\pi]-Z I[\tilde{\pi}]+Z I[\hat{\pi}]
    = I[\pi].
  \]
  所以 $\hat{\pi}_{\mathrm{new}}$ 也是 $(\mu,\nu)$ 的最优计划；由 $\pi$ 的唯一性，
  $\hat{\pi}_{\mathrm{new}}=\pi$。消去公共部分得 $Z\hat{\pi}=\pi'=Z\tilde{\pi}$，
  又 $Z>0$，故 $\hat{\pi}=\tilde{\pi}$。\hfill$\square$
  \textbf{例 4.7（条件化仍最优）。}
  设 $(X,Y)$ 是最优计划 $\pi$ 下的随机向量，$Z\subset\X\times\Y$ Borel 且 $\Prob[(X,Y)\in Z]>0$。
  令 $\pi' = \pi|_{Z}$，则 $\pi'\le\pi$，
  定理 4.6 直接给出：给定 $(X,Y)\in Z$ 的条件耦合在相应边际之间最优。
  \TLtakeaway{唯一性传递的关键等式是 $I[\hat{\pi}_{\mathrm{new}}]=I[\pi]$：整体唯一性迫使局部替代等于原局部计划。}
\end{frame}

% ── 例 4.7 的原则解释 ──────────────────────────────────────────────────────────
\begin{frame}{例 4.7：条件化原则的威力}
  % Manual warn-style block (avoids 0.4pt draw-border overflow from \TLwarnblock)
  \begin{tikzpicture}
    \node[draw=TLneg, fill=TLneg!8, rounded corners=3pt,
          inner sep=8pt, text width=\linewidth-17pt, align=justify] {%
      \small\textbf{\color{TLneg}为什么这个原则"简单但有力"}\\[0.3em]
      全书多次用到"取一个最优计划，把注意力限制到某个子集，
      子集上的计划仍最优"这一论证步骤。每次背后都是定理 4.6。
    };
  \end{tikzpicture}

  \textbf{使用场景举例：}
  \begin{itemize}
    \item 把最优计划限制到某个"好的"区域，在该区域上研究其几何结构。
    \item 用条件化把全局问题降维为局部问题，再用局部最优性归纳推进。
    \item 在唯一性证明中：若存在两个最优计划，对它们的"差集"条件化推出矛盾。
  \end{itemize}

  \TLtakeaway{例 4.7：最优耦合对任意正概率事件的条件化也最优（定理 4.6 之概率语言形式）。}
\end{frame}

% ── 定理 4.8：最优成本的凸性 ──────────────────────────────────────────────────
\begin{frame}{第三件好事：最优成本的凸性（定理 4.8）}
  \deflab{定理 4.8（最优成本的凸性，Villani Thm 4.8）。}
  设 $(\Theta,\lambda)$ 为概率空间，$\theta\mapsto\mu_{\theta}\in\Pp(\X)$，
  $\theta\mapsto\nu_{\theta}\in\Pp(\Y)$ 可测，成本 $c\ge a+b$ 如定理 4.1。记混合边际
  $\mu=\int_{\Theta}\mu_{\theta}\dd\lambda(\theta)$，
  $\nu=\int_{\Theta}\nu_{\theta}\dd\lambda(\theta)$。则
  \[
    C(\mu,\nu)
    \;\le\;
    \int_{\Theta} C(\mu_{\theta},\nu_{\theta})\dd\lambda(\theta).
  \]

  \textbf{直觉。} "先混合，再最优传输"的代价不超过"先分别最优传输，再混合"的平均代价。
  混合只会\emph{增加灵活性}（可以选混合后的最优计划），故成本不增。

  \deflab{读法。} 这不是说每个最优计划线性变化；它只比较\emph{最优值}。右端给出一个可行的混合方案，
  左端允许在混合后的更大可行集中重新优化。

  \TLtakeaway{最优传输成本关于边际测度\emph{凸}：混合分布的传输成本不超过分别传输成本的平均。}
\end{frame}

\begin{frame}{定理 4.8 的证明（一）：混合计划有正确边际}
  \textbf{选取局部最优计划。}\proofskipnote 对 $\lambda$-a.e.\ $\theta$，定理 4.1 给出
  $\pi_{\theta}\in\coup{\mu_{\theta}}{\nu_{\theta}}$ 且
  $I[\pi_{\theta}]=C(\mu_{\theta},\nu_{\theta})$。
  Villani 暂时借用推论 5.22，保证可把 $\theta\mapsto\pi_{\theta}$ 选成可测。

  \deflab{补一个事实：测度值积分的边际。}
  令 $\pi=\int_{\Theta}\pi_{\theta}\dd\lambda(\theta)$。任取 $\varphi\in C_b(\X)$，把边际推出积分号：
  \[
    \begin{aligned}
      \int_{\X}\varphi\dd\pf{(\mathrm{proj}_{\X})}{\pi}
      &= \int_{\Theta}\!\int_{\X}\varphi\dd\pf{(\mathrm{proj}_{\X})}{\pi_{\theta}}\dd\lambda(\theta)\\
      &= \int_{\Theta}\!\int_{\X}\varphi\dd\mu_{\theta}\dd\lambda(\theta)
       = \int_{\X}\varphi\dd\mu .
    \end{aligned}
  \]
  故第一边际为 $\mu$；同理第二边际为 $\nu$，所以 $\pi\in\coup{\mu}{\nu}$。

  \TLtakeaway{边际核对要用检验函数：把 $\mathrm{proj}_{\X}$ 的边际积分推过 $\theta$-积分，得到混合后的第一边际 $\mu$；第二边际同理。}
\end{frame}

\begin{frame}{定理 4.8 的证明（二）：Fubini 给出成本上界}
  \textbf{用上帧的可行计划估计最优值。}\proofskipnote 已知
  $\pi=\int_{\Theta}\pi_{\theta}\dd\lambda(\theta)\in\coup{\mu}{\nu}$，所以
  $C(\mu,\nu)\le I[\pi]$。

  \deflab{补一个事实：Fubini 定理。}
  在本处它允许交换参数 $\theta$ 与位置 $(x,y)$ 的积分次序，把混合计划的总成本拆成局部计划成本的平均。

  \[
    \begin{aligned}
      C(\mu,\nu)
      &\le I[\pi]
       = \int_{\X\times\Y} c\dd\pi\\
      &= \int_{\Theta}\!\left(\int_{\X\times\Y} c\dd\pi_{\theta}\right)\dd\lambda(\theta)
       = \int_{\Theta} C(\mu_{\theta},\nu_{\theta})\dd\lambda(\theta).
    \end{aligned}
  \]
  最后一行用 $I[\pi_{\theta}]=C(\mu_{\theta},\nu_{\theta})$，这是每个 $\pi_{\theta}$ 的最优性。
  \TLtakeaway{成本上界只需一个可行计划：混合局部最优计划，再用 Fubini 把其成本拆回局部最优成本的平均。}
\end{frame}

% ── 最优计划的刻画：开放问题（一）───────────────────────────────────────────
\begin{frame}{超越存在性：最优计划的刻画（一）}
  存在性定理告诉我们"最优计划存在"，但对其\emph{结构}只字未提。自然的问题：

  \begin{itemize}
    \item \textbf{唯一性？} 最优计划是否唯一？是否光滑？
    \item \textbf{Monge 耦合？} 是否存在\emph{确定性}最优计划（即推前映射 $T$ 使 $\pf{T}{\mu}=\nu$）？
    \item \textbf{几何刻画？} 有没有直接判断某个耦合"是否最优"的几何条件？
  \end{itemize}

  关于唯一性与几何刻画的答案见第 9--12 章；关于几何刻画的关键工具——循环单调性与 Kantorovich 对偶——将在下一节（§2）展开。
\end{frame}

% ── 最优计划的刻画：开放问题（二）───────────────────────────────────────────
\begin{frame}{超越存在性：为何 Monge 不能用紧性论证（二）}
  \textbf{为什么不能直接对 Monge 用定理 4.1 的紧性论证？}

  Monge 计划的集合（由推前映射产生的耦合）\emph{一般不弱紧}——事实上，它往往在全体耦合 $\coup{\mu}{\nu}$ 中\emph{稠密}！
  因此：
  \begin{itemize}
    \item Monge 问题的下确界 $=$ Kantorovich 问题的最小值（两者相等是\emph{可能}的）。
    \item 但\emph{无法保证} Monge 极小化子存在——下确界可以不被 Monge 计划取到。
  \end{itemize}

  \textbf{对比：}
  \begin{itemize}\small
    \item Kantorovich 松弛允许"分裂质量"，集合 $\coup{\mu}{\nu}$ 弱紧，极小化子必然存在。
    \item Monge 限制不允许分裂，集合不紧，极小化子可以不存在。
  \end{itemize}

  \TLtakeaway{确定性 (Monge) 计划集不紧，因此 Monge 极小子不保证存在；但 Kantorovich 最小值 $=$ Monge 下确界仍可能成立。第 5 章开始刻画。}
\end{frame}

% ── 例 4.9：质量分裂（TikZ 图 4.1）── 文字 ──────────────────────────────────
\begin{frame}{例 4.9：需要分裂质量的例子}
  \deflab{例 4.9（Villani Example 4.9）。}
  取 $\X=\Y=\R^2$，成本 $c(x,y)=\abs{x-y}^2$。
  设：
  \begin{align*}
    \mu &= H^1 \text{ 限制到 } \{0\}\times[-1,1], \\
    \nu &= \tfrac{1}{2}H^1 \text{ 限制到 } \{-1\}\times[-1,1]
           + \tfrac{1}{2}H^1 \text{ 限制到 } \{1\}\times[-1,1].
  \end{align*}
  其中 $H^1$ 为一维 Hausdorff 测度（弧长测度）。

  \textbf{唯一最优计划：} 对每个点 $(0,a)\in\supp\mu$，将其一半质量送到 $(-1,a)$，
  另一半送到 $(1,a)$（水平运输，竖直坐标不变）：
  \[
    \pi^\star = \int_{-1}^{1} \tfrac{1}{2}\bigl[\delta_{((0,a),(-1,a))}+\delta_{((0,a),(1,a))}\bigr]\dd a.
  \]

  \textbf{这\emph{不是} Monge 计划：} 每个 $(0,a)$ 的质量被分裂成两份，
  不存在函数 $T$ 使 $T(0,a)=(-1,a)$ 且同时 $T(0,a)=(1,a)$。

  \TLtakeaway{质量分裂是最优的——这说明即使有唯一最优 Kantorovich 计划，也不一定是 Monge 计划。}
\end{frame}

% ── 例 4.9：TikZ 图 4.1 ────────────────────────────────────────────────────
\begin{frame}{例 4.9 图示：最优耦合 vs 确定性近似（图 4.1 重绘）}
  \begin{columns}[T]
    \begin{column}{0.5\textwidth}
      \centering{\small\textbf{最优耦合（质量分裂）}}\\[6pt]
      \begin{tikzpicture}[scale=1.05,>={Stealth[length=4pt]}]
        % 支撑集
        \draw[TLprimary,very thick] (0,-1.3) -- (0,1.3);
        \draw[TLaccent,thick] (-1.5,-1.3) -- (-1.5,1.3);
        \draw[TLaccent,thick] (1.5,-1.3) -- (1.5,1.3);
        % 箭头：中心 -> 左和右（几个示例点）
        \foreach \y in {-1.0,-0.5,0.0,0.5,1.0}{
          \draw[->,TLink,thin] (0.05,\y) -- (1.45,\y);
          \draw[->,TLink,thin] (-0.05,\y) -- (-1.45,\y);
        }
        % 标注
        \node[font=\scriptsize,TLinkSoft,above] at (0,1.35) {$\mu$};
        \node[font=\scriptsize,TLinkSoft,above] at (-1.5,1.35) {$\nu_L$};
        \node[font=\scriptsize,TLinkSoft,above] at (1.5,1.35) {$\nu_R$};
        \node[font=\scriptsize] at (0,-1.6) {$x=0$};
        \node[font=\scriptsize] at (-1.5,-1.6) {$x=-1$};
        \node[font=\scriptsize] at (1.5,-1.6) {$x=1$};
      \end{tikzpicture}\\[3pt]
      {\scriptsize 每点质量各分一半，水平运输。}
    \end{column}
    \begin{column}{0.5\textwidth}
      \centering{\small\textbf{确定性近似（非最优）}}\\[6pt]
      \begin{tikzpicture}[scale=1.05,>={Stealth[length=4pt]}]
        \draw[TLprimary,very thick] (0,-1.3) -- (0,1.3);
        \draw[TLaccent,thick] (-1.5,-1.3) -- (-1.5,1.3);
        \draw[TLaccent,thick] (1.5,-1.3) -- (1.5,1.3);
        % 交替箭头：显式指定每个点的去向
        \draw[->,TLink,thin] (0.05,-1.0) -- (1.45,-1.0);
        \draw[->,TLink,thin] (-0.05,-0.5) -- (-1.45,-0.5);
        \draw[->,TLink,thin] (0.05,0.0) -- (1.45,0.0);
        \draw[->,TLink,thin] (-0.05,0.5) -- (-1.45,0.5);
        \draw[->,TLink,thin] (0.05,1.0) -- (1.45,1.0);
        \node[font=\scriptsize,TLinkSoft,above] at (0,1.35) {$\mu$};
        \node[font=\scriptsize,TLinkSoft,above] at (-1.5,1.35) {$\nu_L$};
        \node[font=\scriptsize,TLinkSoft,above] at (1.5,1.35) {$\nu_R$};
        \node[font=\scriptsize] at (0,-1.6) {$x=0$};
        \node[font=\scriptsize] at (-1.5,-1.6) {$x=-1$};
        \node[font=\scriptsize] at (1.5,-1.6) {$x=1$};
      \end{tikzpicture}\\[3pt]
      {\scriptsize 每点只送一侧（交替），不分裂，非最优。}
    \end{column}
  \end{columns}
  \vspace{4pt}
  {\scriptsize 据 Villani 图 4.1 重绘。}
  \TLtakeaway{最优计划（左）把质量对称分裂；确定性近似（右）不分裂但成本更高。Monge 下确界 $=$ Kantorovich 最小值，但 Monge 极小子不存在。}
\end{frame}

% ── 习题 1 · Exercises ─────────────────────────────────────────────────────────
\begin{frame}{习题 1 · Exercises}
  \begin{itemize}
    \item \textbf{习题 1.1}~(\diffI)\ 验证独立耦合 $\mu\otimes\nu$ 满足 Kantorovich 问题的边际约束，
      并在 $\iint c\dd\mu\dd\nu<+\infty$ 时说明 $I[\mu\otimes\nu]<+\infty$，
      从而 $C(\mu,\nu)<+\infty$。
      \begin{itemize}\footnotesize
        \item \hint{验证 $\pf{(\mathrm{proj}_{\X})}{\mu\otimes\nu}=\mu$ 和 $\pf{(\mathrm{proj}_{\Y})}{\mu\otimes\nu}=\nu$，
          再利用 Fubini 定理计算 $I[\mu\otimes\nu]=\iint c\dd\mu\dd\nu$。}
      \end{itemize}
    % SOLUTION: mu otimes nu 的第一边际：对任意 Borel A 属于 X，mu otimes nu [A × Y] = mu[A] * nu[Y] = mu[A]，故第一边际为 mu；同理第二边际为 nu。I[mu otimes nu] = int_{X×Y} c d(mu⊗nu) = int_X int_Y c(x,y) dnu(y) dmu(x) (Fubini，c>=0 可以交换) = ∬ c dmu dnu < +∞ 由假设。故独立耦合有限成本，C(μ,ν) ≤ I[μ⊗ν] < +∞。

    \item \textbf{习题 1.2}~(\diffII)\ 设 $\pi^\star$ 是有限离散分布的最优 Kantorovich 计划。
      取子集 $Z=\{(a_i,b_j):\pi^\star_{ij}>0\}$ 中某个矩形子块 $A\times B$，
      令 $\pi'=\pi^\star|_{A\times B}$。用显式"重新分配"论证证明 $\pi'/\pi'[A\times B]$ 是最优的。
      \begin{itemize}\footnotesize
        \item \hint{假设子计划不最优，构造更便宜的局部替换，再代入整体计划，与 $\pi^\star$ 最优性矛盾。}
      \end{itemize}
    % SOLUTION: 设子计划 tilde{pi} = pi'/(pi'[A×B]) 不是最优的。则存在有相同边际但成本更小的计划 hat{pi}。令 Z = pi'[A×B]，构造 pi_new = (pi^star - pi') + Z * hat{pi}。验证：pi_new 的边际等于 mu, nu（因为 pi' 和 hat{pi} 有相同的边际，差了一个因子 Z）；pi_new 的成本 = I[pi^star] - I[pi'] + Z * I[hat{pi}] = I[pi^star] - Z * I[tilde{pi}] + Z * I[hat{pi}] < I[pi^star]（因为 I[hat{pi}] < I[tilde{pi}]）。这与 pi^star 最优性矛盾。

    \item \textbf{习题 1.3}~(\diffIII)\ 用例 4.9 说明：存在 $\mu,\nu\in\Pp(\R^2)$ 使得
      Monge 下确界等于 Kantorovich 最小值，但不存在 Monge 映射 $T$ 实现 Kantorovich 最小值。
      \begin{itemize}\footnotesize
        \item \hint{例 4.9 中唯一最优 Kantorovich 计划分裂质量，不对应任何映射；
          但确定性近似序列的成本趋向 $C(\mu,\nu)$，说明两个下确界相等。}
      \end{itemize}
    % SOLUTION: 例 4.9 中 C(mu,nu) 由唯一最优计划 pi^star 达到，其值为 1（每点 (0,a) 到 (±1,a) 的平方距离均为 1，总成本为 int_{-1}^{1} 1 da = 2，但归一化后为 1）。任意 Monge 计划 T 满足 pf{T}{mu} = nu，必将 [−1,1] 的一半质量送到 x=−1 段，另一半送到 x=1 段；设在集合 A 上 T(0,a) = (−1,a)，在 [−1,1]\A 上 T(0,a)=(1,a)，A 测度为 1。任意这样的 T 成本等于 1 = C(mu,nu)，但没有一个是连续映射或弱极限——因为 pi^star 不是推前测度。实际上可以验证 pi^star 的条件分布 pi^star_{(0,a)} = 1/2 delta_{(-1,a)} + 1/2 delta_{(1,a)} 不是 delta 测度，故 pi^star 不是任何映射的图。但 Monge 下确界确实等于 Kantorovich 最小值（均为 1），因为上述显式 Monge 映射（取 A 任意测度为 1 的子集）成本恰好也是 1。此例说明：Kantorovich 最小值可以被 Monge 计划的成本序列实现，但极限（唯一最优 pi^star）本身不是 Monge 计划。
  \end{itemize}
\end{frame}

% ── §1 小结 ───────────────────────────────────────────────────────────────────
\begin{frame}{§1 小结：最优耦合的四个基本性质}
  \begin{itemize}\small
    \item \textbf{存在性（定理 4.1）。} Polish 空间上下半连续成本满足下界条件时，最优耦合必然存在。
      证明骨架：Prokhorov 紧性 $+$ 成本泛函下半连续。
    \item \textbf{成本有定义（注记 4.2 / 4.5）。} 下界条件保证 $I[\pi]$ 在 $\R\cup\{+\infty\}$ 有意义；
      独立耦合成本有限 $\iff$ $\iint c\dd\mu\dd\nu<+\infty$。
    \item \textbf{限制性质（定理 4.6）。} 最优计划的任意子计划也最优——"局部也优"。
      条件化同样保持最优性。
    \item \textbf{凸性（定理 4.8）。} 最优成本 $C(\mu,\nu)$ 关于边际测度凸：
      混合分布的传输成本不超过分别传输成本的平均。
  \end{itemize}
  \textbf{遗留问题（§2 起逐步回答）：}
  \begin{itemize}\small
    \item 唯一性？几何刻画？$\to$ §2 循环单调性与 Kantorovich 对偶。
    \item 确定性（Monge）最优计划是否存在？$\to$ 第 9--12 章。
  \end{itemize}
  \TLtakeaway{第 4 章打下了基础：存在性、有限性、限制性、凸性——四块基石，为后续刻画最优计划做好准备。}
\end{frame}
       % §1  Ch.4 最优耦合的基本性质
% =============================================================================
%  §2  Ch.5 循环单调性与 Kantorovich 对偶
%  来源：Villani, Optimal Transport: Old and New, Chapter 5 (pp. 51–92)
% =============================================================================
\TLsection{2 \ 循环单调性与 Kantorovich 对偶}

% ── 本节要回答什么问题 ──────────────────────────────────────────────────────────
\begin{frame}{本节：最优传输的两个核心工具}
  第 1 节证明了最优耦合\emph{存在}，但对其\textbf{结构}只字未提。本节引入两个互补视角：

  \begin{itemize}
    \item \textbf{循环单调性}（几何视角）：最优计划的支撑必须满足的"不可改进"条件。
    \item \textbf{Kantorovich 对偶}（价格视角）：用"竞争价格对"刻画最小成本，
      给出等价的对偶极大化问题。
  \end{itemize}

  \deflab{本章主结论（定理 5.10）。} 在温和条件下：
  \[
    \min_{\pi\in\coup{\mu}{\nu}}\int_{\X\times\Y} c\dd\pi
    \;=\;
    \sup\Bigl\{\int_{\Y}\ctr{\psi}\dd\nu - \int_{\X}\psi\dd\mu \;:\; \psi\;\text{c-凸}\Bigr\},
  \]
  且极小化器与极大化器均存在，并通过\TLterm{c-次微分}相互刻画。

  \TLtakeaway{本节的核心等式：最优传输成本 $=$ Kantorovich 对偶问题最优值。这是全书第一章的终点，也是后续所有几何应用的出发点。}
\end{frame}

% ── 动机：面包店 / 咖啡馆的循环改进 ─────────────────────────────────────────
\begin{frame}{动机：循环重路由——哪些计划可以改进？（一）}
  \textbf{场景。} 你负责把面包从面包店 $x$ 运到咖啡馆 $y$，已写好一份转运计划。
  有人抱怨成本太高，你决定试试\emph{局部改进}。

  \medskip
  \textbf{改进方法（循环重路由）。}
  \begin{enumerate}\small
    \item 选面包店 $x_1$，它目前向远处咖啡馆 $y_1$ 运货。
    \item 把这一篮面包改道到更近的 $y_2$：节省 $c(x_1,y_2)-c(x_1,y_1)$。
    \item 但 $y_2$ 现在多了面包，再把一篮从面包店 $x_2$ 改道到 $y_3$……
    \item ……直到最终从某个 $x_N$ 把面包改送回 $y_1$，形成一个\emph{闭环}。
  \end{enumerate}

  \medskip
  \deflab{可改进的条件。} 上述重路由使总成本严格降低，当且仅当：
  \[
    c(x_1,y_2)+c(x_2,y_3)+\cdots+c(x_N,y_1)
    \;<\;
    c(x_1,y_1)+c(x_2,y_2)+\cdots+c(x_N,y_N).
  \]

  \TLtakeaway{若支撑中存在可改进的循环，计划\emph{绝对不是最优的}。反之，"无循环可改进"是最优的必要条件——这就是循环单调性。}
\end{frame}

% ── 动机：图示（图 5.1 重绘）─────────────────────────────────────────────────
\begin{frame}{动机：循环重路由图示（图 5.1）}
  \begin{columns}[T]
    \begin{column}{0.55\textwidth}
      \begin{tikzpicture}[scale=0.95,>={Stealth[length=5pt]}]
        % 面包店节点（左侧，圆形）
        \foreach \i/\y in {1/2.4, 2/1.2, 3/0.0}{
          \node[draw=TLprimary,fill=TLprimaryTint,circle,minimum size=16pt,
                inner sep=1pt,font=\scriptsize] (x\i) at (0,\y) {$x_{\i}$};
        }
        % 咖啡馆节点（右侧，矩形）
        \foreach \j/\y in {1/2.4, 2/1.2, 3/0.0}{
          \node[draw=TLaccent,fill=TLaccentLight,rectangle,minimum size=16pt,
                inner sep=2pt,font=\scriptsize] (y\j) at (4.2,\y) {$y_{\j}$};
        }
        % 原始计划（实线）
        \draw[->,TLink,thick] (x1) -- (y1) node[midway,above,font=\scriptsize,TLinkSoft]{原};
        \draw[->,TLink,thick] (x2) -- (y2);
        \draw[->,TLink,thick] (x3) -- (y3);
        % 重路由循环（虚线，橙色）
        \draw[->,TLaccent,dashed,thick] (x1) to[bend left=18] (y2)
              node[near end,above,font=\scriptsize,TLaccent]{改};
        \draw[->,TLaccent,dashed,thick] (x2) to[bend right=18] (y3);
        \draw[->,TLaccent,dashed,thick] (x3) to[bend right=25] (y1);
        % 图例
        \node[font=\scriptsize,TLinkSoft] at (0,-0.6) {面包店};
        \node[font=\scriptsize,TLinkSoft] at (4.2,-0.6) {咖啡馆};
      \end{tikzpicture}
    \end{column}
    \begin{column}{0.42\textwidth}
      \textbf{图例说明：}
      \begin{itemize}\small
        \item 实线 = 原始计划（$x_i\to y_i$）。
        \item 虚线 = 重路由后的循环：$x_1\to y_2,\ x_2\to y_3,\ x_3\to y_1$。
        \item 若虚线总成本 $<$ 实线总成本，原计划可被改进。
        \item 循环单调的计划：对\emph{所有}有限循环，重路由均不能降低成本。
      \end{itemize}
      \vspace{6pt}
      {\scriptsize 据 Villani 图 5.1 重绘。}
    \end{column}
  \end{columns}
  \TLtakeaway{若存在一个循环使重路由更便宜，原计划必然\emph{不最优}。循环单调性精确刻画了"无利可图的循环"。}
\end{frame}

% ── 定义 5.1：c-循环单调性 ────────────────────────────────────────────────────
\begin{frame}{定义 5.1：c-循环单调集与 c-循环单调计划}
  \deflab{定义 5.1（c-循环单调性，Villani Def 5.1）。}
  设 $\X,\Y$ 为任意集合，$c:\X\times\Y\to(-\infty,+\infty]$。
  子集 $\Gamma\subset\X\times\Y$ 称为\TLterm{c-循环单调的}，若对任意 $N\in\N$
  与任意 $(x_1,y_1),\ldots,(x_N,y_N)\in\Gamma$，均有
  \[
    \sum_{i=1}^{N} c(x_i,y_i) \;\le\; \sum_{i=1}^{N} c(x_i,y_{i+1}),
    \qquad (\text{约定 }y_{N+1}=y_1).
  \]
  一个转运计划称为 c-循环单调的，若它集中（concentrated）在某个 c-循环单调集上。

  \begin{itemize}\small
    \item \textbf{直觉：} c-循环单调计划是"无法通过重路由有限次循环来节省成本"的计划——某种意义上的\emph{局部极小化器}。
    \item \textbf{必要性（直觉）：} 最优计划必然 c-循环单调。若不然，存在一个更便宜的循环重路由，与最优性矛盾。
    \item \textbf{充分性（非显然）：} 反过来，c-循环单调是否保证最优？这正是定理 5.10 的核心内容。
  \end{itemize}

  \TLtakeaway{c-循环单调性 = "无法通过循环重排来节省总成本"。它是最优计划的\emph{必要}条件，在温和假设下也是\emph{充分}条件（定理 5.10）。}
\end{frame}

% ── 1D 例子：c-循环单调性与单调重排 ─────────────────────────────────────────
\begin{frame}{可算的例子：1D 中 $c=|x-y|^2$ 的 c-循环单调性}
  \textbf{设定。} $\X=\Y=\R$，$c(x,y)=|x-y|^2$，转运计划 $\pi$。

  \medskip
  \deflab{命题（1D 等价）。} 在此成本下，$\pi$ c-循环单调当且仅当其支撑是\emph{单调不减}的子集（即：$(x_1,y_1),(x_2,y_2)\in\supp\pi$，$x_1<x_2$ 则 $y_1\le y_2$）。

  \medskip
  \textbf{证明简述（$N=2$ 情形）。} 对 $(x_1,y_1),(x_2,y_2)\in\supp\pi$，c-循环单调要求：
  \[
    |x_1-y_1|^2+|x_2-y_2|^2 \;\le\; |x_1-y_2|^2+|x_2-y_1|^2.
  \]
  展开化简得 $(x_1-x_2)(y_1-y_2)\ge0$，即 $x_1,x_2$ 与 $y_1,y_2$ 的大小关系相同——这正是单调性。

  \medskip
  \deflab{推论。} 对 $c=|x-y|^2$，唯一的最优（c-循环单调）计划由\TLterm{单调重排}给出：
  若 $\mu,\nu$ 为 $\R$ 上的概率测度，最优传输映射为 $T=F_{\nu}^{-1}\circ F_{\mu}$（分布函数的反函数复合）。

  \TLtakeaway{1D 二次成本：c-循环单调 $\Leftrightarrow$ 支撑单调不减 $\Rightarrow$ 最优映射 = 单调重排 $T=F_\nu^{-1}\circ F_\mu$。}
\end{frame}

% ── 对偶 Kantorovich 问题：价格视角 ─────────────────────────────────────────
\begin{frame}{价格视角：对偶 Kantorovich 问题（一）}
  \textbf{经济直觉（Villani 的"中介公司"类比）。}
  设一家中介公司愿意接管所有运输，以价格 $\psi(x)$ 在面包店 $x$ 收购面包，
  以价格 $\phi(y)$ 在咖啡馆 $y$ 出售。为保持\emph{竞争性}，必须满足：
  \[
    \phi(y) - \psi(x) \;\le\; c(x,y) \qquad \forall\, x\in\X,\; y\in\Y.
  \]
  中介的利润为 $\int_{\Y}\phi\dd\nu - \int_{\X}\psi\dd\mu$；他们希望\emph{最大化}利润。

  \medskip
  \deflab{对偶 Kantorovich 问题（式 5.3）。}
  \[
    \sup\left\{\int_{\Y}\phi\dd\nu - \int_{\X}\psi\dd\mu \;:\;
    \phi(y)-\psi(x)\le c(x,y),\; \psi\in L^1(\mu),\; \phi\in L^1(\nu)\right\}.
  \]

  \TLtakeaway{"成本视角"（搬多少钱）与"价格视角"（值多少钱）通过对偶相联系：中介能争取的最大利润 $=$ 自运的最小成本。}
\end{frame}

% ── 弱对偶不等式 ────────────────────────────────────────────────────────────
\begin{frame}{价格视角：弱对偶不等式（命题）}
  \deflab{命题（弱对偶，Villani 式 5.4）。}
  对任意竞争价格对 $\phi(y)-\psi(x)\le c(x,y)$ 以及任意耦合 $\pi\in\coup{\mu}{\nu}$，
  \[
    \int_{\Y}\phi\dd\nu - \int_{\X}\psi\dd\mu
    = \int_{\X\times\Y}[\phi(y)-\psi(x)]\dd\pi
    \le \int_{\X\times\Y} c\dd\pi.
  \]
  \textbf{证明。} 等式由边际条件给出（$\int\phi\dd\pi=\int\phi\dd\nu$，$\int\psi\dd\pi=\int\psi\dd\mu$）；
  不等式由约束 $\phi-\psi\le c$ 及 $\pi\ge0$ 直接得到。

  \TLtakeaway{弱对偶：中介能争取的最大利润不超过自运成本。若某对 $(\pi,\psi,\phi)$ 使等号成立，则同时找到了最优计划和最优价格。}
\end{frame}

% ── 弱对偶推论 ───────────────────────────────────────────────────────────────
\begin{frame}{价格视角：弱对偶推论（二）}
  \deflab{推论（弱对偶不等式）。}
  \[
    \sup_{\phi-\psi\le c}\left[\int_{\Y}\phi\dd\nu-\int_{\X}\psi\dd\mu\right]
    \;\le\;
    \inf_{\pi\in\coup{\mu}{\nu}}\int_{\X\times\Y} c\dd\pi.
  \]
  \begin{itemize}
    \item 对偶问题的\emph{任何}可行解给出原始问题的\emph{下界}——这是"弱对偶"。
    \item \textbf{强对偶}（等号成立）是定理 5.10(i) 的深刻内容，需要额外论证。
    \item 实用技巧：若能展示某对 $(\pi,\psi,\phi)$ 使两边都等于同一个值，则两者同时最优。
  \end{itemize}

  \TLtakeaway{弱对偶：$\text{对偶上界} \le \text{原始下界}$。强对偶等号的成立是 Kantorovich 理论的核心定理（定理 5.10）。}
\end{frame}

% ── c-变换 ──────────────────────────────────────────────────────────────────
\begin{frame}{c-变换的定义}
  \textbf{如何找到最优价格？} 给定 $\psi$，最佳的 $\phi(y)$ 应取 $\psi(x)+c(x,y)$ 的下确界：

  \deflab{定义（c-变换）。} 函数 $\psi:\X\to\R\cup\{+\infty\}$ 的\TLterm{c-变换}为
  \[
    \ctr{\psi}(y) \;:=\; \inf_{x\in\X}\bigl[\psi(x)+c(x,y)\bigr],\qquad y\in\Y.
  \]
  （记号 $\ctr{\psi}=\psi^c$。）对称地，$\phi:\Y\to\R\cup\{-\infty\}$ 的 c-变换为
  $\phi^c(x):=\sup_{y\in\Y}[\phi(y)-c(x,y)]$。

  \begin{itemize}
    \item 直觉：$\ctr{\psi}(y)$ 是"在 $y$ 处，考虑所有可能的 $x$ 来源后，诚实竞争价格的上界"。
    \item 由于 $\ctr{\psi}(y)-\psi(x)\le c(x,y)$，$(\psi,\ctr{\psi})$ 自动是竞争价格对。
  \end{itemize}

  \TLtakeaway{c-变换把"买方价格 $\psi$"转化为"最优卖方价格 $\ctr{\psi}$"——是 Legendre 变换的推广，也是对偶理论的核心操作。}
\end{frame}

% ── 紧价格对 ─────────────────────────────────────────────────────────────────
\begin{frame}{紧价格对与对偶问题的约简}
  \deflab{紧价格对（Villani 式 5.5）。} 价格对 $(\psi,\phi)$ 是\TLterm{紧的（tight）}，若
  \[
    \phi(y)=\inf_{x}[\psi(x)+c(x,y)]=\ctr{\psi}(y)
    \quad\text{且}\quad
    \psi(x)=\sup_{y}[\phi(y)-c(x,y)].
  \]
  即：不能再提高卖价，也不能降低买价——价格已达"诚实竞争"的极致。

  \deflab{约简引理。} 从任意竞争价格对 $(\psi,\phi)$ 出发，令 $\phi_1=\ctr{\psi}$，$\psi_1=(\phi_1)^c$，则 $(\psi_1,\phi_1)$ 是紧价格对，且
  $\int\phi_1\dd\nu-\int\psi_1\dd\mu \ge \int\phi\dd\nu-\int\psi\dd\mu$（利润不减）。

  因此，对偶问题可以\emph{不失一般性}地只在紧价格对上取极值。

  \TLtakeaway{一步迭代即可将任意竞争价格对变成紧的：$(\psi,\phi)\mapsto(\psi_1,\ctr{\psi_1})$。对偶问题化为：寻找最优 $\psi$（c-凸性即为其特征）。}
\end{frame}

% ── 定义 5.2：c-凸性 ─────────────────────────────────────────────────────────
\begin{frame}{定义 5.2：c-凸函数}
  \deflab{定义 5.2（c-凸函数，Villani Def 5.2）。}
  设 $c:\X\times\Y\to(-\infty,+\infty]$。函数 $\psi:\X\to\R\cup\{+\infty\}$
  称为\TLterm{c-凸的}，若它不恒为 $+\infty$，且存在 $\zeta:\Y\to\R\cup\{\pm\infty\}$ 使得
  \[
    \psi(x) \;=\; \sup_{y\in\Y}\bigl[\zeta(y)-c(x,y)\bigr] \qquad \forall\,x\in\X.
  \]
  此时 $\psi$ 的 c-变换为 $\ctr{\psi}(y)=\inf_{x}[\psi(x)+c(x,y)]$，
  c-次微分为
  \[
    \partial_c\psi \;:=\; \{(x,y)\in\X\times\Y \;:\; \ctr{\psi}(y)-\psi(x)=c(x,y)\}.
  \]
  $\psi$ 与 $\ctr{\psi}$ 互称为\TLterm{c-共轭函数}。

  \begin{itemize}\small
    \item \textbf{等价描述。} $\psi$ c-凸 $\Leftrightarrow$ $\psi(x)=\inf_{y}[\ctr{\psi}(y)+c(x,y)]$ 对所有 $x$ 成立（即 $\cctr{\psi}=\psi$，见命题 5.8）。
    \item \textbf{c-次微分的含义。} $(x,y)\in\partial_c\psi$ 等价于 $x\in\arg\min_{x'}[\psi(x')+c(x',y)]$——即 $y$ 只接受从\emph{总成本最低}的 $x$ 运来的货物。
  \end{itemize}

  \TLtakeaway{c-凸函数 = "由 $c(x,\cdot)$ 的平移从下方刻画"的函数（图 5.2 直觉）。$\partial_c\psi$ 是最优匹配关系的几何编码。}
\end{frame}

% ── c-凸性图示（图 5.2 重绘）────────────────────────────────────────────────
\begin{frame}{c-凸性的几何直觉（图 5.2 重绘）}
  \begin{columns}[T]
    \begin{column}{0.55\textwidth}
      \begin{tikzpicture}[scale=1.0,>={Stealth[length=4pt]}]
        % x 轴
        \draw[->] (-0.3,0)--(5.2,0) node[right,font=\scriptsize]{$x$};
        \draw[->] (0,-0.3)--(0,3.2) node[above,font=\scriptsize]{值};
        % c-凸函数 psi（折线近似）
        \draw[TLprimary,very thick]
          (0.2,2.4) .. controls (1.2,1.0) and (2.0,0.5) .. (2.6,0.6)
          .. controls (3.2,0.8) and (4.0,1.8) .. (4.8,2.6)
          node[right,font=\scriptsize,TLprimary]{$\psi$};
        % 三个"工具"曲线：zeta(y_i) - c(x, y_i)（细线，从下方贴近）
        \foreach \yi/\col/\pos in {0/TLaccent/0.8, 1/TLpos/2.6, 2/TLinkSoft/4.2}{
          \draw[\col,thin,dashed]
            (0.2,{2.2-0.6*\yi}) .. controls (2,{1.1-0.4*\yi}) and (3,{0.3+0.3*\yi}) .. (4.8,{1.9+0.2*\yi});
        }
        % 接触点（c-次微分元素）
        \fill[TLaccent] (0.9,1.6) circle (2.5pt) node[right,font=\scriptsize,TLaccent]{$y_0\in\partial_c\psi(x_0)$};
        \fill[TLpos] (2.6,0.6) circle (2.5pt) node[right,font=\scriptsize,TLpos]{$y_1\in\partial_c\psi(x_1)$};
        \fill[TLinkSoft] (4.1,1.9) circle (2.5pt) node[left,font=\scriptsize,TLinkSoft]{$y_2\in\partial_c\psi(x_2)$};
      \end{tikzpicture}
    \end{column}
    \begin{column}{0.42\textwidth}
      \textbf{直觉说明：}
      \begin{itemize}\small
        \item 虚线为 $\zeta(y_i)-c(\cdot,y_i)$（形状随 $y_i$ 变化的"工具"）。
        \item $\psi$ 是这些工具的\emph{逐点上确界}。
        \item 每个接触点 $(x_i,y_i)$ 即一个 c-次微分元素：工具恰好从下方"吻到"$\psi$。
        \item Villani："c-凸函数是可以被形状为 $-c(x,\cdot)$ 的工具从下方完整抚摸的函数。"
      \end{itemize}
      \vspace{4pt}
      {\scriptsize 据 Villani 图 5.2 重绘。}
    \end{column}
  \end{columns}
  \TLtakeaway{c-凸函数的图象被一族"成本工具"从下方生成：每个接触点给出一个最优匹配 $(x_i\leftrightarrow y_i)$，即 $\partial_c\psi$ 中的元素。}
\end{frame}

% ── 特殊情形 5.3 与 5.4 ──────────────────────────────────────────────────────
\begin{frame}{c-凸性的两个重要特殊情形（5.3 与 5.4）}
  \deflab{特殊情形 5.3（$c=-x\cdot y$，Villani PC 5.3）。}
  在 $\R^n\times\R^n$ 上取 $c(x,y)=-x\cdot y$，则：
  \begin{itemize}\small
    \item c-变换 $\ctr{\psi}(y)=\inf_{x}[\psi(x)-x\cdot y]=-\sup_{x}[x\cdot y-\psi(x)]$，
      即 $\ctr{\psi}=-\psi^*$（$\psi$ 的\emph{Legendre 变换}的相反数）。
    \item c-凸 $\Leftrightarrow$ $\psi$ 是\emph{凸函数}且下半连续（lsc）。
    \item 这将 Kantorovich 对偶化为经典 Legendre 对偶——后者是前者的特例。
  \end{itemize}

  \medskip
  \deflab{特殊情形 5.4（$c=d$，距离，Villani PC 5.4）。}
  在度量空间 $(\X,d)$ 上取 $c=d$，则：
  \begin{itemize}\small
    \item c-凸 $\Leftrightarrow$ $\psi$ 是\emph{1-Lipschitz 函数}（$|\psi(x)-\psi(x')|\le d(x,x')$）。
    \item 1-Lipschitz 函数的 c-变换是其自身：$\ctr{\psi}=\psi$。
    \item 由此可以得到 Kantorovich–Rubinstein 公式（定理 5.10 的推论）：
      $W_1(\mu,\nu)=\sup_{\psi\text{ 1-Lip}}\int\psi\dd\mu-\int\psi\dd\nu$。
  \end{itemize}

  \TLtakeaway{c-凸性统一了两种经典对偶：$c=-x\cdot y$ 给出 Legendre 变换；$c=d$ 给出 Kantorovich–Rubinstein 公式。}
\end{frame}

% ── 命题 5.8：c-凸性的等价刻画 ────────────────────────────────────────────────
\begin{frame}{命题 5.8：c-凸性的等价刻画与双重 c-变换}
  \deflab{命题 5.8（Villani Prop 5.8）。}
  对任意函数 $\psi:\X\to\R\cup\{+\infty\}$，定义其\TLterm{二次 c-变换}为 $\cctr{\psi}=(\ctr{\psi})^c$，即
  \[
    \cctr{\psi}(x) = \sup_{y\in\Y}\Bigl[\inf_{x'\in\X}\bigl(\psi(x')+c(x',y)\bigr)-c(x,y)\Bigr].
  \]
  则 $\psi$ 是 c-凸的，当且仅当 $\cctr{\psi}=\psi$（即经过两次 c-变换后函数不变）。

  \medskip
  \textbf{证明。} 一般地，对任意 $\phi:\Y\to\R\cup\{-\infty\}$，有 $(\phi^c)^{cc}=\phi^c$。
  这是因为：取 $x'=x$ 可得 $\le$；取 $y'=y$ 可得 $\ge$，两个方向均由定义直接给出。
  由此：若 $\psi=\zeta^c$（c-凸），则 $\cctr{\psi}=\zeta^{ccc}=\zeta^c=\psi$。
  反之若 $\cctr{\psi}=\psi$，则 $\psi=(\ctr{\psi})^c$ 是 $\ctr{\psi}$ 的 c-变换，故 c-凸。

  \medskip
  \deflab{注记 5.9。} 命题 5.8 是 Legendre 对偶的推广：取 $c=-x\cdot y$ 时恢复
  $\psi^{**}=\psi$（凸的充要条件）。

  \TLtakeaway{c-凸 $\Leftrightarrow$ $\cctr{\psi}=\psi$：两次 c-变换后不变。这推广了 Legendre 对偶的"双共轭等于自身"——整个理论的基本代数事实。}
\end{frame}

% ── 定理 5.10（一）：对偶等式陈述 ────────────────────────────────────────────
\begin{frame}{定理 5.10：Kantorovich 对偶（一）——对偶等式}
  \deflab{定理 5.10（Kantorovich 对偶，Villani Thm 5.10）。}
  设 $\X,\Y$ 为 Polish 空间，$\mu\in\Pp(\X)$，$\nu\in\Pp(\Y)$，$c:\X\times\Y\to\R\cup\{+\infty\}$ 下半连续，且
  $c(x,y)\ge a(x)+b(y)$ 对某上半连续的 $a\in L^1(\mu)$，$b\in L^1(\nu)$。则

  \deflab{(i) 存在对偶等式：}
  \[
    \min_{\pi\in\coup{\mu}{\nu}}\int c\dd\pi
    \;=\;\sup_{\substack{\psi\in L^1(\mu),\,\phi\in L^1(\nu)\\\phi-\psi\le c}}\!\!\!\!
      \Bigl[\int_{\Y}\phi\dd\nu-\int_{\X}\psi\dd\mu\Bigr]
    \;=\;\sup_{\psi\text{ c-凸}}\Bigl[\int_{\Y}\ctr{\psi}\dd\nu-\int_{\X}\psi\dd\mu\Bigr].
  \]
  且对偶问题的上确界可只在 $\psi$ 为 c-凸、$\phi=\ctr{\psi}$ 为 c-凹的对上取得。

  \begin{itemize}\small
    \item 注意左端是\emph{最小值}（最优耦合存在，由定理 4.1）。
    \item 右端对偶问题的\emph{上确界}在温和条件下也是最大值（见定理 5.10(iii)）。
  \end{itemize}

  \TLtakeaway{最优传输成本 $=$ Kantorovich 对偶最优值：成本的最小化与价格利润的最大化在最优处\emph{完全吻合}。}
\end{frame}

% ── 定理 5.10（二）：几何刻画 ────────────────────────────────────────────────
\begin{frame}{定理 5.10：Kantorovich 对偶（二）——几何刻画}
  \deflab{(ii) 最优性的等价条件（Villani Thm 5.10(ii)）。}
  设 $c$ 实值，$C(\mu,\nu)<+\infty$。则存在 c-循环单调 Borel 集 $\Gamma\subset\X\times\Y$，
  使得对任意 $\pi\in\coup{\mu}{\nu}$，以下五个命题等价：
  \begin{enumerate}\small
    \item[(a)] $\pi$ 最优（使 $\int c\dd\pi$ 达到最小值）；
    \item[(b)] $\pi$ 是 c-循环单调的（集中在 c-循环单调集上）；
    \item[(c)] 存在 c-凸函数 $\psi$，使得 $\pi$-a.s.\ $\ctr{\psi}(y)-\psi(x)=c(x,y)$；
    \item[(d)] 存在 $\psi:\X\to\R\cup\{+\infty\}$，$\phi:\Y\to\R\cup\{-\infty\}$，满足 $\phi-\psi\le c$ 对全集成立，且等号 $\pi$-a.s.\ 成立；
    \item[(e)] $\pi$ 集中在 $\Gamma$ 上（同一个 $\Gamma$ 对所有最优 $\pi$ 共享）。
  \end{enumerate}

  \TLtakeaway{等价链：$\text{(a)最优} \Leftrightarrow \text{(b)循环单调} \Leftrightarrow \text{(c)集中在}\;\partial_c\psi \Leftrightarrow \text{(d)竞争价格等号}\;\Leftrightarrow\;\text{(e)集中在}\;\Gamma$。这五个刻画是本章最核心的内容。}
\end{frame}

% ── 等价链的实用意义 ─────────────────────────────────────────────────────────
\begin{frame}{定理 5.10 的实用推论（注记 5.13）}
  \textbf{如何用定理 5.10 判断最优性？}

  \begin{itemize}
    \item \textbf{验证 $\pi$ 最优。} 构造竞争价格对 $(\psi,\phi)$（$\phi-\psi\le c$），
      在 $\pi$ 的支撑上验证等号 $\phi(y)-\psi(x)=c(x,y)$。
    \item \textbf{验证 $(\psi,\phi)$ 最优。} 构造支撑上 $\phi-\psi=c$ 的可行耦合 $\pi$。
    \item \textbf{几何。} 最优 $(\psi,\phi)$ 固定后：质量从 $x$ 到 $y$ 当且仅当
      $x\in\arg\min_{x'}[\psi(x')+c(x',y)]$（咖啡馆只接受总成本最低面包店的货）。
  \end{itemize}

  \TLtakeaway{定理 5.10 的实用法则：构造 $(\pi,\psi,\phi)$ 并在支撑上验证 $\phi-\psi=c$，即可同时证明 $\pi$ 和 $(\psi,\phi)$ 均最优。}
\end{frame}

% ── 定理 5.10 等价条件的注意点 ───────────────────────────────────────────────
\begin{frame}{定理 5.10 等价条件辨析}
  % Manual warn-style block (avoids 0.4pt draw-border overflow from \TLwarnblock)
  \begin{tikzpicture}
    \node[draw=TLneg, fill=TLneg!8, rounded corners=3pt,
          inner sep=8pt, text width=\linewidth-17pt, align=justify] {%
      \small\textbf{\color{TLneg}易混淆：私有 c-单调集 vs 共享集 $\Gamma$}\\[0.3em]
      等价条件 (b)：$\pi$ 集中在\emph{某个} c-循环单调集上——该集合随 $\pi$ 变化（私有）。\\
      等价条件 (e)：所有最优 $\pi$ 共享同一集合 $\Gamma=\partial_c\psi$（$\psi$ = 固定的 Kantorovich 势）。
    };
  \end{tikzpicture}

  \begin{itemize}
    \item 若 $c$ 连续，最小共享集 $\Gamma_{\min}$（所有最优支撑的并的闭包）和最大共享集
      $\Gamma_{\max}$（所有最优 $\psi$ 的 c-次微分的交）均有意义。
    \item 注记 5.11：$c$ 不连续时，"集中在 $\Gamma$" 不等同于 "$\mathrm{supp}\,\pi\subset\Gamma$"。
    \item 注记 5.12：$\Gamma$ 一般不唯一；连续成本时可唯一确定 $\Gamma_{\min}$ 和 $\Gamma_{\max}$。
  \end{itemize}

  \TLtakeaway{公共集 $\Gamma$ 编码所有最优计划共享的几何信息；$\Gamma$ 不唯一，任意选取其一均可用于最优性刻画。}
\end{frame}

% ── 证明思路（一）：构造 Kantorovich 势 ──────────────────────────────────────
\begin{frame}{定理 5.10 证明思路（一）：从 c-循环单调集构造势函数}
  \textbf{核心构造（Rüschendorf/Rockafellar）。}
  给定一个 c-循环单调集 $\Gamma$，固定基点 $(x_0,y_0)\in\Gamma$，令 $\psi(x_0)=0$，
  对任意 $x\in\X$ 定义（Villani 式 5.17）：
  \[
    \psi(x) := \sup_{m,\,(x_1,y_1),\ldots,(x_m,y_m)\in\Gamma}
    \bigl[c(x_0,y_0)-c(x_1,y_0)+c(x_1,y_1)-c(x_2,y_1)
    +\cdots+c(x_m,y_m)-c(x,y_m)\bigr].
  \]

  \textbf{关键验证。}
  \begin{itemize}\small
    \item \textbf{$\psi(x_0)=0$。} 由 c-循环单调性，上式对 $(x_1,y_1)=(x_0,y_0)$ 的所有链均 $\le 0$，而取 $m=1,(x_1,y_1)=(x_0,y_0)$ 时为 $0$，故 $\psi(x_0)=0$。
    \item \textbf{$\psi$ 是 c-凸函数。} 通过引入辅助函数 $\zeta(y)=\sup_{\text{链}}\{\cdots+c(x_m,y)\}$，可将 $\psi$ 改写为 $\psi(x)=\sup_{y}[\zeta(y)-c(x,y)]$，这正是 c-凸性定义。
    \item \textbf{$\Gamma\subset\partial_c\psi$。} 对 $(x,y)\in\Gamma$：利用 $\psi$ 的定义，可验证 $\ctr{\psi}(y)=\psi(x)+c(x,y)$。
  \end{itemize}

  \TLtakeaway{构造关键：c-循环单调性保证了 $\psi(x_0)=0$（有界不发散）；由此得到 c-凸函数 $\psi$，其 c-次微分恰好包含 $\Gamma$。}
\end{frame}

% ── 证明思路（二）：对偶等式的建立 ──────────────────────────────────────────
\begin{frame}{定理 5.10 证明思路（二）：对偶等式的完成}
  \textbf{承接前页的构造（$\psi$ c-凸，$\phi=\ctr{\psi}$，$\pi$ 集中在 $\partial_c\psi$）。}

  \textbf{关键计算。} 由边际条件与 $\phi(y)-\psi(x)=c(x,y)$（$\pi$-a.s.）：
  \[
    \int_{\X\times\Y} c(x,y)\dd\pi
    \;=\; \int_{\X\times\Y}[\phi(y)-\psi(x)]\dd\pi
    \;=\; \int_{\Y}\phi\dd\nu - \int_{\X}\psi\dd\mu.
  \]
  因此：$\int c\dd\pi = \int\ctr{\psi}\dd\nu - \int\psi\dd\mu \le \sup_{\psi'\text{ c-凸}}\Bigl[\int\ctr{\psi'}\dd\nu-\int\psi'\dd\mu\Bigr]$，
  但由弱对偶，后者 $\le\int c\dd\pi$。两个方向均成立，故\emph{等号}处处成立。

  \begin{itemize}\small
    \item \textbf{连续成本的情形：} 上述论证通过"从离散逼近 $\mu_n,\nu_n$ 出发，取弱极限"严格化。
    \item \textbf{下半连续成本：} 进一步通过 $c=\lim c_k$（单调逼近，$c_k$ 有界连续）完成。
    \item \textbf{可测性：} $\psi,\ctr{\psi}$ 的可测性需要额外论证（利用 $c$ 的下半连续性和 Luzin 定理），参见 Villani 推论 5.22。本讲义在此处省略该步骤。
  \end{itemize}

  \TLtakeaway{对偶等式的证明方向：先用 c-循环单调集构造 $\psi$，再通过边际条件与等号条件把积分恒等式"锁住"。}
\end{frame}

% ── Rockafellar 定理（经典情形）────────────────────────────────────────────
\begin{frame}{Rockafellar 定理：$c=-x\cdot y$ 的经典情形}
  取 $c(x,y)=-x\cdot y$ 在 $\R^n\times\R^n$ 上，则定理 5.10 化为经典的 Rockafellar 定理：

  \deflab{Rockafellar 定理（经典凸分析）。}
  $\Gamma\subset\R^n\times\R^n$ 是循环单调的（即"凸分析意义下的循环单调"）当且仅当存在凸 lsc 函数 $f:\R^n\to\R\cup\{+\infty\}$，使得 $\Gamma\subset\partial f$（$f$ 的次微分的图象）。

  \textbf{与定理 5.10 的联系：}
  \begin{itemize}\small
    \item "c-循环单调"（$c=-x\cdot y$）$=$ Rockafellar 意义下的"循环单调"。
    \item "c-凸函数"$=$ 凸 lsc 函数（由特殊情形 5.3）。
    \item "c-次微分 $\partial_c\psi$"$=$ 凸分析次微分 $\partial f$（即 $\{(x,y):f(x)+f^*(y)=x\cdot y\}$）。
    \item 定理 5.10 是 Rockafellar 定理的\emph{严格推广}：将凸分析中"凸函数 $\leftrightarrow$ 次微分"替换为"c-凸函数 $\leftrightarrow$ c-次微分"。
  \end{itemize}

  \TLtakeaway{Rockafellar 定理 = 定理 5.10 在 $c=-x\cdot y$ 时的特殊情形：循环单调集 $\Leftrightarrow$ 凸函数的次微分。最优传输把这一经典结果推广到了任意成本函数。}
\end{frame}

% ── 定理 5.10(iii)：最优解存在的条件 ────────────────────────────────────────
\begin{frame}{定理 5.10(iii)：对偶最优解的存在性条件}
  \deflab{(iii) 对偶最优解（Villani Thm 5.10(iii)）。}
  在定理 5.10(i) 的条件下，若还有逐点上界
  \[
    c(x,y) \;\le\; c_{\X}(x)+c_{\Y}(y), \qquad (c_{\X},c_{\Y})\in L^1(\mu)\times L^1(\nu),
  \]
  则对偶问题的\emph{上确界是最大值}（即存在最优价格对），且
  \[
    \min_{\pi\in\coup{\mu}{\nu}}\int c\dd\pi
    \;=\;\max_{\psi\text{ c-凸}}\Bigl[\int_{\Y}\ctr{\psi}\dd\nu-\int_{\X}\psi\dd\mu\Bigr],
  \]
  最优 $\psi$（即 Kantorovich 势）满足 $\psi\in L^1(\mu)$，$\ctr{\psi}\in L^1(\nu)$。

  \textbf{证明思路。} 前页构造的 $\psi$ 在此条件下有界（可积）：利用上界
  $c_{\X}(x)+c_{\Y}(y)$ 建立 $\psi(x)\ge\phi(y_0)-c_{\Y}(y_0)-c_{\X}(x)$（下界为 $\mu$-可积函数），
  $\phi(y)\le\psi(x_0)+c_{\X}(x_0)+c_{\Y}(y)$（上界为 $\nu$-可积函数），因此两者均可积。

  \TLtakeaway{若成本有逐点可积上界，则对偶问题的"sup"变为"max"：Kantorovich 势 $\psi$ 不仅存在，而且可积、最优——三件事同时成立。}
\end{frame}

% ── 稳定性：定理 5.20 ─────────────────────────────────────────────────────────
\begin{frame}{稳定性：最优计划在弱收敛下保持（定理 5.20）}
  \deflab{定理 5.20（最优传输的稳定性，Villani Thm 5.20）。}
  设 $c\in C(\X\times\Y)$ 连续，$\inf c>-\infty$，$(c_k)$ 一致收敛到 $c$，
  $\mu_k\wkto\mu$，$\nu_k\wkto\nu$。对每个 $k$，设 $\pi_k$ 是 $(\mu_k,\nu_k)$ 之间关于 $c_k$ 的最优耦合，
  且 $\int c_k\dd\pi_k<+\infty$。则：
  \begin{enumerate}\small
    \item $(\pi_k)$ 存在弱收敛子列，极限 $\pi\in\coup{\mu}{\nu}$ 是 c-循环单调的。
    \item 若还有 $\liminf_k\int c_k\dd\pi_k<+\infty$，则 $C(\mu,\nu)<+\infty$，$\pi$ 是 $(\mu,\nu)$ 之间的最优耦合。
  \end{enumerate}

  \textbf{证明要点：}
  \begin{itemize}\small
    \item \emph{存在子列：} $\mu_k,\nu_k$ 弱收敛 $\Rightarrow$ 胶着 $\Rightarrow$ Lemma 4.4 给出 $(\pi_k)$ 胶着 $\Rightarrow$ Prokhorov 抽取子列。
    \item \emph{极限 c-循环单调：} 每个 $\pi_k^{\otimes N}$ 集中在 $c_k$-循环单调集上，该集在 $c_k\to c$ 后取极限为 c-循环单调闭集，弱极限继承此性质。
    \item \emph{最优性：} 由 $c\dd\pi\le\liminf c_k\dd\pi_k$ 及定理 5.10(ii) 给出。
  \end{itemize}

  \TLtakeaway{最优传输计划对边际测度的弱收敛稳定：极限仍是最优的。这保证了最优传输是一个"连续"的运算。}
\end{frame}

% ── 推论 5.23：传输映射的稳定性 ──────────────────────────────────────────────
\begin{frame}{推论 5.23：传输映射的稳定性}
  \deflab{推论 5.23（传输映射稳定，Villani Cor 5.23）。}
  设 $T_k:\X\to\Y$ 是关于 $c_k$ 在 $(\mu,\nu_k)$ 之间的最优传输映射（Monge 映射），
  $\nu_k\wkto\nu$。若 $T:\X\to\Y$ 是关于 $c$ 在 $(\mu,\nu)$ 之间的\emph{唯一}最优传输映射，则
  \[
    T_k \;\To\; T \quad \text{依概率}\ (\mu\text{-概率}):
    \quad
    \forall\eps>0,\quad \mu\bigl[\{x:\,d(T_k(x),T(x))\ge\eps\}\bigr]\;\To\;0.
  \]

  \textbf{证明思路。} 由定理 5.20，$\pi_k=(Id,T_k)_{\#}\mu\wkto(Id,T)_{\#}\mu=\pi$。
  利用 Lusin 定理，$T$ 在某紧子集 $K\subset\X$ 上连续（$\mu[K]\ge1-\delta$）；
  再由 $\pi_k\to\pi$ 弱收敛与 $\{d(T_k(x),T(x))\ge\eps\}$ 对应集合的闭性，得到 $\mu$-概率趋向 0。

  \medskip
  \deflab{注记 5.24。} 唯一性假设\emph{不可缺}：若最优映射不唯一，$T_k$ 可能振荡。
  若 $\mu_k\to\mu$ 而不是 $\mu$ 固定，稳定性一般失效（见注记 5.25 中的反例）。

  \TLtakeaway{若极限最优映射\emph{唯一}，则近似问题的最优映射按概率收敛到它——传输的定性性质在弱极限下保持。}
\end{frame}

% ── Brenier 定理（定理 5.30 准则）────────────────────────────────────────────
\begin{frame}{定理 5.30：Monge 问题的可解性准则}
  \deflab{定理 5.30（Monge 问题的可解性准则，Villani Thm 5.30）。}
  在定理 5.10 的条件下，若还满足：
  \begin{enumerate}\small
    \item $C(\mu,\nu)<+\infty$；
    \item 对任意 c-凸函数 $\psi$，使得 $\partial_c\psi(x)$ 包含\emph{多于一个}元素的 $x$ 的集合为 $\mu$-零测。
  \end{enumerate}
  则存在\emph{唯一的}（律意义下）最优耦合，且它是确定性的 Monge 映射：存在可测 $T:\X\to\Y$ 使得 $\pi=(Id,T)_{\#}\mu$ 是唯一的最优耦合。

  \textbf{证明思路。} 存在 c-凸 $\psi$ 使最优 $\pi$ 集中在 $\partial_c\psi$。由条件 (ii)，$\partial_c\psi(x)$ 几乎处处单点，
  可定义 $T(x)$ 为该唯一的 $y$；任何最优 $\pi$ 均集中在 $T$ 的图上。

  \TLtakeaway{定理 5.30 是"c-次微分几乎处处单点 $\Rightarrow$ 最优映射存在且唯一"的精确表述，为 Brenier 定理奠定基础。}
\end{frame}

% ── Brenier 定理的标准应用 ───────────────────────────────────────────────────
\begin{frame}{Brenier 定理：二次成本下的结论}
  \textbf{标准应用。} 取 $c(x,y)=|x-y|^2/2$ 在 $\R^n$ 上，$\mu\in\Pac(\R^n)$（绝对连续）。

  \deflab{Brenier 定理。}
  \begin{itemize}
    \item c-凸函数 = 凸 lsc 函数（由特殊情形 5.3，$c=|x-y|^2/2$ 等价于 $-x\cdot y$）。
    \item 定理 5.30 条件 (ii) 自动满足：$\mu$ 绝对连续时，凸函数的次梯度集 $\mu$-a.e.\ 单点。
    \item 结论：唯一最优 Monge 映射为
      \[
        T = \nabla\varphi, \qquad \varphi\;\text{凸函数},
      \]
      即最优映射是某个凸函数的梯度。
  \end{itemize}

  \TLtakeaway{Brenier 定理：$c=|x-y|^2/2$，$\mu$ 绝对连续 $\Rightarrow$ 最优映射唯一，形如 $T=\nabla\varphi$（凸函数的梯度）。这是最优传输最著名的结果，是 Monge-Ampère 方程的出发点。}
\end{frame}

% ── Brenier：1D 例子 ─────────────────────────────────────────────────────────
\begin{frame}{Brenier 定理：1D 单调重排（具体例子）}
  \textbf{离散例子。} 取 $\mu=\tfrac{1}{3}(\delta_{0}+\delta_{1}+\delta_{2})$，$\nu=\tfrac{1}{3}(\delta_{1}+\delta_{3}+\delta_{5})$，$c=|x-y|^2/2$。

  最优 Monge 映射：$T(0)=1,\;T(1)=3,\;T(2)=5$（单调）。
  总成本 $=\tfrac{1}{3}\cdot\tfrac{1}{2}(1+4+9)=\tfrac{7}{3}$。
  非单调映射如 $T(0)=5,T(1)=3,T(2)=1$ 的成本 $=\tfrac{1}{3}\cdot\tfrac{1}{2}(25+4+1)=5>\tfrac{7}{3}$，更差。

  \textbf{1D 一般情形（$\mu$ 无原子）。} 设 $F_\mu,F_\nu$ 为分布函数，$F_\nu^{-1}$ 为广义逆，则最优映射为
  \[
    T(x) = F_{\nu}^{-1}(F_{\mu}(x)).
  \]
  直觉：把 $\mu$ 的第 $p$ 分位数对应到 $\nu$ 的第 $p$ 分位数——始终单调。

  \TLtakeaway{1D 单调重排 $T=F_\nu^{-1}\circ F_\mu$ 是 Brenier 定理最简单的具体化：最优映射 = 分布函数的广义逆复合，即"按分位数对齐"。}
\end{frame}

% ── Brenier：高斯情形 ─────────────────────────────────────────────────────────
\begin{frame}{Brenier 定理：高斯情形与矩阵几何平均}
  \textbf{高斯情形（$\R^n$）。} 设 $\mu=\mathcal{N}(0,A)$，$\nu=\mathcal{N}(0,B)$，$A,B$ 正定，$c=|x-y|^2/2$。

  最优映射为线性映射 $T(x)=Sx$，其中
  \[
    S = A^{-1/2}\bigl(A^{1/2}BA^{1/2}\bigr)^{1/2}A^{-1/2}.
  \]
  \begin{itemize}\small
    \item $S$ 满足 $S A S = B$（映射把协方差 $A$ 变为 $B$）；
    \item $SA$ 对称正定，故 $S=\nabla\varphi$ 对某凸函数 $\varphi$（确认 Brenier 定理结论）；
    \item 高斯情形的 $W_2$ 距离：$W_2^2(\mu,\nu)=\mathrm{tr}(A+B-2(A^{1/2}BA^{1/2})^{1/2})$；
    \item $S$ 即两协方差矩阵的\emph{几何平均}——高斯情形的 Riemannian 测地线正中点。
  \end{itemize}

  \TLtakeaway{高斯情形：最优映射 $T=Sx$（线性），$S$ 由矩阵几何平均给出。$W_2$ 距离有闭合公式。这是无穷维 Wasserstein 几何在高斯族上的精确计算。}
\end{frame}

% ── 习题 2 · Exercises ────────────────────────────────────────────────────────
\begin{frame}{习题 2 · Exercises（提示见本页，解答见附录）}
  \begin{itemize}
    \item \textbf{习题 2.1}~(\diffI)\ 设 $\X=\Y=\{1,2,3\}$，$c(i,j)=|i-j|$，$\mu=\nu=\frac13\sum_i\delta_i$。
      写出所有耦合中最优的（Monge）耦合 $\pi^\star$，并验证其支撑是 c-循环单调的。
      \begin{itemize}\footnotesize
        \item \hint{最优 Monge 映射为恒等映射；验证对任意 2-元循环 $(i,i),(j,j)$，改道后不减少成本。}
      \end{itemize}
      % SOLUTION: 最优耦合为恒等映射 pi^star = (1/3)(delta_{(1,1)}+delta_{(2,2)}+delta_{(3,3)})，总成本为 0。对任意 (i,i),(j,j) in supp(pi^star)，c(i,i)+c(j,j)=0 <= c(i,j)+c(j,i)=2|i-j|，c-循环单调性成立。

    \item \textbf{习题 2.2}~(\diffII)\ 设 $\psi(x)=\frac12 x^2$ 在 $\R$ 上，$c(x,y)=\frac12|x-y|^2$。
      计算 $\ctr{\psi}(y)$，并确定 $\partial_c\psi$。
      \begin{itemize}\footnotesize
        \item \hint{$\ctr{\psi}(y)=\inf_x[\frac12 x^2+\frac12|x-y|^2]$，对 $x$ 求导令其为零；$\partial_c\psi=\{(x,y):\ctr{\psi}(y)-\psi(x)=c(x,y)\}$。}
      \end{itemize}
      % SOLUTION: ctr{psi}(y) = inf_x [x^2/2 + (x-y)^2/2] = inf_x [x^2 - xy + y^2/2]。导数 = 2x - y = 0，即 x = y/2，代入得 ctr{psi}(y) = y^2/4 - y^2/2 + y^2/2 = y^2/4。验证：ctr{psi}(y) - psi(x) = y^2/4 - x^2/2 = c(x,y) = (x-y)^2/2 iff y^2/4 - x^2/2 = x^2/2 - xy + y^2/2 iff 0 = x^2 - xy + y^2/4 = (x-y/2)^2，即 x=y/2。故 partial_c psi = {(x,y): x = y/2}，即最优映射 T(x)=2x 把 mu 推前到 nu 时对应 y=2x，x=y/2。

    \item \textbf{习题 2.3}~(\diffIII)\ 设 $\mu\in\Pac(\R^n)$，$\nu\in\Pp(\R^n)$，$c=|x-y|^2/2$。
      利用定理 5.30 论证：若两个最优耦合均为 Monge 耦合（对应映射 $T_1,T_2$），则 $T_1=T_2$（$\mu$-a.s.）。
      \begin{itemize}\footnotesize
        \item \hint{两个最优 Monge 耦合均集中在同一个 $\partial_c\psi$（同一个 Kantorovich 势）的图上；利用 $\mu$ 绝对连续时 $\partial_c\psi(x)$ 几乎处处单点。}
      \end{itemize}
      % SOLUTION: 由定理 5.10(ii)，存在 c-凸函数 psi 使得 pi_k = (Id,T_k)_# mu 均集中在 partial_c psi 上，即 T_k(x) in partial_c psi(x) mu-a.s. 对 k=1,2。由 mu 绝对连续以及 c=|x-y|^2/2 的凸性，partial_c psi(x) 对 mu-a.e. x 为单点（这正是定理 5.30 条件(ii)的内容）。故 T_1(x)=T_2(x) mu-a.s.，即最优 Monge 映射在 mu-a.e. 意义下唯一。
  \end{itemize}
\end{frame}

% ── §2 小结 ──────────────────────────────────────────────────────────────────
\begin{frame}{§2 小结：最优传输的三棱柱——循环单调、对偶、势函数}
  本节建立了以下等价链（定理 5.10 的核心）：
  \[
    \pi\text{ 最优}
    \;\Longleftrightarrow\;
    \pi\text{ c-循环单调}
    \;\Longleftrightarrow\;
    \pi\text{ 集中在 }\partial_c\psi
    \;\Longleftrightarrow\;
    \text{竞争价格等号成立}
  \]

  \begin{itemize}\small
    \item \textbf{循环单调性（几何）。} 无法通过重路由有限循环来节省成本。
    \item \textbf{Kantorovich 对偶（价格）。} 最优传输成本 $=$ 最优价格利润；
      最优价格由 c-凸 Kantorovich 势 $\psi$ 给出，$\phi=\ctr{\psi}$。
    \item \textbf{c-次微分（结构）。} 最优计划集中在 $\partial_c\psi$：
      质量从 $x$ 到 $y$ 当且仅当 $\psi(x)+c(x,y)=\ctr{\psi}(y)$。
    \item \textbf{稳定性（定理 5.20）。} 对边际测度弱收敛稳定。
    \item \textbf{Brenier 定理（应用）。} $c=|x-y|^2/2$，$\mu$ 绝对连续 $\Rightarrow$ 唯一最优映射 $T=\nabla\varphi$（凸函数的梯度）。
  \end{itemize}

  \TLtakeaway{§2 是本讲义的"制高点"：三条路通向最优性（几何、价格、势函数），在定理 5.10 处交汇。后续各节（§3--§5）将在此基础上展开度量、曲率和热方程的应用。}
\end{frame}
     % §2  Ch.5 循环单调性与 Kantorovich 对偶
% =============================================================================
%  §3  Ch.6 Wasserstein 距离
%  来源：Villani, Optimal Transport: Old and New, Chapter 6 (pp. 93–108)
% =============================================================================
\TLsection{3 \ Wasserstein 距离}

% ── 动机 ─────────────────────────────────────────────────────────────────────
\begin{frame}{动机：从最优传输成本到概率测度上的距离}
  第 4--5 章构造了最优传输成本
  \[
    C(\mu,\nu) = \inf_{\pi\in\coup{\mu}{\nu}} \int_{\X\times\X} c(x,y)\dd\pi(x,y).
  \]
  $C(\mu,\nu)$ 定量刻画了把分布 $\mu$ 运输到 $\nu$ 的代价。

  \textbf{自然问题：} $C(\mu,\nu)$ 是概率测度空间上的\emph{距离}吗？

  \begin{columns}[T]
    \begin{column}{0.5\textwidth}
      \textbf{一般情形下 $C$ 的问题：}
      \begin{itemize}
        \item 成本 $c$ 不一定是距离，$C$ 未必满足三角不等式。
        \item $c$ 不对称时 $C$ 也不对称。
        \item $C(\mu,\nu)$ 可以是 $+\infty$。
      \end{itemize}
    \end{column}
    \begin{column}{0.5\textwidth}
      \textbf{当 $c=d^{p}$ 时的修正：}
      \begin{itemize}
        \item 令 $c(x,y)=d(x,y)^{p}$，$p\ge1$。
        \item 取 $p$ 次根：$\Wp{p}(\mu,\nu)=C(\mu,\nu)^{1/p}$。
        \item 本章证明 $\Wp{p}$ 是\TLterm{Wasserstein 距离}。
      \end{itemize}
    \end{column}
  \end{columns}

  \TLtakeaway{以 $d^{p}$ 为成本、取 $p$ 次根后，最优传输成本升格为概率测度空间上的真正距离。}
\end{frame}

% ── 定义 6.1 ─────────────────────────────────────────────────────────────────
\begin{frame}{定义 6.1：Wasserstein 距离}
  \deflab{定义 6.1（Wasserstein 距离，Villani Def 6.1）。}

  设 $(\X,d)$ 为 Polish 度量空间，$p\in[1,\infty)$。
  对 $\X$ 上的两个概率测度 $\mu,\nu$，\TLterm{$p$ 阶 Wasserstein 距离}定义为
  \[
    \Wp{p}(\mu,\nu) \;:=\;
    \left(\inf_{\pi\in\coup{\mu}{\nu}} \int_{\X\times\X} d(x,y)^{p}\dd\pi(x,y)\right)^{1/p}.
  \]
  等价地，若 $(X,Y)$ 为联合随机变量，$\law(X)=\mu$，$\law(Y)=\nu$，则
  \[
    \Wp{p}(\mu,\nu)
    = \inf\Bigl\{\bigl(\E[d(X,Y)^{p}]\bigr)^{1/p}:\;
      \law(X)=\mu,\;\law(Y)=\nu\Bigr\}.
  \]

  \deflab{特殊情形 6.2（Kantorovich–Rubinstein 距离）。}
  $\Wp{1}$ 也常称为 \TLterm{Kantorovich–Rubinstein 距离}（简称 KR 距离）。

  \TLtakeaway{以 $d^{p}$ 为成本的最优传输值取 $p$ 次根，即得 $p$ 阶 Wasserstein 距离 $\Wp{p}$。}
\end{frame}

% ── 例 6.3 ───────────────────────────────────────────────────────────────────
\begin{frame}{例 6.3：Wasserstein 距离继承基础度量}
  \deflab{例 6.3（Villani Example 6.3）。}
  对任意 $x,y\in\X$ 及任意 $p\ge1$，
  \[
    \Wp{p}(\delta_{x},\delta_{y}) \;=\; d(x,y).
  \]

  \textbf{验证。} $\coup{\delta_{x}}{\delta_{y}}=\{\delta_{(x,y)}\}$（唯一耦合），因此
  \[
    \Wp{p}(\delta_{x},\delta_{y})^{p}
    = \int_{\X\times\X} d(x',y')^{p}\dd\delta_{(x,y)}(x',y')
    = d(x,y)^{p}.
  \]
  取 $p$ 次根即得。结果不依赖 $p$——这是 Dirac 测度的特殊性。

  \textbf{含义。} 映射 $x\mapsto\delta_{x}$ 是 $(\X,d)$ 到 $(\Ppp{p}(\X),\Wp{p})$ 的\emph{等距嵌入}。

  \TLtakeaway{$\Wp{p}(\delta_x,\delta_y)=d(x,y)$：Wasserstein 距离忠实保留基础度量，Dirac 的嵌入是等距的。}
\end{frame}

% ── 距离公理（一）对称性与非退化性 ──────────────────────────────────────────
\begin{frame}{$\Wp{p}$ 是距离（一）：对称性与非退化性}
  \textbf{目标：} 验证 $\Wp{p}$ 满足距离的三条公理。

  \textbf{公理 1（对称性）：} $\Wp{p}(\mu,\nu)=\Wp{p}(\nu,\mu)$。

  \textbf{证明。} 若 $\pi\in\coup{\mu}{\nu}$，令 $\tilde{\pi}$ 为 $\pi$ 在 $(x,y)\mapsto(y,x)$ 下的像，则 $\tilde{\pi}\in\coup{\nu}{\mu}$，且 $\int d^{p}\dd\tilde{\pi}=\int d^{p}\dd\pi$（$d$ 对称）。两边下确界相等。$\square$

  \textbf{公理 2（非退化性）：} $\Wp{p}(\mu,\nu)=0 \iff \mu=\nu$。

  \textbf{证明（$\Leftarrow$）。} 取对角耦合 $\pi=(\Id,\Id)_{\#}\mu$，则 $\int d(x,y)^{p}\dd\pi=0$，故 $\Wp{p}(\mu,\nu)=0$。

  \textbf{证明（$\Rightarrow$）。} 若 $\Wp{p}(\mu,\nu)=0$，存在耦合使 $\int d(x,y)^{p}\dd\pi=0$，故 $\pi$ 集中在对角线 $\{y=x\}$ 上，从而 $\nu=\pf{\Id}{\mu}=\mu$。$\square$

  \TLtakeaway{对称性由 $d$ 的对称性继承；非退化性等价于最优耦合集中在对角线上 iff 两测度相同。}
\end{frame}

% ── 距离公理（二）粘合引理 ───────────────────────────────────────────────────
\begin{frame}{$\Wp{p}$ 是距离（二）：粘合引理的陈述}
  \deflab{粘合引理（Gluing Lemma，Villani Chapter 1）。}

  设 $\mu_{1},\mu_{2},\mu_{3}\in\Pp(\X)$。记 $\pi^{12}$ 为 $(\mu_{1},\mu_{2})$ 的一个耦合，
  $\pi^{23}$ 为 $(\mu_{2},\mu_{3})$ 的一个耦合（二者共享中间边际 $\mu_{2}$）。

  则存在定义在同一概率空间上的三元随机变量 $(X_{1},X_{2},X_{3})$ 使得
  \[
    \law(X_{1},X_{2}) = \pi^{12}, \quad
    \law(X_{2},X_{3}) = \pi^{23}.
  \]
  特别地，$\law(X_{1},X_{3})$ 是 $(\mu_{1},\mu_{3})$ 的一个耦合。

  \textbf{构造。} 把 $\pi^{(23)}$ 沿 $\mu_2$ disintegrate 为条件测度 $\pi^{(23)}_{x_2}$（见 §0），定义
  \[
    \pi^{(123)}(\dd x_{1}\dd x_{2}\dd x_{3})
    = \pi^{(12)}(\dd x_{1}\dd x_{2})\cdot\pi^{(23)}_{x_{2}}(\dd x_{3})
  \]
  （即给定 $X_2$ 时让 $X_1,X_3$ 条件独立）。

  \TLtakeaway{粘合引理：沿公共边际 $\mu_2$ 把两耦合"串联"成三元分布，$(X_1,X_3)$ 是 $(\mu_1,\mu_3)$ 的耦合。}
\end{frame}

% ── 距离公理（三）三角不等式 ────────────────────────────────────────────────
\begin{frame}{$\Wp{p}$ 是距离（三）：三角不等式证明}
  \textbf{公理 3（三角不等式）：} $\Wp{p}(\mu_{1},\mu_{3})\le\Wp{p}(\mu_{1},\mu_{2})+\Wp{p}(\mu_{2},\mu_{3})$。

  \textbf{证明。}
  设 $(X_{1},X_{2})$ 是 $(\mu_{1},\mu_{2})$ 的最优耦合，$(Z_{2},Z_{3})$ 是 $(\mu_{2},\mu_{3})$ 的最优耦合。

  由粘合引理，存在 $(X_{1},X_{2},X_{3})$ 使 $\law(X_{1},X_{3})$ 是 $(\mu_{1},\mu_{3})$ 的耦合，故
  \[
    \Wp{p}(\mu_{1},\mu_{3})
    \;\le\; \bigl(\E[d(X_{1},X_{3})^{p}]\bigr)^{1/p}.
  \]
  由 $d(X_{1},X_{3})\le d(X_{1},X_{2})+d(X_{2},X_{3})$ 及 $L^{p}$ Minkowski 不等式：
  \[
    \bigl(\E[d(X_{1},X_{3})^{p}]\bigr)^{1/p}
    \;\le\; \bigl(\E[d(X_{1},X_{2})^{p}]\bigr)^{1/p}
           + \bigl(\E[d(X_{2},X_{3})^{p}]\bigr)^{1/p}.
  \]
  右边两项分别等于 $\Wp{p}(\mu_{1},\mu_{2})$ 和 $\Wp{p}(\mu_{2},\mu_{3})$。$\square$

  \TLtakeaway{三步：粘合最优耦合 $\to$ $d$ 的三角不等式 $\to$ $L^{p}$ Minkowski 不等式。}
\end{frame}

% ── 定义 6.4：Wasserstein 空间 ────────────────────────────────────────────────
\begin{frame}{定义 6.4：Wasserstein 空间 $\Ppp{p}(\X)$}
  \deflab{定义 6.4（Wasserstein 空间，Villani Def 6.4）。}

  设 $(\X,d)$ 为 Polish 空间，$p\in[1,\infty)$，取定基点 $x_{0}\in\X$。\TLterm{$p$ 阶 Wasserstein 空间}定义为
  \[
    \Ppp{p}(\X) \;:=\;
    \left\{\mu\in\Pp(\X):\; \int_{\X} d(x_{0},x)^{p}\dd\mu(x) < +\infty\right\}.
  \]
  该定义不依赖 $x_{0}$ 的选取。

  \textbf{$\Wp{p}$ 在 $\Ppp{p}(\X)$ 上有限的证明。} 对 $\mu,\nu\in\Ppp{p}(\X)$，由
  \[
    d(x,y)^{p} \;\le\; 2^{p-1}\bigl(d(x,x_{0})^{p}+d(x_{0},y)^{p}\bigr)
  \]
  得 $\int d(x,y)^{p}\dd\pi\le 2^{p-1}\bigl(\int d(x_{0},x)^{p}\dd\mu+\int d(x_{0},y)^{p}\dd\nu\bigr)<+\infty$，故 $\Wp{p}(\mu,\nu)<+\infty$。$\square$

  \TLtakeaway{$\Ppp{p}(\X)$ 是 $p$ 阶矩有限的概率测度全体，$\Wp{p}$ 在其上为真正有限的距离。}
\end{frame}

% ── 定理 6.18 ────────────────────────────────────────────────────────────────
\begin{frame}{$\Ppp{p}(\X)$ 本身也是 Polish 空间（定理 6.18）}
  \deflab{定理 6.18（Wasserstein 空间的拓扑，Villani Thm 6.18）。}

  设 $\X$ 为 Polish 空间，$p\in[1,\infty)$。则 $\bigl(\Ppp{p}(\X),\Wp{p}\bigr)$ 也是 Polish 空间（完备可分度量空间），且有限支撑测度在其中稠密。

  \textbf{证明思路（可分性）。} 取 $\X$ 的稠密子集 $D$，有理系数、支撑于 $D$ 中有限多点的测度 $\mathcal{P}_{0}$ 在 $\Ppp{p}(\X)$ 中稠密。

  \textbf{为什么这很重要？}
  \begin{itemize}
    \item \textbf{完备性：} Cauchy 序列有极限，§4 的测地线曲线在此收敛。
    \item \textbf{可分性：} Wasserstein 空间上可以做测度论，为 §4 的置换插值铺路。
    \item 若 $\X$ 紧，则 $\Ppp{p}(\X)$ 也紧；若 $\X$ 仅局部紧，则 $\Ppp{p}(\X)$ 不局部紧。
  \end{itemize}

  \TLtakeaway{Polish 空间上的 Wasserstein 空间仍是 Polish 空间：完备性和可分性被继承，为 §4 的几何结构奠基。}
\end{frame}

% ── 注记 6.5：KR 对偶 ───────────────────────────────────────────────────────
\begin{frame}{注记 6.5：$\Wp{1}$ 的 Kantorovich–Rubinstein 对偶}
  \deflab{注记 6.5（KR 对偶公式，Villani Rem 6.5）。}

  由定理 5.10(i) 与特殊情形 5.4，对任意 $\mu,\nu\in\Ppp{1}(\X)$：
  \[
    \Wp{1}(\mu,\nu)
    \;=\; \sup\left\{\int_{\X}\psi\dd\mu - \int_{\X}\psi\dd\nu
          :\; \|\psi\|_{\mathrm{Lip}}\le1\right\}.
  \]
  这是 \TLterm{KR 对偶公式}，全书最常用的计算工具之一。

  \textbf{应用举例（协方差不等式）。} 若 $f$ 是关于 $\mu$ 的概率密度，则
  \[
    \int f\dd\mu\int g\dd\mu - \int fg\dd\mu
    \;\le\; \|g\|_{\mathrm{Lip}}\cdot\Wp{1}(f\mu,\mu).
  \]

  \textbf{含义。} 上界将函数的协方差与 $\Wp{1}$ 距离联系起来，是 $\Wp{1}$ 在概率估计中的典型应用。

  \TLtakeaway{KR 对偶：$\Wp{1}$ 等于 1-Lipschitz 函数积分差的上确界，使得计算转化为寻找最优 Lipschitz 检验函数。}
\end{frame}

% ── 注记 6.6：阶的排序 ───────────────────────────────────────────────────────
\begin{frame}{注记 6.6：Wasserstein 距离的阶的排序}
  \deflab{注记 6.6（阶的排序，Villani Rem 6.6）。}

  由 Hölder 不等式：对 $1\le p\le q$，
  \[
    \Wp{p}(\mu,\nu) \;\le\; \Wp{q}(\mu,\nu).
  \]
  特别地，$\Wp{1}$ 是所有 Wasserstein 距离中\emph{最弱}的。

  \begin{center}\small
    \begin{tabular}{lll}
      \toprule
      & $\Wp{1}$ & $\Wtwo$ \\
      \midrule
      强度 & 弱（最易满足） & 强 \\
      估计难度 & 容易（KR 对偶） & 较难 \\
      几何信息 & 有限 & 丰富（黎曼风格） \\
      典型应用 & PDE、信息论 & OT 几何、置换插值 \\
      \bottomrule
    \end{tabular}
  \end{center}

  \TLtakeaway{$\Wp{1}\le\Wp{p}$：阶越高距离越强；$p=1$ 最灵活易估，$p=2$ 最能捕捉几何结构。}
\end{frame}

% ── 计算例（一）：一维分位公式——陈述与证明思路 ───────────────────────────────
\begin{frame}{计算例（一）：一维分位函数公式（陈述）}
  \deflab{一维闭合公式。}

  设 $\mu,\nu\in\Ppp{p}(\R)$，$F_{\mu}^{-1},F_{\nu}^{-1}$ 为分位函数（广义逆 CDF）。则
  \[
    \Wp{p}(\mu,\nu)^{p}
    \;=\; \int_{0}^{1} \bigl|F_{\mu}^{-1}(t) - F_{\nu}^{-1}(t)\bigr|^{p}\dd t.
  \]

  \textbf{证明思路。} $\R$ 上最优传输由\emph{单调重排}实现：$T=F_{\nu}^{-1}\circ F_{\mu}$ 是最优 Monge 映射，且
  \[
    \int_{\R} |x-T(x)|^{p}\dd\mu(x)
    = \int_{0}^{1}|F_{\mu}^{-1}(t)-F_{\nu}^{-1}(t)|^{p}\dd t.
  \]
  最优性：$T$ 单调故满足 $c$-循环单调性（第 5 章）；分位公式是其代价的等价表达。

  \TLtakeaway{一维情形：最优传输就是单调重排；Wasserstein 距离是分位函数差的 $L^{p}$ 范数。}
\end{frame}

% ── 计算例（一）：数值例 ────────────────────────────────────────────────────
\begin{frame}{计算例（一）续：均匀分布平移的 $\Wp{p}$}
  \textbf{数值例（均匀分布的平移）。}

  设 $\mu=\mathrm{Uniform}[0,1]$，$\nu=\mathrm{Uniform}[a,a+1]$（$a\in\R$）。

  分位函数：$F_{\mu}^{-1}(t)=t$，$F_{\nu}^{-1}(t)=a+t$。代入一维公式：
  \[
    \Wp{p}(\mu,\nu)^{p}
    = \int_{0}^{1}\abs{t-(a+t)}^{p}\dd t
    = \int_{0}^{1}\abs{a}^{p}\dd t
    = \abs{a}^{p},
  \]
  故 $\Wp{p}(\mu,\nu)=\abs{a}$，与 $p$ 无关。

  \textbf{含义。} 两个均匀分布相差平移 $a$，最优方案就是每个质量点向右（或左）平移 $a$。
  Wasserstein 距离恰好等于平移量 $\abs{a}$，与 $p$ 无关。

  \TLtakeaway{均匀分布平移量 $a$ 的 Wasserstein 距离等于 $|a|$，与 $p$ 无关——单纯平移的代价就是平移距离。}
\end{frame}

% ── 计算例（二）：两个高斯之间的 W_2 ────────────────────────────────────────
\begin{frame}{计算例（二）：高斯分布之间的 $\Wtwo$（推导）}
  \deflab{一维高斯的 $\Wtwo$ 公式。}

  设 $\mu=\mathcal{N}(m_{1},\sigma_{1}^{2})$，$\nu=\mathcal{N}(m_{2},\sigma_{2}^{2})$（$\R$ 上），则
  \[
    \Wtwo(\mu,\nu)^{2} \;=\; (m_{1}-m_{2})^{2} + (\sigma_{1}-\sigma_{2})^{2}.
  \]

  \textbf{推导（分位公式）。} 令 $\Phi^{-1}$ 为标准正态分位函数，则
  $F_{\mu}^{-1}(t)=m_{1}+\sigma_{1}\Phi^{-1}(t)$，$F_{\nu}^{-1}(t)=m_{2}+\sigma_{2}\Phi^{-1}(t)$。代入：
  \[
    \Wtwo(\mu,\nu)^{2}
    = \int_{0}^{1}\bigl[(m_{1}-m_{2})+(\sigma_{1}-\sigma_{2})\Phi^{-1}(t)\bigr]^{2}\dd t.
  \]
  展开，利用 $\int_{0}^{1}\Phi^{-1}(t)\dd t=0$（均值零）与 $\int_{0}^{1}[\Phi^{-1}(t)]^{2}\dd t=1$（方差一）：
  \[
    = (m_{1}-m_{2})^{2} + 2(m_{1}-m_{2})(\sigma_{1}-\sigma_{2})\cdot0
      + (\sigma_{1}-\sigma_{2})^{2}\cdot1
    = (m_{1}-m_{2})^{2} + (\sigma_{1}-\sigma_{2})^{2}.\;\square
  \]

  \TLtakeaway{$\Wtwo^2=(m_1-m_2)^2+(\sigma_1-\sigma_2)^2$：均值偏移与尺度偏移相互独立地贡献。}
\end{frame}

% ── 计算例（三）：一维最优单调匹配图示 ──────────────────────────────────────
\begin{frame}{计算例图示：一维密度的最优单调匹配}
  \begin{columns}[T]
    \begin{column}{0.54\textwidth}
      \centering
      \begin{tikzpicture}[scale=0.82,>={Stealth[length=4pt]}]
        \draw[->,TLink,thin] (-0.3,0) -- (4.8,0) node[right,font=\scriptsize] {$x$};
        \draw[->,TLink,thin] (0,-0.15) -- (0,2.1);
        % mu: Uniform[0,2]
        \draw[TLprimary,very thick] (0,0) -- (0,1.35) -- (2,1.35) -- (2,0);
        \node[TLprimary,font=\scriptsize,above] at (1,1.35) {$\mu$};
        % nu: Uniform[1.5,3.5]
        \draw[TLaccent,very thick] (1.5,0) -- (1.5,1.35) -- (3.5,1.35) -- (3.5,0);
        \node[TLaccent,font=\scriptsize,above] at (2.5,1.35) {$\nu$};
        % 单调匹配箭头
        \foreach \x/\y in {0.2/1.7, 0.6/2.1, 1.0/2.5, 1.4/2.9, 1.8/3.3}{
          \draw[->,TLinkSoft,thin,dashed] (\x,0.08) -- (\y,0.08);
        }
        \node[font=\scriptsize,TLinkSoft] at (1.8,0.55) {单调匹配};
      \end{tikzpicture}\\[3pt]
      {\scriptsize 据 Villani 图 6.1 重绘。}
    \end{column}
    \begin{column}{0.43\textwidth}
      \small
      \textbf{设置。} $\mu=\mathrm{Uniform}[0,2]$，$\nu=\mathrm{Uniform}[1.5,3.5]$，$p=1$。

      分位函数：$F_{\mu}^{-1}(t)=2t$，$F_{\nu}^{-1}(t)=1.5+2t$。

      \[
        \Wp{1}(\mu,\nu)=\int_{0}^{1}1.5\,\dd t=1.5.
      \]

      最优方案：平移 $T(x)=x+1.5$。每个质量点向右移 $1.5$ 单位。
    \end{column}
  \end{columns}
  \TLtakeaway{一维最优传输即单调重排；每个分位点点对点对齐，Wasserstein 距离是分位差的 $L^{p}$ 范数。}
\end{frame}

% ── 定义 6.8 ─────────────────────────────────────────────────────────────────
\begin{frame}{定义 6.8：$\Ppp{p}(\X)$ 中的弱收敛（等价条件）}
  \deflab{定义 6.8（Wasserstein 意义下的弱收敛，Villani Def 6.8）。}

  设 $(\mu_{k})\subset\Ppp{p}(\X)$，$\mu\in\Ppp{p}(\X)$，取定 $x_{0}\in\X$。以下等价：
  \begin{enumerate}\small
    \item[(i)] $\mu_{k}\wkto\mu$ 且 $\displaystyle\int_{\X} d(x_{0},x)^{p}\dd\mu_{k}\To\int_{\X} d(x_{0},x)^{p}\dd\mu$。
    \item[(ii)] $\mu_{k}\wkto\mu$ 且 $\displaystyle\limsup_{k}\int_{\X} d(x_{0},x)^{p}\dd\mu_{k} \le \int_{\X} d(x_{0},x)^{p}\dd\mu$。
    \item[(iii)] $\mu_{k}\wkto\mu$ 且 $\displaystyle\lim_{R\to\infty}\limsup_{k}\int_{d(x_{0},x)\ge R} d(x_{0},x)^{p}\dd\mu_{k}=0$（质量不逃逸）。
    \item[(iv)] 对所有 $|\varphi(x)|\le C(1+d(x_{0},x)^{p})$ 的连续 $\varphi$，$\int\varphi\dd\mu_{k}\To\int\varphi\dd\mu$。
  \end{enumerate}

  \TLtakeaway{$\Ppp{p}$ 中弱收敛 = 普通弱收敛 + $p$ 阶矩收敛（等价条件 (i)）；条件 (iii) 意为矩均匀可积。}
\end{frame}

% ── 定理 6.9：陈述 ──────────────────────────────────────────────────────────
\begin{frame}{定理 6.9：$\Wp{p}$ 刻画 $\Ppp{p}$ 中的弱收敛}
  \deflab{定理 6.9（$\Wp{p}$ 距离化 $\Ppp{p}$，Villani Thm 6.9）。}

  设 $(\X,d)$ 为 Polish 空间，$p\in[1,\infty)$，$(\mu_{k})\subset\Ppp{p}(\X)$，$\mu\in\Pp(\X)$。则
  \[
    \Wp{p}(\mu_{k},\mu)\To 0
    \;\iff\;
    (\mu_{k})\;\text{在}\;\Ppp{p}(\X)\;\text{中弱收敛到}\;\mu.
  \]

  \textbf{重要性。} Wasserstein 收敛比普通弱收敛更强——同时要求 $p$ 阶矩收敛。

  \textbf{反例。} 令 $\mu_{k}=(1-1/k)\delta_{0}+(1/k)\delta_{k^{2}}$，$\mu=\delta_{0}$。则 $\mu_{k}\wkto\mu$，但 $\int|x|\dd\mu_{k}=k\To\infty\neq0$，故 $\Wp{1}(\mu_{k},\mu)\To\infty$。

  \TLtakeaway{定理 6.9：$\Wp{p}$ 收敛等价于弱收敛加 $p$ 阶矩收敛，Wasserstein 拓扑严格强于弱拓扑。}
\end{frame}

% ── 定理 6.9 证明（一）──────────────────────────────────────────────────────
\begin{frame}{定理 6.9 的证明（一）：$\Wp{p}\to0$ 蕴含弱收敛与矩收敛}
  \textbf{方向 $(\Rightarrow)$：} 设 $\Wp{p}(\mu_{k},\mu)\To0$。

  \textbf{弱收敛。} 由引理 6.14，$(\mu_{k})$ 胶着，存在弱收敛子列 $\mu_{k'}\wkto\mu'$；用引理 4.3 知 $\Wp{p}(\mu',\mu)=0$，故 $\mu'=\mu$。

  \textbf{矩的上界。} 设 $\pi_{k}$ 为 $(\mu_{k},\mu)$ 的最优耦合。对 $\eps>0$ 存在 $C_{\eps}>0$ 使
  \[
    d(x_{0},x)^{p} \;\le\; (1+\eps)d(x_{0},y)^{p} + C_{\eps}d(x,y)^{p}.
  \]
  对 $\pi_{k}$ 积分（利用边际性质）：
  \[
    \int d(x_{0},x)^{p}\dd\mu_{k}
    \;\le\; (1+\eps)\int d(x_{0},y)^{p}\dd\mu + C_{\eps}\Wp{p}(\mu_{k},\mu)^{p}.
  \]
  令 $k\to\infty$ 再令 $\eps\to0$：$\limsup_{k}\int d(x_{0},x)^{p}\dd\mu_{k}\le\int d(x_{0},x)^{p}\dd\mu$（条件 (ii) 成立）。$\square$

  \TLtakeaway{$(\Rightarrow)$：胶着性保证弱极限；对不等式积分给出 $p$ 阶矩的上界，从而矩收敛。}
\end{frame}

% ── 定理 6.9 证明（二）──────────────────────────────────────────────────────
\begin{frame}{定理 6.9 的证明（二）：弱收敛加矩收敛蕴含 $\Wp{p}\to0$}
  \textbf{方向 $(\Leftarrow)$：} 设 $\mu_{k}\wkto\mu$ 且 $p$ 阶矩收敛。

  设 $\pi_{k}$ 为最优耦合；由引理 4.4 胶着，取子列 $\pi_{k}\wkto\pi$，定理 5.20 给出 $\pi=(\Id,\Id)_{\#}\mu$（对角耦合）。把 $\Wp{p}(\mu_k,\mu)^p\le\int d^p\dd\pi_k$ 拆成两段：
  \begin{itemize}\small
    \item \emph{截断项} $\int\min(d^p,R^p)\dd\pi_k\to\int\min(d^p,R^p)\dd\pi=0$：被积函数有界连续，$\pi_k\wkto\pi$ 且 $\pi$ 在对角线上（$d=0$）。
    \item \emph{尾部项} $\int(d^p-R^p)_+\dd\pi_k$：由 $(d^p{-}R^p)_+\le 2^p[d(x,x_0)^p\mathbf{1}_{d(x,x_0)\ge R/2}+d(x_0,y)^p\mathbf{1}_{d(x_0,y)\ge R/2}]$，先 $k\to\infty$ 再 $R\to\infty$ 趋于 $0$（矩均匀可积）。
  \end{itemize}
  两段相加：$\displaystyle\limsup_{k}\Wp{p}(\mu_{k},\mu)^{p} = 0$。\hfill$\square$

  \TLtakeaway{$(\Leftarrow)$：最优耦合弱极限是对角耦合；截断分解利用矩均匀可积，得 Wasserstein 收敛。}
\end{frame}

% ── 推论 6.11 与注记 6.10 ─────────────────────────────────────────────────────
\begin{frame}{注记 6.10 与推论 6.11：矩收敛与 $\Wp{p}$ 的连续性}
  \deflab{注记 6.10（矩收敛，Villani Rem 6.10）。}

  由定理 6.9，Wasserstein 收敛蕴含 $p$ 阶矩的收敛。更强地，在局部紧长度空间中，映射
  \[
    \mu\;\mapsto\;\Bigl(\int_{\X} d(x_{0},x)^{p}\dd\mu\Bigr)^{1/p}
  \]
  关于 $\Wp{p}$ 是 $1$-Lipschitz 的（命题 7.29）。

  \deflab{推论 6.11（连续性，Villani Cor 6.11）。}

  若 $\mu_{k}\To\mu$，$\nu_{k}\To\nu$ 均在 $\Ppp{p}(\X)$ 中，则
  $\Wp{p}(\mu_{k},\nu_{k})\To\Wp{p}(\mu,\nu)$。

  \textbf{证明。} 由三角不等式：
  \[
    |\Wp{p}(\mu_{k},\nu_{k})-\Wp{p}(\mu,\nu)|
    \le\Wp{p}(\mu_{k},\mu)+\Wp{p}(\nu_{k},\nu)\To0.\;\square
  \]

  \TLtakeaway{Wasserstein 收敛蕴含矩收敛；$\Wp{p}$ 关于两个参数均在 Wasserstein 拓扑下连续。}
\end{frame}

% ── 注记 6.12 与推论 6.13 ─────────────────────────────────────────────────────
\begin{frame}{注记 6.12 与推论 6.13：下半连续性与弱拓扑等价性}
  \deflab{注记 6.12（弱拓扑下仅下半连续，Villani Rem 6.12）。}

  若仅有普通弱收敛，则只能得到
  \[
    \Wp{p}(\mu,\nu) \;\le\; \liminf_{k\to\infty}\Wp{p}(\mu_{k},\nu_{k}).
  \]
  即：$\Wp{p}$ 在 $\Pp(\X)$ 上关于弱拓扑\emph{下半连续}，但一般不连续。

  \deflab{推论 6.13（有界距离下等价，Villani Cor 6.13）。}

  设 $\tilde{d}=d/(1+d)$（有界且与 $d$ 诱导相同拓扑）。对应于 $\tilde{d}$ 的 Wasserstein 收敛等价于 $\Pp(\X)$ 上的普通弱收敛。

  \textbf{原因。} 有界化后 $\int\tilde{d}(x_{0},x)^{p}\dd\mu\le1$ 对所有 $\mu$ 自动成立，矩条件退化为自动满足。

  \TLtakeaway{弱拓扑下 $\Wp{p}$ 下半连续（不连续）；有界距离下 Wasserstein 收敛等价于弱收敛。}
\end{frame}

% ── 为何偏爱 Wasserstein 距离 ─────────────────────────────────────────────────
\begin{frame}{为何偏爱 Wasserstein 距离？}
  弱收敛有多种度量化方式（Lévy–Prokhorov、有界 Lipschitz、弱-$*$ 等），Wasserstein 距离有以下优势（Villani pp.97--99）：

  \begin{enumerate}\small
    \item \textbf{计及大距离。} 与弱-$*$ 距离相比，$\Wp{p}$ 对质量逃逸到远处敏感，更精确。
    \item \textbf{天然适应 OT 框架。} 问题含输运结构（PDE、流体力学）时，Wasserstein 语言直接适用。
    \item \textbf{丰富对偶性。} $\Wp{1}$ 的 KR 对偶使计算转化为 Lipschitz 函数积分，技术上便利。
    \item \textbf{上界易构造。} 任意耦合给出 $\Wp{p}$ 上界；$C$-Lipschitz 映射 $f$ 诱导 $C$-Lipschitz 映射 $\mu\mapsto\pf{f}{\mu}$。
    \item \textbf{编码底层几何。} $x\mapsto\delta_{x}$ 等距嵌入（例 6.3）；Wasserstein 空间曲率反映 $\X$ 的曲率（§4 以后）。
  \end{enumerate}

  \TLtakeaway{Wasserstein 距离的价值：对大距离敏感 + KR 对偶易计算 + 自然编码底层几何，而非仅是弱收敛的又一度量化。}
\end{frame}

% ── 定理 6.15：全变差控制 ────────────────────────────────────────────────────
\begin{frame}{定理 6.15：全变差控制 Wasserstein 距离}
  \deflab{定理 6.15（全变差上界，Villani Thm 6.15）。}

  设 $\mu,\nu$ 为 $(\X,d)$ 上的概率测度，$p\in[1,\infty)$，$x_{0}\in\X$，$1/p+1/p'=1$。则
  \[
    \Wp{p}(\mu,\nu)
    \;\le\; 2^{1/p'}\!\left(\int_{\X} d(x_{0},x)^{p}\dd|\mu-\nu|(x)\right)^{1/p}.
  \]
  \deflab{特殊情形 6.16（有界直径）。} 若 $\X$ 直径 $\le D$，则 $\Wp{1}(\mu,\nu)\le D\cdot\|\mu-\nu\|_{TV}$。

  \textbf{证明思路。} 构造耦合
  $\pi=(\Id,\Id)_{\#}(\mu\wedge\nu)+\frac{1}{a}(\mu-\nu)^{+}\otimes(\mu-\nu)^{-}$（固定公共质量，剩余乘积耦合），再利用 $d(x,y)^{p}\le2^{p-1}(d(x,x_{0})^{p}+d(x_{0},y)^{p})$ 估计。

  \TLtakeaway{全变差控制 Wasserstein 须用 $d(x_0,x)^p$ 加权；裸全变差仅在有界空间中适用。}
\end{frame}

% ── §3 小结（一） ─────────────────────────────────────────────────────────────
\begin{frame}{§3 小结（一）：定义与距离公理}
  \begin{itemize}\small
    \item \textbf{定义 6.1（$\Wp{p}$）。} 成本 $d^{p}$ 的最优传输值取 $p$ 次根。满足距离三公理，三角不等式经粘合引理 + $L^{p}$ Minkowski 不等式证明。
    \item \textbf{例 6.3。} $\Wp{p}(\delta_{x},\delta_{y})=d(x,y)$：Dirac 测度等距嵌入。
    \item \textbf{定义 6.4（$\Ppp{p}(\X)$）。} $p$ 阶矩有限的概率测度；$\Wp{p}$ 在其上有限；$\Ppp{p}(\X)$ 本身是 Polish 空间（定理 6.18）。
    \item \textbf{注记 6.5。} KR 对偶：$\Wp{1}=\sup_{\|\psi\|_\mathrm{Lip}\le1}(\int\psi\dd\mu-\int\psi\dd\nu)$。
    \item \textbf{注记 6.6。} 阶序 $\Wp{1}\le\Wp{p}$；$p=1$ 最灵活，$p=2$ 最能反映几何。
  \end{itemize}

  \TLtakeaway{$\Wp{p}$ 将最优传输代价升格为真正的距离；$\Ppp{p}(\X)$ 是自然定义域，本身也是 Polish 空间。}
\end{frame}

% ── §3 小结（二） ─────────────────────────────────────────────────────────────
\begin{frame}{§3 小结（二）：收敛刻画与计算工具}
  \begin{itemize}\small
    \item \textbf{定理 6.9。} $\Wp{p}$ 收敛等价于弱收敛加 $p$ 阶矩收敛。
    \item \textbf{推论 6.11。} $\Wp{p}$ 关于两个参数在 Wasserstein 拓扑下连续。
    \item \textbf{注记 6.12。} 仅弱收敛时 $\Wp{p}$ 下半连续（不连续）。
    \item \textbf{推论 6.13。} 有界距离下 Wasserstein 收敛等价于弱收敛。
    \item \textbf{计算工具：} 一维分位公式；高斯 $\Wtwo^{2}=(m_1{-}m_2)^{2}+(\sigma_1{-}\sigma_2)^{2}$；定理 6.15 全变差上界。
  \end{itemize}

  \textbf{预告：} §4 利用 $(\Ppp{2}(\X),\Wtwo)$ 的 Polish 性质，构造测地线与 McCann 插值。

  \TLtakeaway{定理 6.9 是本章核心：$\Wp{p}$ 刻画弱收敛 + 矩收敛，为 §4 在概率测度空间上做几何提供度量基础。}
\end{frame}

% ── 习题 3 · Exercises ─────────────────────────────────────────────────────────
\begin{frame}{习题 3 · Exercises}
  \begin{itemize}
    \item \textbf{习题 3.1}~(\diffI)\
      设 $\mu=\mathrm{Bernoulli}(p_{0})$，$\nu=\mathrm{Bernoulli}(q_{0})$（在 $\{0,1\}$ 上，$0\le p_{0}\le q_{0}\le1$）。
      用一维分位公式计算 $\Wp{1}(\mu,\nu)$。
      \begin{itemize}\footnotesize
        \item \hint{$F_{\mu}^{-1}(t)=\mathbf{1}_{t\ge1-p_{0}}$；差在区间 $[1-q_{0},1-p_{0})$ 上为 $1$，其余为 $0$。}
      \end{itemize}
      % SOLUTION: F_mu^{-1}=0 on [0,1-p_0), 1 on [1-p_0,1]. F_nu^{-1}=0 on [0,1-q_0), 1 on [1-q_0,1]. 差在 [1-q_0,1-p_0) 上为 1，其余 0（p_0<=q_0 故 1-q_0<=1-p_0）。W_1(mu,nu)=int_0^1|差|dt=(1-p_0)-(1-q_0)=q_0-p_0。直觉：把 q_0-p_0 比例的"0"搬到"1"，每单位距离 1，总代价 q_0-p_0。

    \item \textbf{习题 3.2}~(\diffII)\
      构造 $\mu_{k}\wkto\mu$ 但 $\Wp{1}(\mu_{k},\mu)\to\infty$ 的例子，并说明违反了定理 6.9 的哪个条件。
      \begin{itemize}\footnotesize
        \item \hint{令 $\mu_{k}=(1-1/k)\delta_{0}+(1/k)\delta_{k^{2}}$，计算一阶矩；用 KR 对偶取检验函数 $\psi(x)=\min(x,R)$ 估计下界。}
      \end{itemize}
      % SOLUTION: mu_k 弱收敛到 delta_0（phi 有界时 (1-1/k)phi(0)+(1/k)phi(k^2)→phi(0)）。一阶矩 int|x|dmu_k=(1/k)k^2=k→∞，违反定理 6.9 条件 (i) 的矩收敛。取 psi(x)=min(x,R)/R（1/R-Lipschitz），KR 对偶：W_1(mu_k,delta_0)≥R·(1/k·min(k^2,R)/R - 0)；令 R=k 得 W_1≥k→∞。

    \item \textbf{习题 3.3}~(\diffIII)\
      设 $\mu=\mathcal{N}(0,I_{n})$，$\nu=\mathcal{N}(m,\Sigma)$（$\R^{n}$ 上，$\Sigma$ 正定）。
      利用多维高斯 $\Wtwo$ 公式
      $\Wtwo^{2}=|m_1{-}m_2|^{2}+\mathrm{tr}(\Sigma_1+\Sigma_2-2(\Sigma_1^{1/2}\Sigma_2\Sigma_1^{1/2})^{1/2})$，
      计算 $\Wtwo(\mu,\nu)^{2}$，并在 $n=1$ 时验证与正文高斯例吻合。
      \begin{itemize}\footnotesize
        \item \hint{代入 $\Sigma_1=I_n$；化简 $(I\Sigma I)^{1/2}=\Sigma^{1/2}$，再用 $\mathrm{tr}(I+\Sigma-2\Sigma^{1/2})=\mathrm{tr}((I-\Sigma^{1/2})^{2})$。}
      \end{itemize}
      % SOLUTION: 代入 Sigma_1=I_n, m_1=0。(I^{1/2}Sigma I^{1/2})^{1/2}=Sigma^{1/2}。W_2^2=|m|^2+tr(I+Sigma-2Sigma^{1/2})=|m|^2+tr((I-Sigma^{1/2})^2)=|m|^2+sum_i(1-sqrt(lambda_i))^2，lambda_i 为 Sigma 的特征值。当 n=1，Sigma=sigma^2，W_2^2=m^2+(1-sigma)^2，与正文（m_1=0,sigma_1=1,m_2=m,sigma_2=sigma）的 W_2^2=(m_1-m_2)^2+(sigma_1-sigma_2)^2 完全吻合。
  \end{itemize}
  \TLtakeaway{3.1 用分位差计算 Bernoulli 距离；3.2 构造矩发散的弱收敛反例；3.3 验证多维高斯公式。}
\end{frame}
 % §3  Ch.6 Wasserstein 距离
% =============================================================================
%  §4  Ch.7 位移插值（Displacement Interpolation）
%  来源：Villani, Optimal Transport: Old and New, Chapter 7 (pp. 114–160)
%  主要定理：定理 7.21
% =============================================================================
\TLsection{4 \ 位移插值}

% ── 动机（一）：线性插值的缺陷 ──────────────────────────────────────────────
\begin{frame}{动机：线性插值的缺陷}
  \textbf{自然问题。} 给定两个概率分布 $\mu_{0}$ 和 $\mu_{1}$，如何构造"合理的"中间分布？

  \textbf{线性插值：} $\mu_{t}^{\mathrm{lin}}=(1-t)\mu_{0}+t\mu_{1}$
  \begin{itemize}
    \item 数学上简单，但\emph{物理上不合理}：质量被"传送"（teleport）而非移动。
    \item $\mu_{0}=\delta_{0}$，$\mu_{1}=\delta_{1}$ 时，$\mu_{1/2}^{\mathrm{lin}}=\tfrac{1}{2}\delta_{0}+\tfrac{1}{2}\delta_{1}$——质量\emph{分裂}到两处。
    \item 两高斯分布的线性插值在 $t=\tfrac{1}{2}$ 产生\emph{双峰}（bimodal），与直觉相悖。
  \end{itemize}

  \textbf{位移插值（McCann 1995）：} 让每单位质量沿\TLterm{最优路径}从 $\mu_{0}$ 移动到 $\mu_{1}$；
  $\mu_{t}$ 是所有质量在 $t$ 时刻位置的分布。

  \TLtakeaway{线性插值"传送"（teleport）质量产生双峰；位移插值沿最优路径整体平移，保留分布形状。}
\end{frame}

% ── 动机（图示）：线性 vs 位移插值 ───────────────────────────────────────────
\begin{frame}{图示：线性插值 vs 位移插值（对比）}
  \begin{columns}[T]
    \begin{column}{0.48\textwidth}
      \centering
      \textbf{线性插值}（产生双峰）\\[4pt]
      \begin{tikzpicture}[scale=0.72,>={Stealth[length=3pt]}]
        % t=0：单峰在 x=0（与位移列同一端点）
        \draw[TLprimary,thick,domain=-1.0:1.0,smooth,samples=25] plot(\x,{1.0*exp(-3*\x*\x)});
        \node[TLprimary,font=\scriptsize,below] at (0,-0.1) {$t=0$};
        \begin{scope}[shift={(0,-1.8)}]
          % t=1/2：½在 0 + ½在 1.6，即两端点的平均 → 双峰
          \draw[TLprimary!50!TLaccent,thick,domain=-1.0:2.6,smooth,samples=50]
            plot(\x,{0.5*exp(-3*\x*\x)+0.5*exp(-3*(\x-1.6)*(\x-1.6))});
          \node[TLprimary!50!TLaccent,font=\scriptsize,below] at (0.8,-0.1) {$t=\tfrac12$（双峰）};
        \end{scope}
        \begin{scope}[shift={(0,-3.6)}]
          % t=1：单峰在 x=1.6（与位移列同一端点）
          \draw[TLaccent,thick,domain=0.6:2.6,smooth,samples=25] plot(\x,{1.0*exp(-3*(\x-1.6)*(\x-1.6))});
          \node[TLaccent,font=\scriptsize,below] at (1.6,-0.1) {$t=1$};
        \end{scope}
      \end{tikzpicture}
    \end{column}
    \begin{column}{0.48\textwidth}
      \centering
      \textbf{位移插值}（平滑平移）\\[4pt]
      \begin{tikzpicture}[scale=0.72,>={Stealth[length=3pt]}]
        \draw[TLprimary,thick,domain=-1.0:1.0,smooth,samples=25]
          plot(\x,{1.0*exp(-3*\x*\x)});
        \node[TLprimary,font=\scriptsize,below] at (0,-0.1) {$t=0$};
        \begin{scope}[shift={(0,-1.8)}]
          \draw[TLprimary!50!TLaccent,thick,domain=-0.2:1.8,smooth,samples=25]
            plot(\x,{1.0*exp(-3*(\x-0.8)*(\x-0.8))});
          \node[TLprimary!50!TLaccent,font=\scriptsize,below] at (0.8,-0.1) {$t=\tfrac12$（单峰）};
        \end{scope}
        \begin{scope}[shift={(0,-3.6)}]
          \draw[TLaccent,thick,domain=0.6:2.6,smooth,samples=25]
            plot(\x,{1.0*exp(-3*(\x-1.6)*(\x-1.6))});
          \node[TLaccent,font=\scriptsize,below] at (1.6,-0.1) {$t=1$};
        \end{scope}
      \end{tikzpicture}
    \end{column}
  \end{columns}
  \vspace{2pt}
  {\scriptsize 据 Villani 重绘。}

  \TLtakeaway{位移插值中单峰始终保持单峰并平滑平移；线性插值在中间产生双峰——这正是最优传输几何直觉的来源。}
\end{frame}

% ── 作用量泛函与 Lagrange 量 ─────────────────────────────────────────────────
\begin{frame}{作用量泛函与 Lagrange 量（7.2）}
  \deflab{作用量泛函（Villani 第 7 章）。}
  \deflab{补一个事实：切丛。} $TM$ 把位置 $x\in M$ 与速度 $v\in T_{x}M$ 配成点 $(x,v)$。
  \deflab{补一个事实：Lagrange 量。} $L(x,v,t)$ 表示时刻 $t$ 以速度 $v$ 经过 $x$ 的瞬时运动成本。
  \TLterm{作用量}把瞬时成本沿曲线累加：
  \[
    A(\gamma) = \int_{0}^{1} L(\gamma_{t},\dot{\gamma}_{t},t)\,\dd t.
  \]
  \textbf{典型形式：} $L=\tfrac{|v|^{2}}{2}$，$A(\gamma)=\int_{0}^{1}\tfrac{|\dot{\gamma}_{t}|^{2}}{2}\dd t$（动能积分）。

  \textbf{最小化曲线与成本：} 同端点下 $A(\gamma)\le A(\tilde{\gamma})$；关联成本
  $c(x,y):=\inf_{\gamma_{0}=x,\,\gamma_{1}=y}A(\gamma)$。
  \textbf{例 7.1。} $L=|v|$ 时 $A(\gamma)=\mathrm{length}(\gamma)$，$c(x,y)=|x-y|$。

  \TLtakeaway{作用量 = Lagrange 量的时间积分；最小化对应"最省力"路径；关联成本 $c(x,y)$ 是最优作用量值。}
\end{frame}

% ── 例 7.2 与例 7.4 ──────────────────────────────────────────────────────────
\begin{frame}{例 7.2 \& 7.4：凸 Lagrange 量与 Riemann 测地线}
  \deflab{例 7.2（Villani Example 7.2）。} 取 $L(x,v,t)=c(v)$，$c$ \emph{严格凸}，则最小化曲线具有\emph{常速度}。

  \textbf{证明（Jensen 不等式）：}
  \[
    A(\gamma)=\int_{0}^{1}c(\dot{\gamma}_{t})\dd t
    \;\ge\; c\!\left(\int_{0}^{1}\dot{\gamma}_{t}\dd t\right)
    = c(\gamma_{1}-\gamma_{0}).
  \]
  等号成立当且仅当 $\dot{\gamma}_{t}=\mathrm{const}$（常速直线）。$\square$

  \vspace{3pt}
  \deflab{例 7.4（Villani Example 7.4）。} 设 $M$ 为 Riemann 流形，取 $L(x,v,t)=|v|^{p}$，$p\ge1$：
  \begin{itemize}
    \item \textbf{$p>1$（严格凸）：} 最小化曲线是\TLterm{常速测地线}，满足
      $d(\gamma_{s},\gamma_{t})=|t-s|\cdot d(\gamma_{0},\gamma_{1})$，关联成本 $c=d^{p}$。
    \item \textbf{$p=1$（线性，非严格凸）：} 最小化曲线是测地曲线，允许任意重参数化。
  \end{itemize}

  \TLtakeaway{严格凸 Lagrange 量（$p>1$）通过 Jensen 不等式迫使最小化曲线等速；$c(x,y)=d(x,y)^p$ 是 Wasserstein 距离 $\Wp{p}$ 的成本函数。}
\end{frame}

% ── 度量空间中的测地线 ──────────────────────────────────────────────────────
\begin{frame}{度量空间中的测地线（例 7.8–7.9）}
  \deflab{例 7.8（Villani Example 7.8）。}
  设 $(\X,d)$ 为度量空间，曲线 $\gamma:[0,1]\to\X$ 的\TLterm{长度}为
  \[
    \mathrm{length}(\gamma):=\sup_{N}\sup_{0=t_{0}<\cdots<t_{N}=1}
    \sum_{i=0}^{N-1}d(\gamma_{t_{i}},\gamma_{t_{i+1}}).
  \]
  \TLterm{测地线}：长度等于端点距离的曲线（可变速，可重参数化为等速）。

  \deflab{例 7.9（Villani Example 7.9）。} 取 $L(x,v,t)=|v|^{p}$（$p>1$）。则

  $\gamma$ 为\TLterm{常速最小测地线} $\iff$ $d(\gamma_{s},\gamma_{t})=|t-s|\cdot d(\gamma_{0},\gamma_{1})$，$\forall s,t$.

  关联成本：$c(x,y)=d(x,y)^{p}$，最小化曲线是常速最小测地线。

  \deflab{测地长度空间（Villani p.298）。}
  若 $(\X,d)$ 中任意两点之间存在最短测地线，则称其为\TLterm{测地长度空间}。
  例：$\R^{n}$（直线段）、完备 Riemann 流形（Hopf–Rinow 定理）。

  \TLtakeaway{$p>1$ 的 $L^p$ 作用量选出常速最小测地线；测地长度空间保证最短路存在——二者共同构成位移插值理论的空间假设。}
\end{frame}

% ── 位移插值的构造 ────────────────────────────────────────────────────────────
\begin{frame}{位移插值的构造（一）：路径空间测度}
  \textbf{设置：} $(\X,d)$ 测地长度空间，$\mu_{0},\mu_{1}\in\Ppp{p}(\X)$，
  $\pi$ 为 $(\mu_{0},\mu_{1})$ 最优耦合（成本 $c=d^{p}$）。

  \textbf{第一步：} 对每对 $(x,y)$，选常速最小测地线 $\gamma^{x,y}:[0,1]\to\X$，
  满足 $d(\gamma^{x,y}_{s},\gamma^{x,y}_{t})=|t-s|\cdot d(x,y)$。

  \textbf{第二步：} 令 $e_{t}(\gamma):=\gamma(t)$，构造路径空间概率测度
  \[
    \Pi:=\int_{\X\times\X}\delta_{\gamma^{x,y}}\dd\pi(x,y).
  \]

  \textbf{第三步：} \TLterm{位移插值}为 $\mu_{t}:=(e_{t})_{\#}\Pi$。

  \TLtakeaway{位移插值 = 沿最优耦合对每对质量点取测地线后推前到 $t$ 时刻；$\Pi$ 是"随机测地线"的分布。}
\end{frame}

\begin{frame}{位移插值的构造（二）：$\R^{n}$ 中的显式公式}
  \textbf{特殊情形：} $\X=\R^{n}$，$c(x,y)=|x-y|^{2}$。当 $\mu_0\in\Pac(\R^n)$（绝对连续，见 §0）时，Brenier 定理（§2）保证最优映射 $T=\nabla\varphi$ 存在。

  直线段 $\gamma^{x,y}_{t}=(1-t)x+tT(x)$ 是常速最小测地线（$\R^n$ 中直线唯一），故
  \[
    \mu_{t}=\bigl((1-t)\Id+tT\bigr)_{\#}\mu_{0}.
  \]
  每质量点 $x$ 从 $x$ 以常速移向 $T(x)$，$t$ 时刻位于 $(1-t)x+tT(x)$。

  \begin{center}
  \begin{tikzpicture}[scale=1.0,>={Stealth[length=4pt]}]
    \fill[TLprimary] (0,0) circle (3pt) node[below,font=\scriptsize] {$x$};
    \fill[TLaccent]  (4,0.6) circle (3pt) node[below,font=\scriptsize] {$T(x)$};
    \draw[->,TLprimaryDark,thick] (0,0) -- (4,0.6);
    \fill[TLprimary!60!TLaccent] (2,0.3) circle (2.5pt) node[above,font=\scriptsize] {$\gamma_{1/2}^x$};
  \end{tikzpicture}
  \end{center}
  {\scriptsize 据 Villani 重绘。}

  \TLtakeaway{$\R^n$ 中的位移插值：每质量点沿直线匀速移动，$\mu_t=((1-t)\Id+tT)_\#\mu_0$；直线段唯一性保证构造的唯一性。}
\end{frame}

% ── 定义 7.19–7.20：动力学耦合 ──────────────────────────────────────────────
\begin{frame}{定义 7.19–7.20：动力学耦合与动力学最优耦合}
  \deflab{定义 7.19（动力学转移计划，Villani Def 7.19）。}
  \TLterm{动力学转移计划}是路径空间 $C([0,1];\X)$ 上的概率测度 $\Pi$。$\Pi$ 称为 $(\mu_{0},\mu_{1})$ 的\TLterm{动力学耦合}若
  $(e_{0})_{\#}\Pi=\mu_{0}$，$(e_{1})_{\#}\Pi=\mu_{1}$。

  \deflab{定义 7.20（动力学最优耦合，Villani Def 7.20）。}
  \TLterm{动力学最优转移计划}是集中在作用量最小化曲线全体 $\Gamma$ 上的概率测度 $\Pi$，使得
  $\pi_{0,1}:=(e_{0},e_{1})_{\#}\Pi$ 是 $(\mu_{0},\mu_{1})$ 的最优（静态）转移计划。

  等价地，$\Pi$ 是"随机作用量最小化曲线"的分布，端点构成最优耦合。

  \textbf{关系（路径空间 $\to$ 联合分布 $\to$ 边际）：}
  \begin{center}
  \begin{tikzpicture}[scale=0.85,>={Stealth[length=4pt]}]
    \node[draw=TLprimary,rounded corners=3pt,inner sep=5pt,font=\small] (A) at (0,0) {$\Pi\in\Pp(C([0,1];\X))$};
    \node[draw=TLaccent,rounded corners=3pt,inner sep=5pt,font=\small] (B) at (5,0) {$\pi_{0,1}=(e_0,e_1)_\#\Pi$};
    \draw[->,TLinkSoft,thick] (A) -- (B) node[midway,above,font=\scriptsize] {推前};
    \node[font=\scriptsize,TLinkSoft,below=6pt of A] {随机测地线};
    \node[font=\scriptsize,TLinkSoft,below=6pt of B] {最优静态耦合};
  \end{tikzpicture}
  \end{center}

  \TLtakeaway{$\Pi$（路径空间）$\to$ $\pi_{0,1}$（端点耦合，静态最优）$\to$ $\mu_t=(e_t)_\#\Pi$（位移插值）。}
\end{frame}

% ── 定理 7.21（陈述）────────────────────────────────────────────────────────
\begin{frame}{定理 7.21：位移插值（陈述）}
  \deflab{定理 7.21（位移插值，Villani Thm 7.21）。}

  设 $(\X,d)$ 为 Polish 空间，$(A)^{0,1}$ 为强制 Lagrange 作用量，成本 $c^{s,t}$ 连续，
  $C(\mu_{0},\mu_{1})<+\infty$。对连续路径 $(\mu_{t})_{0\le t\le1}$，以下三条\emph{等价}：
  \begin{enumerate}
    \item[(i)] $\mu_{t}=(e_{t})_{\#}\Pi$，其中 $\Pi$ 为某动力学最优耦合。
    \item[(ii)] 对任意 $t_{1}<t_{2}<t_{3}\in[0,1]$：
      $C^{t_{1},t_{2}}(\mu_{t_{1}},\mu_{t_{2}})+C^{t_{2},t_{3}}(\mu_{t_{2}},\mu_{t_{3}})
      =C^{t_{1},t_{3}}(\mu_{t_{1}},\mu_{t_{3}})$。
    \item[(iii)] $(\mu_{t})$ 是 $\Pp(\X)$ 上作用量泛函 $\mathcal{A}^{s,t}(\mu)$ 的最小化路径。
  \end{enumerate}
  满足上述条件的 $(\mu_{t})$ 称为\TLterm{位移插值}，至少存在一条。

  \TLtakeaway{定理 7.21 三条等价刻画：(i) 动力学最优耦合的边际，(ii) 测地型加法等式，(iii) 作用量最小化路径——三者互相等价。}
\end{frame}

% ── 推论 7.22（一）：陈述 ────────────────────────────────────────────────────
\begin{frame}{推论 7.22（一）：位移插值是 Wasserstein 测地线}
  \deflab{推论 7.22（Villani Corollary 7.22）。}

  设 $(\X,d)$ 为完备可分局部紧测地长度空间，$p>1$，$\mu_{0},\mu_{1}\in\Ppp{p}(\X)$。

  对连续路径 $(\mu_{t})_{0\le t\le1}\subset\Pp(\X)$，以下\emph{等价}：
  \begin{enumerate}
    \item[(i)] $\mu_{t}=(e_{t})_{\#}\Pi$，$\Pi$ 集中在常速最小测地线上且 $(e_{0},e_{1})_{\#}\Pi$ 为最优耦合。
    \item[(ii)] $(\mu_{t})_{0\le t\le1}$ 是 $\bigl(\Ppp{p}(\X),\Wp{p}\bigr)$ 中的\TLterm{常速测地线}：
      \[
        \Wp{p}(\mu_{s},\mu_{t})=|t-s|\cdot\Wp{p}(\mu_{0},\mu_{1}),\quad\forall s,t\in[0,1].
      \]
  \end{enumerate}
  对任意 $\mu_{0},\mu_{1}$，至少存在一条这样的测地线。

  \TLtakeaway{推论 7.22：Wasserstein 测地线 = "测地线的分布"（$\mu_t=\law(\gamma_t)$，$\gamma$ 为随机最优测地线）——联结最优传输与 Wasserstein 几何。}
\end{frame}

% ── 推论 7.22（二）：核心口号与意义 ─────────────────────────────────────────
\begin{frame}{推论 7.22（二）：核心口号与意义}
  \textbf{核心口号：}
  \[
    (\mu_{t})\text{ 是 }\Ppp{p}(\X)\text{ 中的测地线}
    \;\iff\;
    \mu_{t}=\law(\gamma_{t}),
  \]
  其中 $\gamma$ 是 $\X$ 中的随机常速最小测地线，端点构成最优耦合。

  \textbf{深层含义：}
  \begin{itemize}
    \item 最优传输（静态）$\longleftrightarrow$ Wasserstein 空间中的测地线（动态）。
    \item 推论 7.22 回答"$\Ppp{p}(\X)$ 中的直线是什么？"——就是位移插值！
    \item 后续章节（Ricci 曲率下界、位移凸性）以此为出发点。
  \end{itemize}

  \deflab{推论 7.23（唯一性，Villani Cor 7.23）。}
  若 (a) 最优传输计划唯一；(b) $\pi$-a.s. 两端点间最小化曲线唯一，则位移插值\emph{唯一}。

  \TLtakeaway{最优传输 $\iff$ Wasserstein 测地线：是定理 7.21 最深刻的几何含义，为后续章节奠定基础。}
\end{frame}

% ── 定理 7.21 证明思路 ────────────────────────────────────────────────────────
\begin{frame}{定理 7.21 证明思路（一）：三时刻引理}
  \textbf{目标。}\proofskipnote 证明 (ii) $\Rightarrow$ (i)：由三时刻加法等式构造路径空间上的动力学最优耦合 $\Pi$。
  \deflab{补一个事实：三时刻粘合。}
  若两个相邻最优耦合共享中间边际 $\mu_{t_{2}}$，则 Gluing Lemma（§0/§3）给出三元随机变量
  $(\gamma_{t_{1}},\gamma_{t_{2}},\gamma_{t_{3}})$，并保留两段边际耦合。
  \deflab{关键引理。}
  若 $(\gamma_{t_{1}},\gamma_{t_{2}})$ 与 $(\gamma_{t_{2}},\gamma_{t_{3}})$ 各自最优，且 (ii) 成立，则
  $(\gamma_{t_{1}},\gamma_{t_{3}})$ 也是最优耦合：
  \[
    \mathbb{E}\,c^{t_{1},t_{3}}(\gamma_{t_{1}},\gamma_{t_{3}})
    \le C^{t_{1},t_{2}}+C^{t_{2},t_{3}}
    \overset{\mathrm{(ii)}}{=}C^{t_{1},t_{3}}.
  \]
  左端又不小于最优值 $C^{t_{1},t_{3}}$，故只能处处等号，端点耦合最优。
  \deflab{补一个事实：限制性质。}
  {\small 同一条最小化曲线穿过中点时，它在每个子区间仍最小化；这是作用量最小化曲线的基本性质——最小化曲线的任意子段仍是最小化曲线（restriction/Markov property）。}
  \textbf{反向。} (i) $\Rightarrow$ (ii) 由 $\Pi$-a.s. 的最小化曲线直接积分计算得到。
  \TLtakeaway{三时刻引理把局部最优耦合和加法等式升级为跨端点最优；这是二进构造能够保持全局最优性的核心。}
\end{frame}

\begin{frame}{定理 7.21 证明思路（二）：二进构造与极限}
  \deflab{补一个事实：二进时刻。}
  二进时刻是 $j/2^{k}$；它们在 $[0,1]$ 中稠密，所以先控制这些时刻，再用连续性传到所有时刻。\proofskipnote
  \textbf{二进逼近（由 (ii) 构造 (i)）：}
  \begin{enumerate}
    \item 取 $(\mu_{0},\mu_{1})$ 的最优耦合 $\pi^{(0)}$。
    \item 在 $0,\tfrac{1}{2},1$ 上粘合两段最优耦合；上一帧引理保证端点耦合仍最优。
    \item 归纳到所有 $j/2^{k}$：反复粘合三时刻最优耦合，得到一致的离散随机路径。
    \item 在相邻二进点之间选择作用量最小化曲线，得到分段最小化路径的法律 $\Pi^{(k)}$。
  \end{enumerate}
  \deflab{本帧暂缓的技术点。}
  从 $\Pi^{(k)}$ 过渡到 $\Pi\in\Pp(C([0,1];\X))$ 需要测地线的可测选择，以及路径空间上的 tightness/compactness；Villani 第 7 章处理这些细节，主线只需二进边际一致且极限仍支撑在最小化曲线上。
  \TLtakeaway{(ii) $\Rightarrow$ (i) 的实质是：二进时刻反复粘合三时刻最优耦合，再用路径空间紧性取得连续随机最小化曲线。}
\end{frame}

% ── 计算例（一）：两个 Dirac 测度 ──────────────────────────────────────────
\begin{frame}{计算例（一）：两个 Dirac 测度的位移插值}
  \textbf{设置：} $\mu_{0}=\delta_{x}$，$\mu_{1}=\delta_{y}$，$x,y\in\X$，$p>1$。

  \textbf{最优耦合：} $\coup{\delta_{x}}{\delta_{y}}=\{\delta_{(x,y)}\}$（唯一），故 $\pi=\delta_{(x,y)}$，
  路径空间测度 $\Pi=\delta_{\gamma^{x,y}}$（唯一测地线上的 Dirac 测度）。

  \textbf{位移插值：}
  \[
    \mu_{t}=(e_{t})_{\#}\delta_{\gamma^{x,y}}=\delta_{\gamma^{x,y}_{t}}.
  \]
  中间分布仍是 Dirac 测度，集中在测地线 $\gamma^{x,y}$ 的 $t$ 时刻位置。

  \textbf{验证推论 7.22：}
  \[
    \Wp{p}(\mu_{s},\mu_{t})
    =\Wp{p}(\delta_{\gamma_{s}},\delta_{\gamma_{t}})
    =d(\gamma_{s},\gamma_{t})
    =(t-s)\,d(x,y)
    =(t-s)\,\Wp{p}(\mu_{0},\mu_{1}).
  \]
  常速测地线条件完美成立。

  \TLtakeaway{Dirac 测度的位移插值：中间分布仍为 Dirac 测度集中在测地线上——是定理 7.21 最简单的完整验证。}
\end{frame}

% ── 计算例（二）：一维高斯分布 ───────────────────────────────────────────────
\begin{frame}{计算例（二）：$\R$ 上两个高斯分布的位移插值}
  \textbf{设置：} $\mu_{0}=\mathcal{N}(m_{0},\sigma_{0}^{2})$，$\mu_{1}=\mathcal{N}(m_{1},\sigma_{1}^{2})$，$c(x,y)=(x-y)^{2}$。

  \textbf{最优映射：} 一维高斯的单调最优映射为仿射映射
  $T(x)=m_{1}+\tfrac{\sigma_{1}}{\sigma_{0}}(x-m_{0})$。

  \textbf{位移插值的显式公式：} $T_{t}:=(1-t)\Id+tT$，则
  \[
    \mu_{t}=(T_{t})_{\#}\mu_{0}=\mathcal{N}(m_{t},\sigma_{t}^{2}),
  \]
  其中
  \[
    m_{t}=(1-t)m_{0}+tm_{1},\qquad
    \sigma_{t}=(1-t)\sigma_{0}+t\sigma_{1}.
  \]

  \textbf{要点：} 中间分布仍是高斯！均值和标准差均线性插值，\emph{不出现双峰}。

  \textbf{对比线性插值：} $(1-t)\mathcal{N}(m_{0},\sigma_{0}^{2})+t\mathcal{N}(m_{1},\sigma_{1}^{2})$
  在 $m_{0}\ne m_{1}$ 时产生双峰高斯混合。

  \TLtakeaway{高斯的位移插值保持高斯性，均值和标准差均线性变化——与产生双峰的线性插值形成鲜明对比。}
\end{frame}

% ── 计算例（二）续：计算验证 ────────────────────────────────────────────────
\begin{frame}{计算例（二）续：高斯插值的验证}
  \textbf{验证 $\mu_{t}=\mathcal{N}(m_{t},\sigma_{t}^{2})$。}

  映射 $T_{t}(x)=(1-t)x+t\bigl(m_{1}+\tfrac{\sigma_{1}}{\sigma_{0}}(x-m_{0})\bigr)$ 整理为
  \[
    T_{t}(x)=\underbrace{\bigl(1-t+t\tfrac{\sigma_{1}}{\sigma_{0}}\bigr)}_{\alpha_{t}}x
             +\underbrace{t\bigl(m_{1}-\tfrac{\sigma_{1}}{\sigma_{0}}m_{0}\bigr)}_{\beta_{t}}.
  \]
  这是仿射映射，故 $(T_{t})_{\#}\mu_{0}$ 为正态分布。计算参数：
  \begin{align*}
    m_{t} &= T_{t}(m_{0})=(1-t)m_{0}+tm_{1},\\
    \sigma_{t}^{2} &= \alpha_{t}^{2}\,\sigma_{0}^{2}
    =\bigl((1-t)\sigma_{0}+t\sigma_{1}\bigr)^{2}.\quad\square
  \end{align*}

  \textbf{$W_{2}$ 距离：} $\Wtwo(\mu_{0},\mu_{1})^{2}=(m_{1}-m_{0})^{2}+(\sigma_{1}-\sigma_{0})^{2}$。

  \TLtakeaway{仿射映射的推前保持高斯性；方差公式 $\sigma_t=(1-t)\sigma_0+t\sigma_1$ 直接来自仿射斜率 $\alpha_t$。}
\end{frame}

% ── 计算例（三）：平移 ──────────────────────────────────────────────────────
\begin{frame}{计算例（三）：平移的位移插值}
  \textbf{设置：} $\mu_{1}=(\tau_{a})_{\#}\mu_{0}$，$\tau_{a}(x)=x+a$，$a\in\R^{n}$，$c=|\cdot|^{2}$。

  \textbf{最优映射：} $T(x)=x+a$（Brenier 定理：仿射映射为最优）。

  \textbf{位移插值：} $T_{t}(x)=x+ta$，故
  \[
    \mu_{t}=(\tau_{ta})_{\#}\mu_{0}.
  \]
  全体质量同步匀速平移距离 $ta$——整个分布形状不变，仅位置变化。

  \textbf{验证常速测地线条件：}
  \[
    \Wp{2}(\mu_{s},\mu_{t})
    =\Wp{2}\bigl((\tau_{sa})_{\#}\mu_{0},(\tau_{ta})_{\#}\mu_{0}\bigr)
    =|(t-s)a|
    =|t-s|\cdot\Wp{2}(\mu_{0},\mu_{1}).\quad\square
  \]

  \textbf{对比（与 Dirac 例子）：} 若 $\mu_{0}=\delta_{0}$，则 $\mu_{t}=\delta_{ta}$（Dirac 测度沿直线运动）；
  一般分布同样以同样速度整体平移，$\Wp{2}(\mu_{0},\mu_{1})=|a|$。

  \TLtakeaway{平移：整个分布匀速平移 $ta$，每质量点走直线——最直观的常速测地线验证，$\Wp{2}(\mu_0,\mu_1)=|a|$。}
\end{frame}

% ── Benamou–Brenier 前瞻 ─────────────────────────────────────────────────────
\begin{frame}{Benamou–Brenier 公式（前瞻）}
  \textbf{动力学视角（Benamou–Brenier，1999）。} 在 $\R^{n}$ 中，寻找速度场 $v$ 与路径 $(\mu_{t})$。

  \deflab{补一个事实：连续性方程。}
  $\partial_{t}\mu_{t}+\nabla\cdot(v_{t}\mu_{t})=0$ 表示质量守恒：速度场 $v_{t}$ 只搬运质量，不创造也不消灭质量。

  在此约束下最小化动能
  $\mathcal{A}(\mu,v)=\int_{0}^{1}\!\int_{\R^{n}}|v_{t}|^{2}\dd\mu_{t}\dd t$。

  \deflab{Benamou–Brenier 公式（Villani §7 文献注记）。}
  \[
    \Wtwo(\mu_{0},\mu_{1})^{2}
    =\min_{\partial_{t}\mu+\nabla\cdot(v\mu)=0}\mathcal{A}(\mu,v).
  \]

  \textbf{联系与意义。} 最优流就是位移插值路径
  $\mu_{t}=((1-t)\Id+tT)_{\#}\mu_{0}$；公式给出 $\Wtwo$ 的 PDE 表示、数值算法入口与 Otto 几何解释。

  \TLtakeaway{Benamou–Brenier：$W_2$ = 最小动能流；联结 PDE、数值计算与 Riemann 几何。}
\end{frame}

% ── §4 小结 ──────────────────────────────────────────────────────────────────
\begin{frame}{§4 小结：位移插值的核心结论}
  \begin{enumerate}
    \item \textbf{作用量框架：} $A(\gamma)=\int_{0}^{1}L(\gamma_{t},\dot{\gamma}_{t},t)\dd t$；
    严格凸 $L$（$p>1$）的最小化曲线为常速测地线（Jensen），关联成本 $c=d^{p}$。
    \item \textbf{构造：} 对最优耦合 $\pi$ 作测地线提升得 $\Pi$，$\mu_{t}=(e_{t})_{\#}\Pi$；
    $\R^{n}$ 中：$\mu_{t}=((1-t)\Id+tT)_{\#}\mu_{0}$。
    \item \textbf{定理 7.21：} 位移插值 $\iff$ 三时刻等号关系 $\iff$ 作用量最小化路径。
    \item \textbf{推论 7.22：} 位移插值 = $\Wp{p}$ 常速测地线：$\Wp{p}(\mu_{s},\mu_{t})=|t-s|\Wp{p}(\mu_{0},\mu_{1})$。
    \item \textbf{核心口号：} "概率测度空间的测地线 = 测地线的分布。"
  \end{enumerate}

  \textbf{下节预告：} §5 收缩性与位移凸性——位移插值的正则性和曲率下界。

  \TLtakeaway{位移插值统一了最优传输问题与 Wasserstein 空间的测地线几何，是整个 Part I 的几何核心。}
\end{frame}

% ── 习题 4 · Exercises（一）─────────────────────────────────────────────────
\begin{frame}{习题 4 · Exercises（一）}
  \diffI\ \textbf{习题 4.1.} 设 $\mu_{0}=\delta_{0}$，$\mu_{1}=\delta_{1}$，$\X=\R$，$p=2$。
  计算线性插值 $\mu_{t}^{\mathrm{lin}}=(1-t)\delta_{0}+t\delta_{1}$ 和位移插值 $\mu_{t}=\delta_{t}$ 之间的
  Wasserstein 距离 $\Wtwo(\mu_{t}^{\mathrm{lin}},\mu_{t})$。

  \hint{位移插值为 $\mu_{t}=\delta_{t}$；对于线性插值，利用 $\Wtwo^{2}((1-t)\delta_{0}+t\delta_{1},\,\delta_{t})$ 用最优匹配计算。}
  % SOLUTION: The only coupling of ((1-t)delta_0 + t*delta_1, delta_t) sends mass (1-t) from 0 to t and mass t from 1 to t. W_2^2 = (1-t)*t^2 + t*(1-t)^2 = t(1-t)(t+(1-t)) = t(1-t). So W_2 = sqrt(t(1-t)), maximized at t=1/2 with value 1/2.

  \vspace{6pt}
  \diffII\ \textbf{习题 4.2.} 设 $\mu_{0}=\mathcal{N}(0,1)$，$\mu_{1}=\mathcal{N}(0,4)$，$c(x,y)=(x-y)^{2}$。
  \begin{itemize}
    \item[(a)] 求最优传输映射 $T$。
    \item[(b)] 求位移插值 $\mu_{t}$ 的均值和方差。
    \item[(c)] 计算 $\Wtwo(\mu_{0},\mu_{1})$。
  \end{itemize}
  \hint{(a) $T(x)=(\sigma_{1}/\sigma_{0})x=2x$；(b) $m_{t}=0$，$\sigma_{t}=(1-t)\cdot1+t\cdot2=1+t$；(c) 由单调映射 $\Wtwo^{2}=\E[(X-T(X))^{2}]=\E[(X-2X)^{2}]=\E[X^{2}]=\sigma_{0}^{2}=1$，即 $\Wtwo=|\sigma_{1}-\sigma_{0}|=1$。}
  % SOLUTION: T(x)=2x; mu_t=N(0,(1+t)^2); W_2=|sigma_1-sigma_0|=|2-1|=1.

  \TLtakeaway{习题 4.1 揭示线性与位移插值的几何差距；习题 4.2 通过具体数字掌握一维高斯位移插值的全过程。}
\end{frame}

% ── 习题 4 · Exercises（二）────────────────────────────────────────────────
\begin{frame}{习题 4 · Exercises（二）}
  \diffIII\ \textbf{习题 4.3.} 设 $\mu_{0}$，$\mu_{1}$ 是 $\R^{n}$ 上均值为零、协方差矩阵分别为
  $\Sigma_{0}$，$\Sigma_{1}$ 的高斯分布。
  \begin{itemize}
    \item[(a)] 利用 Brenier 定理证明最优传输映射为线性映射 $T(x)=Ax$，其中
      \[
        A=\Sigma_{0}^{-1/2}\!\left(\Sigma_{0}^{1/2}\Sigma_{1}\Sigma_{0}^{1/2}\right)^{1/2}\!\Sigma_{0}^{-1/2}.
      \]
    \item[(b)] 写出位移插值 $\mu_{t}$ 的协方差矩阵 $\Sigma_{t}$。
    \item[(c)] 写出 $\Wtwo(\mu_{0},\mu_{1})^{2}$ 的迹公式。
  \end{itemize}
  \hint{(a) Brenier 定理：最优映射是凸函数的梯度 $T=\nabla\varphi$；高斯情形 $\varphi$ 是二次型，故 $T$ 线性且对称正定。(b) $\Sigma_{t}=((1-t)I+tA)\Sigma_{0}((1-t)I+tA)^{\top}$。(c) $\Wtwo^{2}=\mathrm{tr}(\Sigma_{0})+\mathrm{tr}(\Sigma_{1})-2\,\mathrm{tr}\bigl((\Sigma_{0}^{1/2}\Sigma_{1}\Sigma_{0}^{1/2})^{1/2}\bigr)$。}
  % SOLUTION: T(x)=Ax with A as given; mu_t=N(0,Sigma_t) with Sigma_t=((1-t)I+tA)Sigma_0((1-t)I+tA)^T; W_2^2=tr(Sigma_0)+tr(Sigma_1)-2tr((Sigma_0^{1/2}Sigma_1 Sigma_0^{1/2})^{1/2}).

  \TLtakeaway{高斯位移插值的矩阵公式是习题 4.2 的高维推广，训练对 Brenier 定理和矩阵平方根的运用。}
\end{frame}
% §4  Ch.7 位移插值
% =============================================================================
%  §5  Ch.8 Monge–Mather 缩短原理
%  来源：Villani, Optimal Transport: Old and New, Chapter 8 (pp. 164–198)
%  主要定理：定理 8.1 / 推论 8.2（Mather 缩短引理），定理 8.5，定理 8.7
% =============================================================================
\TLsection{5 \ Monge--Mather 缩短原理}

% ── 动机（一）：最优传输轨迹不相交 ──────────────────────────────────────────
\begin{frame}{动机：最优传输轨迹不相交（Monge 观察）}
  \textbf{核心问题。} 如果两条最优运输路径相交，会发生什么？

  设 $c(x,y)=|x-y|^{2}$，$x_{1},x_{2}$ 是出发点，$y_{1},y_{2}$ 是目标点。

  \deflab{Monge 观察（Villani 第 8 章）。}
  \[
    |x_{1}-y_{1}|^{2}+|x_{2}-y_{2}|^{2} \;\le\; |x_{1}-y_{2}|^{2}+|x_{2}-y_{1}|^{2}.
  \]
  等价条件（展开后）：
  \[
    \langle x_{1}-x_{2},\,y_{1}-y_{2}\rangle \;\ge\; 0.
  \]

  \textbf{含义：} 相交赋值（$x_{1}\to y_{2}$，$x_{2}\to y_{1}$）的成本$\ge$不相交赋值（$x_{1}\to y_{1}$，$x_{2}\to y_{2}$）的成本。

  \textbf{推论：} 若存在更优的"不相交"配对，相交路径不可能是最优的——可以"解开"来降低总成本。

  \TLtakeaway{Monge 观察：对二次成本，最优赋值必为单调不相交赋值。交叉路径可通过"解开"严格降低总成本。}
\end{frame}

% ── 动机（二）：图示 ──────────────────────────────────────────────────────────
\begin{frame}{图示：相交 vs 不相交赋值（据 Villani 图 8.1 重绘）}
  \begin{center}
  \begin{tikzpicture}[scale=1.0, >={Stealth[length=5pt]}]
    \begin{scope}[shift={(-3.2,0)}]
      \node[font=\small,TLinkSoft] at (0,1.6) {\textbf{相交赋值}};
      \fill[TLprimary]  (-0.7, 0.9) circle (2.5pt) node[left,font=\scriptsize]{$x_1$};
      \fill[TLprimary]  ( 0.7, 0.9) circle (2.5pt) node[right,font=\scriptsize]{$x_2$};
      \fill[TLaccent]   (-0.7,-0.9) circle (2.5pt) node[left,font=\scriptsize]{$y_1$};
      \fill[TLaccent]   ( 0.7,-0.9) circle (2.5pt) node[right,font=\scriptsize]{$y_2$};
      \draw[->,TLprimary,thick] (-0.7, 0.9) -- ( 0.7,-0.9);
      \draw[->,TLaccent, thick] ( 0.7, 0.9) -- (-0.7,-0.9);
      \fill[TLneg] (0,0) circle (2.5pt);
      \node[TLneg,font=\scriptsize,right=1pt] at (0,0) {交叉};
    \end{scope}
    \draw[->,TLinkSoft,very thick] (-0.8,0) -- (0.8,0)
      node[midway,above,font=\scriptsize]{解开};
    \begin{scope}[shift={(3.2,0)}]
      \node[font=\small,TLinkSoft] at (0,1.6) {\textbf{不相交（最优）}};
      \fill[TLprimary]  (-0.7, 0.9) circle (2.5pt) node[left,font=\scriptsize]{$x_1$};
      \fill[TLprimary]  ( 0.7, 0.9) circle (2.5pt) node[right,font=\scriptsize]{$x_2$};
      \fill[TLaccent]   (-0.7,-0.9) circle (2.5pt) node[left,font=\scriptsize]{$y_1$};
      \fill[TLaccent]   ( 0.7,-0.9) circle (2.5pt) node[right,font=\scriptsize]{$y_2$};
      \draw[->,TLprimary,thick] (-0.7, 0.9) -- (-0.7,-0.9);
      \draw[->,TLaccent, thick] ( 0.7, 0.9) -- ( 0.7,-0.9);
      \node[TLpos,font=\scriptsize] at (0,-1.3) {成本更低};
    \end{scope}
  \end{tikzpicture}
  \end{center}
  {\scriptsize 据 Villani 重绘。}

  \textbf{数值验证（$x_1=0,\,x_2=1,\,y_1=0,\,y_2=1$）：}
  相交成本 $=|0-1|^2+|1-0|^2=2$；不相交成本 $=|0-0|^2+|1-1|^2=0$。节省了 $2$。

  \TLtakeaway{将相交赋值"解开"得到不相交赋值，总成本严格下降。这是第 8 章所有正则性结论的几何核心。}
\end{frame}

% ── 关键恒等式 (8.2) ──────────────────────────────────────────────────────────
\begin{frame}{关键恒等式：中间时刻控制全程（Villani 式 8.2）}
  \textbf{设置。} 令 $\gamma_{i}(t)=(1-t)x_{i}+ty_{i}$，$i=1,2$，$t\in[0,1]$。

  \deflab{恒等式 8.2（Villani）。}
  \begin{align*}
    &|\gamma_{1}(t)-\gamma_{2}(t)|^{2} \\
    &= (1-t)^{2}|x_{1}-x_{2}|^{2}+t^{2}|y_{1}-y_{2}|^{2} \\
    &\quad+t(1-t)\bigl[|x_{1}-y_{2}|^{2}+|x_{2}-y_{1}|^{2}-|x_{1}-y_{1}|^{2}-|x_{2}-y_{2}|^{2}\bigr].
  \end{align*}

  \textbf{若 $(x_1,y_1),(x_2,y_2)$ 满足最优性条件（括号内 $\ge 0$），则对任意 $t_0\in(0,1)$：}
  \[
    \sup_{0\le t\le 1}|\gamma_{1}(t)-\gamma_{2}(t)|
    \;\le\; \max\!\left(\tfrac{1}{t_{0}},\tfrac{1}{1-t_{0}}\right)
    \cdot|\gamma_{1}(t_{0})-\gamma_{2}(t_{0})|.
  \]

  \textbf{定性结论：} 若两最优轨迹在某中间时刻 $t_0$ 重合，则在整个 $[0,1]$ 上重合。

  \TLtakeaway{式 8.2 是 Monge 无交叉原理的精确定量化：中间时刻 $t_0$ 处的距离控制全程一致距离，乘以 $1/\min(t_0,1-t_0)$。这是第 8 章所有估计的代数基础。}
\end{frame}

% ── 定理 8.1（Mather 缩短引理）─────────────────────────────────────────────
\begin{frame}{定理 8.1：Mather 缩短引理（陈述）}
  \deflab{定理 8.1（Mather 缩短引理，Villani Thm 8.1）。}

  设 $M$ 为光滑 Riemann 流形，$d$ 为测地距离，$c(x,y)$ 由 Lagrange 量 $L(x,v,t)$ 定义。
  设 $x_{1},x_{2},y_{1},y_{2}\in M$ 满足 $c(x_{1},y_{1})+c(x_{2},y_{2})\le c(x_{1},y_{2})+c(x_{2},y_{1})$，
  $\gamma_{1},\gamma_{2}$ 为对应的作用量最小化曲线。假设：

  \begin{itemize}\small
    \item[(i)] 最小化曲线满足 Lipschitz 流（定义 7.6(d)）；
    \item[(ii)] $L$ 关于 $(x,v)$ 在 $\gamma_1,\gamma_2$ 的邻域中是 $C^{1,\alpha}$，$\alpha\in(0,1]$；
    \item[(iii)] $L$ 关于速度 $v$ 是 $(2+\kappa)$-凸的（严格凸的量化版本）。
  \end{itemize}

  则对任意 $t_{0}\in(0,1)$，存在常数 $C_{t_{0}}$ 和指数 $\beta=\beta(\alpha,\kappa)>0$ 使得
  \[
    \sup_{0\le t\le 1}d\!\bigl(\gamma_{1}(t),\gamma_{2}(t)\bigr)
    \;\le\; C_{t_{0}}\;d\!\bigl(\gamma_{1}(t_{0}),\gamma_{2}(t_{0})\bigr)^{\beta}.
  \]

  \TLtakeaway{定理 8.1（一般形式）：严格凸 Lagrange 量 $\Rightarrow$ 中间时刻距离 $\to$ 全程距离，指数 $\beta>0$。假设 (iii) 的 $(2+\kappa)$-凸性是量化的关键。}
\end{frame}

% ── 推论 8.2 ─────────────────────────────────────────────────────────────────
\begin{frame}{推论 8.2：$C^{2}$ Lagrange 量的缩短引理}
  \deflab{推论 8.2（Villani Cor 8.2）。}

  设 $L=L(x,v,t)$ 为 $C^{2}$ Lagrange 量，$\nabla_{v}^{2}L>0$（严格凸），$c$ 为对应成本。
  对任意紧集 $K\subset M$，存在 $C_{K}>0$：只要 $x_{1},y_{1},x_{2},y_{2}\in K$ 满足最优性条件，
  $\gamma_{1},\gamma_{2}$ 为对应的作用量最小化曲线，则对任意 $t_{0}\in(0,1)$，
  \[
    \sup_{0\le t\le 1}d\!\bigl(\gamma_{1}(t),\gamma_{2}(t)\bigr)
    \;\le\; \frac{C_{K}}{\min(t_{0},1-t_{0})}\;d\!\bigl(\gamma_{1}(t_{0}),\gamma_{2}(t_{0})\bigr).
  \]

  \textbf{最重要情形：} $c(x,y)=d(x,y)^{2}$（$L=|v|^{2}$，$\nabla_{v}^{2}L=2I>0$），推论 8.2 完整适用，指数 $\beta=1$。

  \textbf{例 8.4（Villani）。} $c=d^{1+\alpha}$（$0<\alpha<1$）：$L$ 非 $C^2$，但幂函数齐次性仍给出 $\beta=1$（见 Villani 附录）。

  % SOLUTION: 定理 8.1 完整证明含快捷路径构造和 Hölder 估计，见 Villani 第 172--180 页。留作 Part 6 备用幻灯片。

  \TLtakeaway{推论 8.2（最常用形式）：$C^2$、严格凸 Lagrange 量（含 $c=d^2$），$\sup d \le \frac{C_K}{\min(t_0,1-t_0)}\cdot d(t_0)$，指数 $\beta=1$。}
\end{frame}

% ── 定理 8.1 证明思路 ────────────────────────────────────────────────────────
\begin{frame}{定理 8.1 证明思路（概要）}
  \textbf{核心想法：} 若 $\gamma_{1},\gamma_{2}$ 在 $t_{0}$ 时刻相交于 $m_{0}$，且速度分别为 $v_{1},v_{2}$，就在短时间窗内构造"快捷路径"互换连接。

  \textbf{步骤一（局部化）。} 将流形问题化归为局部 Euclid 计算：紧性覆盖 $\gamma_{1}\cup\gamma_{2}$ 邻域，在小球中引入坐标。

  \deflab{补一个事实：$(2+\kappa)$-凸性。}
  这是 $L$ 关于速度的量化严格凸性：存在 $c_{\mathrm{gap}}>0$，使得
  \[
    L\!\left(x,\frac{v_{1}+v_{2}}{2},t\right)
    \le \frac{1}{2}L(x,v_{1},t)+\frac{1}{2}L(x,v_{2},t)
    -c_{\mathrm{gap}}\abs{v_{1}-v_{2}}^{2+\kappa}.
  \]
  这就是 Villani 使用的 midpoint-gap 形式。

  \TLtakeaway{定理 8.1 的局部证明只用两个工具：把曲线放进 Euclid 坐标，以及把速度严格凸性读成 midpoint-gap。}
\end{frame}

\begin{frame}{定理 8.1 证明思路（快捷路径主项）}
  \textbf{步骤二（快捷路径）。}
  在 $[t_{0}-\tau,t_{0}+\tau]$ 内，将两段曲线剪开并互换连接；两条新路径都近似以平均速度 $(v_{1}+v_{2})/2$ 穿过 $m_{0}$。主阶作用量改变为
  \[
    \Delta A \approx
    2\tau\!\left[2L\!\left(m_{0},\frac{v_{1}+v_{2}}{2},t_{0}\right)
    -L(m_{0},v_{1},t_{0})-L(m_{0},v_{2},t_{0})\right].
  \]

  \TLtakeaway{快捷路径为什么会降作用量：两条原速度 $v_{1},v_{2}$ 被两个平均速度替代，主项正是 midpoint-gap 的负数。}
\end{frame}

\begin{frame}{定理 8.1 证明思路（凸性矛盾）}
  \textbf{步骤三（严格下降）。}
  把上一帧的主项代入 midpoint-gap：若 $v_{1}\ne v_{2}$，则
  \[
    \Delta A
    \lesssim -4\tau c_{\mathrm{gap}}\abs{v_{1}-v_{2}}^{2+\kappa}<0.
  \]
  因而互换后的两条路径总作用量更小，即
  $c(x_{1},y_{2})+c(x_{2},y_{1})<c(x_{1},y_{1})+c(x_{2},y_{2})$，这与定理 8.1 的最优性假设矛盾。

  \textbf{步骤四（从相同速度到相同曲线）。}
  唯一剩下的可能是 $v_{1}=v_{2}$。此时两条曲线在 $t_{0}$ 处有同一位置 $m_{0}$ 和同一速度；最小化曲线满足 Lipschitz 流，故 Cauchy--Lipschitz 唯一性推出 $\gamma_{1}=\gamma_{2}$。

  \deflab{本节诚实省略的技术估计。}
  省略的是 Villani 第 172--180 页中，从快捷路径误差界 (8.21) 与凸性下界 (8.22) 推出速度 Hölder 控制 (8.8) 的计算；主线只需要知道该估计把"相交矛盾"量化成 (8.4)。
  % SOLUTION: 完整速度估计 (8.8) 的推导留作 Part 6 备用幻灯片。

  \TLtakeaway{定性结论来自严格凸性矛盾；定量结论来自 (8.21)--(8.22) 的误差比较，最终得到速度差被位置差的幂控制。}
\end{frame}

% ── 定理 8.5：陈述 ────────────────────────────────────────────────────────────
\begin{frame}{定理 8.5：中间时刻传输映射是 Lipschitz（陈述）}
  \deflab{定理 8.5（Villani Thm 8.5）。}

  设 $M$ 为 Riemann 流形，$c$ 满足推论 8.2 假设，$K\subset M$ 紧，$\Pi$ 为支撑在 $K$ 内的动力学最优传输。

  则 $\Pi$-几乎处处，对任意 $\gamma,\tilde{\gamma}\in\mathrm{supp}(\Pi)$，
  \[
    \sup_{0\le t\le 1}d\!\bigl(\gamma(t),\tilde{\gamma}(t)\bigr)
    \;\le\; \frac{C_{K}}{\min(t_{0},1-t_{0})}\,d\!\bigl(\gamma(t_{0}),\tilde{\gamma}(t_{0})\bigr).
  \]

  \textbf{关键推论（传输映射）。} 若 $(\mu_{t})$ 为紧支撑测度间的位移插值，$t_{0}\in(0,1)$，
  则映射 $T_{t_{0}\to t}\colon\gamma(t_{0})\mapsto\gamma(t)$ 在 $\mu_{t_{0}}$-a.s. 有定义，
  在定义域上是 Lipschitz 连续的（Lipschitz 常数 $C_{K}/\min(t_{0},1-t_{0})$），
  且是 $\mu_{t_{0}}$ 到 $\mu_{t}$ 的 Monge 问题的唯一解。

  \TLtakeaway{定理 8.5：中间时刻 $t_0\in(0,1)$ 的传输映射自动 Lipschitz，即使端点传输映射可能不正则——这是缩短引理的第一个重大应用。}
\end{frame}

% ── 定理 8.5：证明思路 ──────────────────────────────────────────────────────
\begin{frame}{定理 8.5 证明思路}
  \textbf{步骤一（最优性传递）。}
  $\pi=(e_{0},e_{1})_{\#}\Pi$ 为最优耦合，故 $c$-循环单调，因此 $\Pi\otimes\Pi$-a.s.：
  \[
    c(\gamma(0),\gamma(1))+c(\tilde{\gamma}(0),\tilde{\gamma}(1))
    \;\le\; c(\gamma(0),\tilde{\gamma}(1))+c(\tilde{\gamma}(0),\gamma(1)).
  \]

  \textbf{步骤二（推论 8.2）。} 上式即推论 8.2 前提，直接得到 Lipschitz 估计
  $\sup_{t}\,d(\gamma(t),\tilde{\gamma}(t))\le\frac{C_{K}}{\min(t_{0},1-t_{0})}\,d(\gamma(t_{0}),\tilde{\gamma}(t_{0}))$。

  \deflab{补一个事实：良定义就是单值。}
  若同一个中间点可由两条支撑曲线到达，即 $\gamma(t_{0})=\tilde{\gamma}(t_{0})$，则右端为 $0$。

  \textbf{步骤三（良定义与唯一性）。}
  缩短估计强制 $\sup_{t}d(\gamma(t),\tilde{\gamma}(t))=0$，所以 $\gamma=\tilde{\gamma}$。
  因而 $T_{t_{0}\to t}\colon\gamma(t_{0})\mapsto\gamma(t)$ 不依赖曲线选择，是单值良定义的；同一估计给出 Lipschitz 常数。最优唯一性由 §2 定理 5.30（Monge 可解性）给出。

  \TLtakeaway{$c$-循环单调 $\Rightarrow$ 推论 8.2 $\Rightarrow$ 缩短估计；中间点若相同，则整条曲线相同，所以 $T_{t_{0}\to t}$ 是单值良定义的。}
\end{frame}

\begin{frame}{例 8.6：中间映射为何显然 Lipschitz}
  \textbf{例 8.6（Villani）。} $c=|\cdot|^{2}$，$\mu_{0}=\delta_{0}$，$\mu=\law(X)$：
  随机测地线为 $\gamma(t)=t\gamma(1)$，所以
  \[
    \mu_{t}=\law(tX),
    \qquad
    T_{t_{0}\to t}=\frac{t}{t_{0}}\Id .
  \]
  该映射的 Lipschitz 常数是 $t/t_{0}$；这正是定理 8.5 在二次成本最简单情形下的显式模型。

  \TLtakeaway{例 8.6 把抽象结论化成公式：中间位置 $t_{0}X$ 唯一决定任意时刻 $tX$，映射就是线性缩放。}
\end{frame}

% ── 定理 8.7：陈述 ────────────────────────────────────────────────────────────
\begin{frame}{定理 8.7：位移插值保持绝对连续性（陈述）}
  \deflab{定理 8.7（绝对连续性，Villani Thm 8.7）。}

  设 $M$ 为 Riemann 流形，$L\in C^{2}$，$\nabla_{v}^{2}L>0$，$c$ 为对应成本。
  设 $\mu_{0},\mu_{1}\in\Pp(M)$，$C(\mu_{0},\mu_{1})<+\infty$，$(\mu_{t})_{0\le t\le 1}$ 为位移插值。

  \begin{center}
  \textbf{若 $\mu_{0}$ 或 $\mu_{1}$ 中至少一个绝对连续于 $M$ 上的体积元，}\\[4pt]
  \textbf{则 $\mu_{t}\in\Pac(M)$ 对所有 $t\in(0,1)$ 成立。}
  \end{center}
  {\small （此处 $\Pac$ 为绝对连续概率测度，定义见 §0。）}

  \textbf{直观。} 在中间时刻 $t_0$，传输映射 $T_{t_0\to 1}$ 是 Lipschitz 的（定理 8.5）。
  Lipschitz 映射的像保持零测集的零测性（Lusin N 条件），从而绝对连续性反向传递：
  $\mu_1\in\Pac$，$T_{t_0\to 1}$ Lipschitz，故 $\mu_{t_0}\in\Pac$。

  \TLtakeaway{定理 8.7：端点之一绝对连续 $\Rightarrow$ 中间时刻的插值测度自动绝对连续。内部正则性从最优传输的组合结构中"免费涌现"。}
\end{frame}

% ── 定理 8.7：证明思路 ──────────────────────────────────────────────────────
\begin{frame}{定理 8.7 证明（紧支撑情形）}
  \textbf{假设 $\mu_{1}\in\Pac$，$\mu_{0},\mu_{1}$ 紧支撑，证明 $\mu_{t_{0}}\in\Pac$，$t_{0}\in(0,1)$。}

  \textbf{步骤一。} 由定理 8.5（Lipschitz 性来自推论 8.2），存在 Lipschitz Monge 映射 $T\colon\mu_{t_{0}}\to\mu_{1}$。

  \deflab{补一个事实：Lusin (N) 性质 / 面积公式。}
  若 $T\colon\R^{n}\to\R^{n}$ 是 $L_{T}$-Lipschitz，则
  $\mathrm{vol}(T(N))\le L_{T}^{n}\mathrm{vol}(N)$；零测集仍送到零测集。流形上用局部坐标。

  \textbf{步骤二（零测集）。} 设 $N$ 为零体积集，并取 $L_{T}=C_{K}/\min(t_{0},1-t_{0})$。于是
  \[
    \mathrm{vol}(N)=0
    \;\Longrightarrow\;
    \mathrm{vol}(T(N))\le L_{T}^{n}\mathrm{vol}(N)=0
    \;\Longrightarrow\;
    \mu_{1}[T(N)]=0.
  \]
  最后一箭头用 $\mu_{1}\in\Pac$，见 §0。

  \TLtakeaway{Lusin (N) / 面积公式在这里只做一件事：$N$ 零体积 $\Rightarrow$ $T(N)$ 零体积 $\Rightarrow$ $\mu_{1}[T(N)]=0$。}
\end{frame}

\begin{frame}{定理 8.7 证明（紧支撑收尾）}
  \textbf{步骤三（绝对连续性）。} 由 $N\subset T^{-1}(T(N))$：
  \[
    \mu_{t_{0}}[N]\;\le\;\mu_{t_{0}}\!\left[T^{-1}(T(N))\right]
    =\pf{T}{\mu_{t_{0}}}[T(N)]=\mu_{1}[T(N)]=0,
  \]
  故 $\mu_{t_{0}}\in\Pac$。

  \textbf{可测性小修。} Villani 指出 $T(N)$ 未必是 Borel 集；严格写法是用一个零体积 Borel 集包含 $T(N)$，上式按该 Borel 包含集解释即可。

  \TLtakeaway{$T$ Lipschitz $+$ $\mu_{1}\in\Pac$ $\Rightarrow$ $T(N)$ 零测 $\Rightarrow$ $\mu_{t_{0}}[N]=0$：绝对连续性沿 Lipschitz 传输反向传递。}
\end{frame}

% ── 定理 8.7：证明续（非紧情形）──────────────────────────────────────────────
\begin{frame}{定理 8.7 证明续（非紧支撑情形）}
  \textbf{步骤四（非紧推广）。} 若 $\mu_{t_0}\notin\Pac$，取体积零但 $\mu_{t_0}$-正测度的集 $Z_{t_0}$；
  对 $\Pi$ 限制到紧子集 $K\subset\{\gamma\colon\gamma(t_0)\in Z_{t_0}\}$，化为上述紧支撑情形，得矛盾。

  \textbf{小结。} 三步骨架：
  \begin{enumerate}\small
    \item 定理 8.5 $\to$ $T$ Lipschitz（推论 8.2 缩短引理）。
    \item Lipschitz 面积公式 $\to$ $T(N)$ 零体积。
    \item $\mu_1\in\Pac$ $\to$ $\mu_1[T(N)]=0$ $\to$ $\mu_{t_0}[N]=0$。
  \end{enumerate}

  \TLtakeaway{绝对连续性沿 Lipschitz 传输的逆向传递：$\mu_1\in\Pac$（见 §0）$+$ $T$ Lipschitz $\Rightarrow$ $\mu_{t_0}\in\Pac$。非紧情形通过紧子集的穷举化归完成。}
\end{frame}

% ── 为什么定理 8.7 重要 ────────────────────────────────────────────────────
\begin{frame}{为什么定理 8.7 重要？内部正则性的来源}
  \textbf{看似悖论。} 端点 $\mu_0$ 可以非常奇异（原子测度、奇异测度），
  但只要 $\mu_1\in\Pac$，中间时刻 $\mu_t$ 就自动绝对连续。

  \begin{center}
  \begin{tikzpicture}[scale=0.9, >={Stealth[length=4pt]}]
    \node[draw=TLprimary,rounded corners=3pt,inner sep=5pt,font=\small,
          fill=TLprimaryTint] (L) at (0,0) {$\mu_0$（可奇异）};
    \node[draw=TLaccent,rounded corners=3pt,inner sep=5pt,font=\small,
          fill=TLaccentLight] (R) at (7.5,0) {$\mu_1\in\Pac$};
    \node[draw=TLpos,rounded corners=3pt,inner sep=5pt,font=\small,
          fill=TLpos!10] (M) at (3.75,0) {$\mu_t\in\Pac$，$\forall\,t\in(0,1)$};
    \draw[->,TLprimary,thick] (L) -- (M);
    \draw[->,TLaccent,thick] (R) -- (M);
  \end{tikzpicture}
  \end{center}

  \textbf{物理直觉。} 质量在中间时刻的位置由 Lipschitz 映射决定，不能聚焦到零测集。

  \textbf{整体逻辑链：}
  \[
    \underbrace{c\text{-循环单调}}_{\text{最优性}}
    \;\overset{\text{推论 8.2}}{\Rightarrow}\;
    \underbrace{T_{t_0\to t}\text{ Lipschitz}}_{\text{定理 8.5}}
    \;\overset{\text{Lusin N}}{\Rightarrow}\;
    \underbrace{\mu_t\in\Pac}_{\text{定理 8.7}}
  \]

  \textbf{前瞻。} 第 9–10 章以此内部绝对连续性为基础研究 Wasserstein 几何的曲率下界。

  \TLtakeaway{"内部正则性免费涌现"：最优传输的组合最优性（无交叉）通过量化（缩短引理）产生分析正则性（$\mu_t\in\Pac$）——这是最优传输几何的核心奇迹之一。}
\end{frame}

% ── 习题 5 · Exercises（一）────────────────────────────────────────────────
\begin{frame}{习题 5 · Exercises（一）}
  \diffI\ \textbf{习题 5.1.} 设 $x_1=0$，$x_2=2$，对 $c(x,y)=(x-y)^2$，
  比较赋值 $(x_1\to 0,\,x_2\to 2)$ 与 $(x_1\to 2,\,x_2\to 0)$ 的总成本，
  验证不相交赋值严格更优。

  \hint{不相交成本：$(0-0)^2+(2-2)^2=0$；相交成本：$(0-2)^2+(2-0)^2=8$。差值 $8=2\langle x_1-x_2,y_2-y_1\rangle$，与 $\langle x_1-x_2,y_1-y_2\rangle\ge 0$ 对应。}
  % SOLUTION: 不相交成本=0，相交成本=8，节省8。一般形式：展开差值得2(x1-x2)(y2-y1)>=0 当 x1<x2,y1<y2 时。

  \vspace{6pt}
  \diffII\ \textbf{习题 5.2.} 设 $(x_1,y_1),(x_2,y_2)\in\R\times\R$ 是某最优耦合
  $\pi$（成本 $c(x,y)=(x-y)^2$）支撑集中的两点。
  证明 $x_1<x_2\;\Rightarrow\;y_1\le y_2$（单调性）。

  \hint{利用 $\pi$ 的 $c$-循环单调性：$c(x_1,y_1)+c(x_2,y_2)\le c(x_1,y_2)+c(x_2,y_1)$。展开后等价于 $(x_1-x_2)(y_1-y_2)\ge 0$。$x_1<x_2$ 时 $x_1-x_2<0$，故 $y_1-y_2\le 0$。}
  % SOLUTION: c-循环单调给出(x1-y1)^2+(x2-y2)^2<=(x1-y2)^2+(x2-y1)^2。展开：(x1-x2)(y1-y2)>=0。x1<x2 时 y1<=y2。

  \TLtakeaway{习题 5.1 数值验证无交叉原理；习题 5.2 代数证明最优耦合支撑集的单调性——二者共同将 Monge 观察从直觉化为精确命题。}
\end{frame}

% ── 习题 5 · Exercises（二）────────────────────────────────────────────────
\begin{frame}{习题 5 · Exercises（二）}
  \diffIII\ \textbf{习题 5.3.} 设 $\mu_0\in\Pac(\R^n)$（绝对连续），$\mu_1$ 为任意概率测度，
  $c(x,y)=|x-y|^2$，$C(\mu_0,\mu_1)<+\infty$。
  利用定理 8.7 证明位移插值 $\mu_t$ 在 $t\in(0,1)$ 上绝对连续。

  \begin{itemize}
    \item[(a)] 识别定理 8.7 适用所需的假设，逐条验证。
    \item[(b)] $\mu_0\in\Pac$ 与 $\mu_1\in\Pac$ 在定理中的角色是否对称？
    \item[(c)] 若改为 $\mu_0$ 是原子测度（非绝对连续），$\mu_1\in\Pac$，结论是否仍成立？
  \end{itemize}

  \hint{(a) 检查：$L=|v|^2$ 是 $C^2$ 且 $\nabla_v^2L=2I>0$；成本有限；$\mu_0\in\Pac$ 满足定理假设之一。(b) 定理 8.7 对 $\mu_0$ 或 $\mu_1$ 绝对连续均成立，故对称。(c) 仍成立，用 $\mu_1\in\Pac$ 的那个方向。}
  % SOLUTION: (a) L=|v|^2满足C^2且Hessian严格正定；C(mu_0,mu_1)<+inf已给；mu_0∈Pac是定理假设之一。定理8.7直接适用，mu_t∈Pac for t∈(0,1)。(b) 定理8.7对"mu_0或mu_1中至少一个绝对连续"均成立，完全对称。(c) 仍成立：用mu_1∈Pac分支，定理8.7的证明对mu_1绝对连续同样有效。

  \TLtakeaway{习题 5.3 训练定理 8.7 假设的精确核查：$L\in C^2$，$\nabla_v^2L>0$，成本有限，端点之一绝对连续——四条缺一不可。}
\end{frame}

% ── §5 小结 ──────────────────────────────────────────────────────────────────
\begin{frame}{§5 小结：无交叉 $\to$ 缩短 $\to$ 内部正则性}
  \textbf{核心思路链：}
  \[
    \text{无交叉}
    \overset{\text{定量化}}{\Rightarrow}
    \text{缩短引理}
    \overset{c\text{-单调}}{\Rightarrow}
    T_{t_0\to t}\text{ Lipschitz}
    \overset{\text{Lusin N}}{\Rightarrow}
    \mu_t\in\Pac
  \]

  \textbf{三个主要结果：}
  \begin{enumerate}\small
    \item \textbf{推论 8.2（缩短引理）：}
      $\sup_t d\le\frac{C_K}{\min(t_0,1-t_0)}\cdot d(t_0)$（$C^2$、严格凸 $L$）。
    \item \textbf{定理 8.5（Lipschitz 传输）：}
      中间时刻 $t_0\in(0,1)$ 的传输映射 $T_{t_0\to t}$ Lipschitz，是唯一 Monge 解。
    \item \textbf{定理 8.7（绝对连续性）：}
      端点之一 $\in\Pac\;\Rightarrow\;\mu_t\in\Pac$ 对所有 $t\in(0,1)$ 成立。
  \end{enumerate}

  \textbf{下节预告。} 第 9–10 章以此内部绝对连续性为出发点，研究 Wasserstein 空间的曲率下界和位移凸性。

  \TLtakeaway{第 8 章精华：最优传输的组合最优性（无交叉）$\to$ 分析正则性（Lipschitz 传输、绝对连续插值）。联结了组合、分析与几何三个视角。}
\end{frame}
  % §5  Ch.8 Monge–Mather 缩短原理
% =============================================================================
%  §6  定性图景（总结）+ 术语表与记号 + 参考文献 + 习题解答（附录）
%  指挥官撰写。校验器锚点：Glossary / Notation / References / Bibliographical notes / Exercises。
% =============================================================================
\TLsection{6 \ 定性图景与总结}

% ── 五章合一（一）────────────────────────────────────────────────────────────
\begin{frame}{定性图景：五章合一（一）——一条主线}
  Villani 第 I 部分的"定性描述"沿一条主线展开：从\emph{解的存在}到\emph{解的几何与正则性}。
  \begin{center}
  \resizebox{\linewidth}{!}{%
  \begin{tikzpicture}[>={Stealth[length=6pt]},node distance=6mm,
      b/.style={draw=TLprimary,fill=TLprimaryTint,rounded corners=3pt,align=center,
        inner sep=5pt,text width=2.5cm,font=\scriptsize}]
    \node[b] (e) {\textbf{存在}\\§1 · Ch.4\\最优耦合存在};
    \node[b,right=of e] (d) {\textbf{刻画}\\§2 · Ch.5\\对偶 + 循环单调};
    \node[b,right=of d] (g) {\textbf{几何}\\§3 · Ch.6\\$\Wp{p}$ 距离};
    \node[b,right=of g] (k) {\textbf{动力学}\\§4 · Ch.7\\位移插值 = 测地线};
    \node[b,right=of k] (r) {\textbf{正则性}\\§5 · Ch.8\\缩短 $\Rightarrow$ 内部光滑};
    \draw[->,TLaccent,thick] (e)--(d);
    \draw[->,TLaccent,thick] (d)--(g);
    \draw[->,TLaccent,thick] (g)--(k);
    \draw[->,TLaccent,thick] (k)--(r);
  \end{tikzpicture}}
  \end{center}
  \begin{itemize}\small
    \item \textbf{存在}（§1）：紧性 + 下半连续 $\Rightarrow$ 最优耦合存在。
    \item \textbf{刻画}（§2）：最优 $\iff$ c-循环单调 $\iff$ 集中在 $\partial_c\psi$ $\iff$ 对偶等号。
    \item \textbf{几何}（§3）：$\Wp{p}$ 把最优成本升格为距离，$\Ppp{p}(\X)$ 是 Polish 空间。
    \item \textbf{动力学}（§4）：位移插值 = $\Wp{p}$ 测地线。\textbf{正则性}（§5）：轨迹不相交 $\Rightarrow$ 内部绝对连续。
  \end{itemize}
  \TLtakeaway{五章是一条因果链：存在 $\to$ 对偶刻画 $\to$ 度量几何 $\to$ 测地线动力学 $\to$ 内部正则性。}
\end{frame}

% ── 五章合一（二）────────────────────────────────────────────────────────────
\begin{frame}{定性图景：五章合一（二）——核心公式一览}
  \begin{center}\small
  \begin{tabular}{ll}
    \toprule
    \textbf{主题} & \textbf{核心命题 / 公式} \\
    \midrule
    Kantorovich 问题 & $C(\mu,\nu)=\inf_{\pi\in\coup{\mu}{\nu}}\int c\dd\pi$ \\
    对偶（Thm 5.10） & $C(\mu,\nu)=\sup_{\psi\text{ c-凸}}\bigl[\int\ctr{\psi}\dd\nu-\int\psi\dd\mu\bigr]$ \\
    最优性刻画 & $\pi$ 最优 $\iff$ $\supp\pi$ c-循环单调 $\iff$ $\pi$ 集中在 $\partial_c\psi$ \\
    Brenier（二次成本） & $T=\nabla\varphi$，$\varphi$ 凸（$\mu$ 绝对连续） \\
    Wasserstein 距离 & $\Wp{p}(\mu,\nu)=\bigl(\inf_\pi\int d^p\dd\pi\bigr)^{1/p}$ \\
    度量化（Thm 6.9） & $\Wp{p}(\mu_k,\mu)\to0 \iff \mu_k\wkto\mu$ 且 $p$ 阶矩收敛 \\
    位移插值（Thm 7.21） & $\mu_t=(e_t)_\#\Pi$，$\Wp{p}(\mu_s,\mu_t)=|t-s|\,\Wp{p}(\mu_0,\mu_1)$ \\
    缩短 / 正则（Thm 8.7） & $\mu_0$ 绝对连续 $\Rightarrow$ $\mu_t$ 绝对连续（$0<t<1$） \\
    \bottomrule
  \end{tabular}
  \end{center}
  \TLtakeaway{这张表是全讲义的"一页纸总览"：把它读懂，即掌握了最优传输定性理论的骨架。}
\end{frame}

% ── 术语表 · Glossary ─────────────────────────────────────────────────────────
\begin{frame}{术语表 · Glossary（一）}
  \begin{center}\small
  \begin{tabular}{ll}
    \toprule
    \textbf{术语（Term）} & \textbf{含义} \\
    \midrule
    耦合 / 转移计划 (coupling) & 以 $\mu,\nu$ 为边际的 $\X\times\Y$ 上概率测度 $\pi$ \\
    推前 (pushforward) $\pf{T}{\mu}$ & $T$ 把质量从 $\mu$ 搬运后得到的测度 \\
    Monge / Kantorovich 问题 & 确定性映射 / 可分裂耦合下的最优传输 \\
    c-循环单调 (c-cyclically monotone) & 无法通过有限循环重路由降低成本的支撑 \\
    c-凸 / c-变换 (c-convex / c-transform) & Legendre 变换的推广；$\ctr{\psi}(y)=\inf_x[\psi(x)+c(x,y)]$ \\
    Kantorovich 势 (potential) & 对偶问题的最优 c-凸函数 $\psi$ \\
    c-次微分 $\partial_c\psi$ & $\{(x,y):\ctr{\psi}(y)-\psi(x)=c(x,y)\}$，最优匹配关系 \\
    Brenier 映射 & 二次成本下的最优映射 $T=\nabla\varphi$（$\varphi$ 凸） \\
    \bottomrule
  \end{tabular}
  \end{center}
  \TLtakeaway{§1--§2 的语言：耦合、推前、对偶、c-凸与 c-次微分——理解它们即理解最优传输的"语法"。}
\end{frame}

\begin{frame}{术语表 · Glossary（二）}
  \begin{center}\small
  \begin{tabular}{ll}
    \toprule
    \textbf{术语（Term）} & \textbf{含义} \\
    \midrule
    Wasserstein 距离 $\Wp{p}$ & 以 $d^p$ 为成本的最优传输值的 $p$ 次根 \\
    Wasserstein 空间 $\Ppp{p}(\X)$ & $p$ 阶矩有限的概率测度，配以 $\Wp{p}$（Polish） \\
    Kantorovich–Rubinstein & $\Wp{1}=\sup_{\norm{\psi}_{\mathrm{Lip}}\le1}(\int\psi\dd\mu-\int\psi\dd\nu)$ \\
    胶着 (tight) & 质量不逃逸到无穷（Prokhorov 紧性判据） \\
    作用量 / Lagrange 量 (action) & $A(\gamma)=\int_0^1 L(\gamma_t,\dot\gamma_t,t)\dd t$ \\
    常速测地线 (constant-speed geodesic) & $d(\gamma_s,\gamma_t)=|t-s|\,d(\gamma_0,\gamma_1)$ \\
    位移插值 (displacement interpolation) & $\mu_t=(e_t)_\#\Pi$，$\Wp{p}$ 空间中的测地线 \\
    缩短原理 (shortening principle) & 最优轨迹不相交 $\Rightarrow$ 内部正则性 \\
    \bottomrule
  \end{tabular}
  \end{center}
  \TLtakeaway{§3--§5 的语言：距离、测地线、作用量与缩短——把静态的传输升格为动态的几何。}
\end{frame}

\begin{frame}{术语表 · Glossary（三）——§0 预备概念}
  \begin{center}\small
  \begin{tabular}{ll}
    \toprule
    \textbf{术语（Term）} & \textbf{含义} \\
    \midrule
    下半连续 (lsc) & $\liminf_{x\to x_0}f(x)\ge f(x_0)$；等价于所有下水平集 $\{f\le a\}$ 为闭集 \\
    上半连续 (usc) & $\limsup_{x\to x_0}f(x)\le f(x_0)$；lsc 的对偶概念 \\
    绝对连续 $\Pac$ & $\mu\ll\mathrm{vol}$（有 Lebesgue 密度）；无原子；Brenier 定理的前提 \\
    Hausdorff 测度 $\mathcal{H}^{k}$ & 度量空间上的 $k$ 维体积；$\mathcal{H}^{1}=$ 弧长，$\mathcal{H}^{n}\propto\mathrm{Leb}^{n}$ \\
    条件测度 / disintegration & 耦合 $\pi$ 沿边际 $\mu$ 分解为 $\{\pi_{x}\}_{x\in\X}$；Gluing Lemma 的基础 \\
    \bottomrule
  \end{tabular}
  \end{center}
  \TLtakeaway{§0 背景术语：lsc 保证最优性问题有解；$\Pac$ 保证 Brenier 映射存在；disintegration 支撑 Gluing Lemma。}
\end{frame}

% ── 记号表 · Notation ────────────────────────────────────────────────────────
\begin{frame}{记号表 · Notation}
  \begin{center}\small
  \begin{tabular}{ll}
    \toprule
    \textbf{记号（Symbol）} & \textbf{含义} \\
    \midrule
    $\X,\Y$ & Polish（完备可分度量）空间 \\
    $\Pp(\X)$，$\Ppp{p}(\X)$ & 概率测度全体；$p$ 阶矩有限者 \\
    $\Pac(\X)$ & 绝对连续概率测度（有 Lebesgue 密度）\\
    $\coup{\mu}{\nu}=\Pi(\mu,\nu)$ & $(\mu,\nu)$ 的全体耦合（转移计划） \\
    $\pf{T}{\mu}$ & 推前测度（$T$ 把 $\mu$ 搬成的测度） \\
    $c(x,y)$，$C(\mu,\nu)$ & 单位成本；最优总成本 \\
    $\ctr{\psi}$，$\cctr{\psi}$，$\partial_c\psi$ & c-变换、二次 c-变换、c-次微分 \\
    $\Wp{p}$，$\Wtwo$ & $p$ 阶、$2$ 阶 Wasserstein 距离 \\
    $\mu_k\wkto\mu$ & 弱（窄）收敛 \\
    $e_t$，$\Pi$，$\mu_t$ & 取值映射 $\gamma\mapsto\gamma_t$；路径测度；位移插值 \\
    \bottomrule
  \end{tabular}
  \end{center}
  \TLtakeaway{所有记号在 §0 预备一节给出严格定义；本表供随时回查。}
\end{frame}

% ── 参考文献 · References ─────────────────────────────────────────────────────
\begin{frame}{参考文献 · References（一）：来源与原始文献}
  \textbf{主要来源（source）。}
  \begin{itemize}\small
    \item C. Villani, \emph{Optimal Transport: Old and New}, Grundlehren der math.\ Wissenschaften \textbf{338}, Springer, 2009. ——本讲义覆盖第 I 部分（第 4--8 章）。
  \end{itemize}
  \textbf{原始文献与经典（primary literature）。}
  \begin{itemize}\small
    \item G. Monge, \emph{Mémoire sur la théorie des déblais et des remblais}, 1781.
    \item L.V. Kantorovich, \emph{On the translocation of masses}, Dokl.\ Akad.\ Nauk SSSR, 1942.
    \item R.T. Rockafellar, \emph{Characterization of subdifferentials of convex functions}, 1966.
    \item Y. Brenier, \emph{Polar factorization and monotone rearrangement of vector-valued functions}, CPAM, 1991.
    \item R.J. McCann, \emph{A convexity principle for interacting gases}, Adv.\ Math., 1997.
    \item J.-D. Benamou, Y. Brenier, \emph{A computational fluid mechanics solution to the Monge--Kantorovich problem}, Numer.\ Math., 2000.
  \end{itemize}
\end{frame}

\begin{frame}{参考文献 · References（二）：教材与综述}
  \textbf{延伸阅读（further reading / textbooks）。}
  \begin{itemize}\small
    \item C. Villani, \emph{Topics in Optimal Transportation}, AMS, 2003.
    \item L. Ambrosio, N. Gigli, G. Savaré, \emph{Gradient Flows in Metric Spaces and in the Space of Probability Measures}, Birkhäuser, 2008.
    \item F. Santambrogio, \emph{Optimal Transport for Applied Mathematicians}, Birkhäuser, 2015.
    \item G. Peyré, M. Cuturi, \emph{Computational Optimal Transport}, Found.\ Trends ML, 2019.
    \item A. Figalli, F. Glaudo, \emph{An Invitation to Optimal Transport, Wasserstein Distances, and Gradient Flows}, EMS, 2021.
  \end{itemize}
  \TLtakeaway{从 Villani 的两本书入门定性理论，AGS 深入梯度流，Santambrogio / Peyré--Cuturi 面向应用与计算。}
\end{frame}

% ── 文献注记 · Bibliographical notes ──────────────────────────────────────────
\begin{frame}{文献注记 · Bibliographical Notes（按章）}
  \begin{itemize}\small
    \item \textbf{第 4 章（存在性）。} 最优耦合的存在性是经典结果；关键工具 Prokhorov 紧性定理见 Billingsley。
    \item \textbf{第 5 章（对偶）。} 对偶源于 Kantorovich (1942)；c-循环单调与 c-凸推广了 Rockafellar (1966) 的凸分析；势函数构造借鉴 Rüschendorf；二次成本的 Brenier (1991)、Knott--Smith 同期独立给出。
    \item \textbf{第 6 章（Wasserstein）。} 距离 $\Wp{1}$ 即 Kantorovich--Rubinstein；"Wasserstein" 之名经 Dobrushin 等推广（实由 Vasershtein 引入）。
    \item \textbf{第 7 章（位移插值）。} 位移插值与位移凸性由 McCann (1997) 提出；动力学（连续性方程）表述见 Benamou--Brenier (2000)，Otto 据此发展 $\Ppp{2}$ 上的形式 Riemann 几何。
    \item \textbf{第 8 章（缩短原理）。} 不相交思想可追溯至 Monge；定量"缩短引理"源于 Mather (1991) 的作用量极小测度理论，由 Bernard--Buffoni 引入最优传输。
  \end{itemize}
  \TLtakeaway{这套理论由 Monge（1781）发端，经 Kantorovich（1942）线性化，到 Brenier/McCann（1990s）几何化，再到 Otto/Villani（2000s）联结 Ricci 曲率——本讲义止于其"定性"地基。}
\end{frame}

% =============================================================================
\appendix

\begin{frame}{附录 · Appendix：习题解答 · Exercise Solutions}
  本附录收录 §0--§5 全部习题的完整解答。建议先独立完成、再对照。
  \TLtakeaway{解答是用来\emph{核对}的，不是用来\emph{替代}思考的——先挣扎，后对照。}
\end{frame}

% ── §0 解答 ──────────────────────────────────────────────────────────────────
\begin{frame}{习题解答 · §0（预备）}
  \textbf{0.1}~(\diffI)\ $T(x)=2x$，$\mu=\mathrm{Unif}[0,1]$。对 $y\in(0,2)$，$\nu[(0,y)]=\mu[2x<y]=\mu[x<y/2]=y/2$，故密度 $g(y)=\tfrac12$ 于 $[0,2]$，即 $\nu=\mathrm{Unif}[0,2]$。

  \medskip
  \textbf{0.2}~(\diffII)\ 乘积测度满足 $(\mu\otimes\nu)[A\times\Y]=\mu[A]\nu[\Y]=\mu[A]$，同理第二边际为 $\nu$，故 $\mu\otimes\nu\in\coup{\mu}{\nu}$。于是 $\coup{\mu}{\nu}\neq\varnothing$（Kantorovich 可行集总非空）。

  \medskip
  \textbf{0.3}~(\diffIII)\ $\mu=\delta_0$ 的质量必须分裂到 $\pm1$，无函数 $T$ 能做到，故 Monge 不可行。唯一耦合 $\pi=\tfrac12\delta_{(0,-1)}+\tfrac12\delta_{(0,1)}$，$C(\mu,\nu)=\tfrac12\cdot1+\tfrac12\cdot1=1$。

  \TLtakeaway{§0：推前是换元；乘积测度保证可行集非空；点质量必须分裂时 Monge 失效而 Kantorovich 仍解。}
\end{frame}

% ── §1 解答 ──────────────────────────────────────────────────────────────────
\begin{frame}{习题解答 · §1（基本性质）}
  \textbf{1.1}~(\diffI)\ $\mu\otimes\nu$ 的边际为 $\mu,\nu$（见 0.2）。由 Fubini（$c\ge0$）：$I[\mu\otimes\nu]=\iint c\dd\mu\dd\nu<+\infty$。故 $C(\mu,\nu)\le I[\mu\otimes\nu]<+\infty$。

  \medskip
  \textbf{1.2}~(\diffII)\ 设子计划 $\tilde\pi=\pi'/\pi'[A\times B]$ 不最优，则有同边际而更便宜的 $\hat\pi$。令 $Z=\pi'[A\times B]$，构造 $\pi_{\mathrm{new}}=(\pi^\star-\pi')+Z\hat\pi$：边际与 $\pi^\star$ 相同，成本 $=I[\pi^\star]-Z\,I[\tilde\pi]+Z\,I[\hat\pi]<I[\pi^\star]$，与 $\pi^\star$ 最优矛盾。

  \medskip
  \textbf{1.3}~(\diffIII)\ 例 4.9 中唯一最优 Kantorovich 计划把每个 $(0,a)$ 的质量对半分到 $(\pm1,a)$，不是任何映射的图，故无 Monge 极小子；但用"几乎处处单侧"的确定性映射逼近，其成本趋于 $C(\mu,\nu)$，故 $\inf_{\text{Monge}}=\min_{\text{Kantorovich}}$。

  \TLtakeaway{§1：独立耦合给出有限上界；限制性质靠"局部替换"的反证；Monge 下确界可等于 Kantorovich 最小值却不被取到。}
\end{frame}

% ── §2 解答 ──────────────────────────────────────────────────────────────────
\begin{frame}{习题解答 · §2（对偶）}
  \textbf{2.1}~(\diffI)\ $c(i,j)=|i-j|$，$\mu=\nu=\tfrac13\sum_i\delta_i$。最优耦合为恒等：$\pi^\star=\tfrac13\sum_i\delta_{(i,i)}$，成本 $0$。对任意 $(i,i),(j,j)\in\supp\pi^\star$：$c(i,i)+c(j,j)=0\le c(i,j)+c(j,i)=2|i-j|$，故 c-循环单调。

  \medskip
  \textbf{2.2}~(\diffII)\ $\psi(x)=\tfrac12x^2$，$c=\tfrac12|x-y|^2$。$\ctr{\psi}(y)=\inf_x[\tfrac12x^2+\tfrac12(x-y)^2]$；令导数 $2x-y=0\Rightarrow x=y/2$，得 $\ctr{\psi}(y)=y^2/4$。条件 $\ctr{\psi}(y)-\psi(x)=c(x,y)$ 化为 $(x-y/2)^2=0$，故 $\partial_c\psi=\{(x,y):x=y/2\}$（对应映射 $y=2x$）。

  \medskip
  \textbf{2.3}~(\diffIII)\ 由 Thm 5.10(ii)，两个最优 Monge 耦合 $(\Id,T_k)_\#\mu$ 都集中在同一 $\partial_c\psi$。$\mu$ 绝对连续 + 二次成本 $\Rightarrow$ $\partial_c\psi(x)$ 几乎处处单点（Thm 5.30 条件），故 $T_1=T_2$（$\mu$-a.s.）。

  \TLtakeaway{§2：恒等耦合显然 c-循环单调；c-变换是逐点极小化；二次成本 + 绝对连续给出唯一 Brenier 映射。}
\end{frame}

% ── §3 解答 ──────────────────────────────────────────────────────────────────
\begin{frame}{习题解答 · §3（Wasserstein）}
  \textbf{3.1}~(\diffI)\ Bernoulli$(p_0)$、Bernoulli$(q_0)$（$p_0\le q_0$）。分位函数 $F_\mu^{-1}=\mathbf{1}_{t\ge1-p_0}$，$F_\nu^{-1}=\mathbf{1}_{t\ge1-q_0}$，二者之差仅在 $[1-q_0,1-p_0)$ 上为 $1$。故 $\Wp{1}(\mu,\nu)=q_0-p_0$。

  \medskip
  \textbf{3.2}~(\diffII)\ $\mu_k=(1-\tfrac1k)\delta_0+\tfrac1k\delta_{k^2}$。对有界 $\varphi$，$\int\varphi\dd\mu_k\to\varphi(0)$，故 $\mu_k\wkto\delta_0$；但一阶矩 $\int|x|\dd\mu_k=\tfrac1k\cdot k^2=k\to\infty$，违反 Thm 6.9 的矩收敛条件，故 $\Wp{1}(\mu_k,\delta_0)\to\infty$。

  \medskip
  \textbf{3.3}~(\diffIII)\ 代入多维高斯公式 $\Sigma_1=I_n,m_1=0$：$(I^{1/2}\Sigma I^{1/2})^{1/2}=\Sigma^{1/2}$，故 $\Wtwo^2=|m|^2+\mathrm{tr}(I+\Sigma-2\Sigma^{1/2})=|m|^2+\mathrm{tr}((I-\Sigma^{1/2})^2)$。当 $n=1$，$\Sigma=\sigma^2$：$\Wtwo^2=m^2+(1-\sigma)^2$，与正文高斯例吻合。

  \TLtakeaway{§3：一维 Wasserstein 是分位差的 $L^p$ 范数；质量逃逸破坏矩收敛；多维高斯公式在 $n=1$ 退化为标量公式。}
\end{frame}

% ── §4 解答 ──────────────────────────────────────────────────────────────────
\begin{frame}{习题解答 · §4（位移插值）}
  \textbf{4.1}~(\diffI)\ 位移插值 $\mu_t=\delta_t$；线性插值 $(1-t)\delta_0+t\delta_1$。唯一耦合把 $(1-t)$ 的质量从 $0$ 送到 $t$、$t$ 的质量从 $1$ 送到 $t$：$\Wtwo^2=(1-t)t^2+t(1-t)^2=t(1-t)$，故 $\Wtwo=\sqrt{t(1-t)}$（在 $t=\tfrac12$ 取最大 $\tfrac12$）。

  \medskip
  \textbf{4.2}~(\diffII)\ $\mu_0=\mathcal{N}(0,1),\mu_1=\mathcal{N}(0,4)$：(a) $T(x)=(\sigma_1/\sigma_0)x=2x$；(b) $m_t=0$，$\sigma_t=(1-t)\cdot1+t\cdot2=1+t$，即 $\mu_t=\mathcal{N}(0,(1+t)^2)$；(c) $\Wtwo=|\sigma_1-\sigma_0|=1$。

  \medskip
  \textbf{4.3}~(\diffIII)\ Brenier：$T=\nabla\varphi$，高斯下 $\varphi$ 为二次型故 $T(x)=Ax$ 线性、$A$ 对称正定，$A=\Sigma_0^{-1/2}(\Sigma_0^{1/2}\Sigma_1\Sigma_0^{1/2})^{1/2}\Sigma_0^{-1/2}$。(b) $\Sigma_t=((1-t)I+tA)\Sigma_0((1-t)I+tA)^{\top}$。(c) $\Wtwo^2=\mathrm{tr}(\Sigma_0)+\mathrm{tr}(\Sigma_1)-2\,\mathrm{tr}((\Sigma_0^{1/2}\Sigma_1\Sigma_0^{1/2})^{1/2})$。

  \TLtakeaway{§4：线性与位移插值的差距是 $\sqrt{t(1-t)}$；一维高斯方差线性插值；高维高斯由 Bures 几何给出。}
\end{frame}

% ── §5 解答 ──────────────────────────────────────────────────────────────────
\begin{frame}{习题解答 · §5（缩短原理）}
  \textbf{5.1}~(\diffI)\ $x_1=0,x_2=2$，$c=(x-y)^2$。不相交 $(0\to0,2\to2)$ 成本 $=0$；相交 $(0\to2,2\to0)$ 成本 $=4+4=8$；节省 $8$。一般地，差值 $=|x_1-y_2|^2+|x_2-y_1|^2-|x_1-y_1|^2-|x_2-y_2|^2=2(x_1-x_2)(y_1-y_2)\ge0$（当 $x_1<x_2,\,y_1<y_2$ 时两因子同号）。

  \medskip
  \textbf{5.2}~(\diffII)\ $c$-循环单调给出 $(x_1-y_1)^2+(x_2-y_2)^2\le(x_1-y_2)^2+(x_2-y_1)^2$，展开即 $(x_1-x_2)(y_1-y_2)\ge0$。故 $x_1<x_2\Rightarrow x_1-x_2<0\Rightarrow y_1-y_2\le0$，即 $y_1\le y_2$（支撑单调）。

  \medskip
  \textbf{5.3}~(\diffIII)\ (a) $L=|v|^2$ 为 $C^2$ 且 $\nabla_v^2L=2I>0$；$C(\mu_0,\mu_1)<+\infty$ 已设；$\mu_0\in\Pac$ 满足"端点之一绝对连续"。故定理 8.7 适用，$\mu_t\in\Pac$（$0<t<1$）。
  (b) 对称：定理 8.7 对 $\mu_0$ \emph{或} $\mu_1$ 绝对连续均成立。
  (c) 仍成立——改用 $\mu_1\in\Pac$ 那一分支即可，$\mu_0$ 是否原子无关紧要。

  \TLtakeaway{§5：无交叉的代数核是 $(x_1-x_2)(y_1-y_2)\ge0$；定理 8.7 的四条假设（$L\in C^2$、$\nabla_v^2L>0$、成本有限、端点之一 $\in\Pac$）缺一不可。}
\end{frame}

%% ── 技术备用幻灯片 (Backup) ─────────────────────────────
\section*{技术备用幻灯片 · Backup（推迟的技术估计）}

\begin{frame}{（接 §5 定理 8.1）速度 Hölder 估计：设置与双边界}
  \deflab{设置。}
  固定内部时刻 $t_{0}$ 与小窗口 $\tau>0$，记
  $X_{i}=\gamma_{i}(t_{0})$，$V_{i}=\dot{\gamma}_{i}(t_{0})$。Villani 的目标是证明
  $|V_{1}-V_{2}|\le C|X_{1}-X_{2}|^{\beta}$，即式 (8.8)。

  \deflab{上界 (8.21)。}
  把两条曲线在 $[t_{0}-\tau,t_{0}+\tau]$ 的中段互换；\emph{节省}的作用量
  $G:=A[\gamma_{1}]+A[\gamma_{2}]-A[\widetilde{\gamma}_{1}]-A[\widetilde{\gamma}_{2}]\ge0$
  由 Taylor 误差从上方控制：
  \[
    G \;\le\;
    C\!\left(
      \frac{|X_{1}-X_{2}|^{1+\alpha}}{\tau^{1+\alpha}}
      +\tau |V_{1}-V_{2}|^{1+\alpha}
    \right).
  \]

  \deflab{下界 (8.22)。}
  假设 (iii) 的 $(2+\kappa)$-凸性给出\emph{同一个} $G$ 的下界：
  \[
    G \;\ge\;
    K\bigl(|V_{1}-V_{2}|-A\tau |X_{1}-X_{2}|\bigr)^{2+\kappa}.
  \]

  \TLtakeaway{上界来自 shortcut 的 Taylor 误差；下界来自速度方向的严格凸性。两者夹住同一个节省量 $G\ge0$。}
\end{frame}

\begin{frame}{（接 §5 定理 8.1）速度 Hölder 估计：比较与优化}
  \deflab{情形一。}
  若 $|V_{1}-V_{2}|\le 2A\tau |X_{1}-X_{2}|$，则速度差已被位置差控制；后面选择允许的 $\tau$ 后即可收尾。

  \deflab{情形二。}
  否则有 $|V_{1}-V_{2}|-A\tau |X_{1}-X_{2}|\ge |V_{1}-V_{2}|/2$。比较上一帧的上下界，得到 (8.23)：
  \[
    |V_{1}-V_{2}|^{2+\kappa}
    \le
    C\!\left(
      \frac{|X_{1}-X_{2}|^{1+\alpha}}{\tau^{1+\alpha}}
      +\tau |V_{1}-V_{2}|^{1+\alpha}
    \right).
  \]

  \deflab{优化。}
  当 $|V_{1}-V_{2}|>0$ 时，取 $\tau\sim|V_{1}-V_{2}|^{1+\kappa-\alpha}$
  以吸收右端第二项（Villani 第 178 页），得到
  \[
    |V_{1}-V_{2}|\le C|X_{1}-X_{2}|^{\beta},
    \qquad \beta=\beta(\alpha,\kappa)>0.
  \]
  在二次成本 $c=d^{2}$ 中，$\alpha=1,\kappa=0$，于是 $\beta=1$。

  \TLtakeaway{这是 §5 "定理 8.1 证明思路（凸性矛盾）" 暂缓的定量估计：速度差由中间位置差的幂控制。}
\end{frame}

\begin{frame}{（接 §2 定理 5.10）可测性补证：势函数为何可测}
  \deflab{问题。}
  链式公式 (5.22) 中的 $\psi$ 是不可数个 usc 函数的上确界，不自动是 Borel 函数。

  \deflab{Egorov + 内正则。}
  取 $0\le c_{\ell}\nearrow c$ 为连续逼近。对每个 $k$，可取紧集 $\Gamma_{k}\subset\Gamma$，使得
  \[
    \pi[\Gamma_{k}]\ge 1-\frac{2}{k},\qquad
    c_{\ell}\to c\ \mathrm{uniformly\ on}\ \Gamma_{k},\qquad
    \Gamma=\bigcup_{k}\Gamma_{k}.
  \]

  \deflab{截断势函数。}
  链长截为 $m$，链点限制在 $\Gamma_{k}$，横跳项用 $c_{\ell}$，得 $\psi_{m,k,\ell}(x)$。
  它关于 $x$ 为 lsc；逐点极限是 Borel，且
  \[
    \psi(x)=\sup_{m\in\N}\sup_{k\in\N}\lim_{\ell\to\infty}\psi_{m,k,\ell}(x).
  \]
  右端只含可数次上确界与极限，故 $\psi$ 可测。

  \TLtakeaway{Egorov + 内正则把不可数上确界转化为可数紧集族上的截断链极限，从而得到 $\psi$ 可测。}
\end{frame}

\begin{frame}{（接 §2 定理 5.10）可测性补证：$\phi$ 与可测选择}
  \deflab{$\phi=\ctr{\psi}$ 的可测版本。}
  分解 $\pi$ 为以第二边际为基的条件测度：
  \[
    \pi(\dd x\,\dd y)=\pi(\dd x\mid y)\,\nu(\dd y),
    \qquad
    \widetilde{\phi}(y)=\int_{\X}[\psi(x)+c(x,y)]\,\pi(\dd x\mid y).
  \]
  则 $\phi(y)=\widetilde{\phi}(y)$ 对 $\nu$-a.e.\ 的 $y$ 成立。

  \deflab{$\partial_{c}\psi$ 的可测版本。}
  除去一个 $\pi$-零集后，$\partial_{c}\psi$ 与下列 Borel 集一致：
  \[
    \{(x,y):\widetilde{\phi}(y)-\psi(x)=c(x,y)\}.
  \]

  \deflab{推论 5.22（最优计划的可测选择）。}
  若 $c$ 连续、$\inf c>-\infty$，且 $\omega\mapsto(\mu_{\omega},\nu_{\omega})$ 可测，则可选可测的最优计划 $\omega\mapsto\pi_{\omega}$。
  证明思路：最优计划全体 $O$ 是 Polish；边际映射 $\Phi:O\to\Pp(\X)\times\Pp(\Y)$ 连续且 onto，纤维紧，故可用可测选择定理。

  \TLtakeaway{这补上 §2 "定理 5.10 证明思路（二c）" 暂缓的可测性缺口：势函数、c-次微分与最优计划选择都可取为可测对象。}
\end{frame}
  % §6  定性图景 + 术语表 + 参考文献 + 习题解答附录

\end{document}
